% Options for packages loaded elsewhere
\PassOptionsToPackage{unicode}{hyperref}
\PassOptionsToPackage{hyphens}{url}
\PassOptionsToPackage{dvipsnames,svgnames,x11names}{xcolor}
%
\documentclass[
  11pt,
]{article}
\usepackage{amsmath,amssymb}
\usepackage{iftex}
\ifPDFTeX
  \usepackage[T1]{fontenc}
  \usepackage[utf8]{inputenc}
  \usepackage{textcomp} % provide euro and other symbols
\else % if luatex or xetex
  \usepackage{unicode-math} % this also loads fontspec
  \defaultfontfeatures{Scale=MatchLowercase}
  \defaultfontfeatures[\rmfamily]{Ligatures=TeX,Scale=1}
\fi
\usepackage{lmodern}
\ifPDFTeX\else
  % xetex/luatex font selection
  \setmainfont[]{Liberation Serif}
  \setsansfont[]{Liberation Sans}
  \setmonofont[]{DejaVu Sans Mono}
\fi
% Use upquote if available, for straight quotes in verbatim environments
\IfFileExists{upquote.sty}{\usepackage{upquote}}{}
\IfFileExists{microtype.sty}{% use microtype if available
  \usepackage[]{microtype}
  \UseMicrotypeSet[protrusion]{basicmath} % disable protrusion for tt fonts
}{}
\makeatletter
\@ifundefined{KOMAClassName}{% if non-KOMA class
  \IfFileExists{parskip.sty}{%
    \usepackage{parskip}
  }{% else
    \setlength{\parindent}{0pt}
    \setlength{\parskip}{6pt plus 2pt minus 1pt}}
}{% if KOMA class
  \KOMAoptions{parskip=half}}
\makeatother
\usepackage{xcolor}
\usepackage[margin=0.8in]{geometry}
\usepackage{color}
\usepackage{fancyvrb}
\newcommand{\VerbBar}{|}
\newcommand{\VERB}{\Verb[commandchars=\\\{\}]}
\DefineVerbatimEnvironment{Highlighting}{Verbatim}{commandchars=\\\{\}}
% Add ',fontsize=\small' for more characters per line
\newenvironment{Shaded}{}{}
\newcommand{\AlertTok}[1]{\textcolor[rgb]{1.00,0.00,0.00}{\textbf{#1}}}
\newcommand{\AnnotationTok}[1]{\textcolor[rgb]{0.38,0.63,0.69}{\textbf{\textit{#1}}}}
\newcommand{\AttributeTok}[1]{\textcolor[rgb]{0.49,0.56,0.16}{#1}}
\newcommand{\BaseNTok}[1]{\textcolor[rgb]{0.25,0.63,0.44}{#1}}
\newcommand{\BuiltInTok}[1]{\textcolor[rgb]{0.00,0.50,0.00}{#1}}
\newcommand{\CharTok}[1]{\textcolor[rgb]{0.25,0.44,0.63}{#1}}
\newcommand{\CommentTok}[1]{\textcolor[rgb]{0.38,0.63,0.69}{\textit{#1}}}
\newcommand{\CommentVarTok}[1]{\textcolor[rgb]{0.38,0.63,0.69}{\textbf{\textit{#1}}}}
\newcommand{\ConstantTok}[1]{\textcolor[rgb]{0.53,0.00,0.00}{#1}}
\newcommand{\ControlFlowTok}[1]{\textcolor[rgb]{0.00,0.44,0.13}{\textbf{#1}}}
\newcommand{\DataTypeTok}[1]{\textcolor[rgb]{0.56,0.13,0.00}{#1}}
\newcommand{\DecValTok}[1]{\textcolor[rgb]{0.25,0.63,0.44}{#1}}
\newcommand{\DocumentationTok}[1]{\textcolor[rgb]{0.73,0.13,0.13}{\textit{#1}}}
\newcommand{\ErrorTok}[1]{\textcolor[rgb]{1.00,0.00,0.00}{\textbf{#1}}}
\newcommand{\ExtensionTok}[1]{#1}
\newcommand{\FloatTok}[1]{\textcolor[rgb]{0.25,0.63,0.44}{#1}}
\newcommand{\FunctionTok}[1]{\textcolor[rgb]{0.02,0.16,0.49}{#1}}
\newcommand{\ImportTok}[1]{\textcolor[rgb]{0.00,0.50,0.00}{\textbf{#1}}}
\newcommand{\InformationTok}[1]{\textcolor[rgb]{0.38,0.63,0.69}{\textbf{\textit{#1}}}}
\newcommand{\KeywordTok}[1]{\textcolor[rgb]{0.00,0.44,0.13}{\textbf{#1}}}
\newcommand{\NormalTok}[1]{#1}
\newcommand{\OperatorTok}[1]{\textcolor[rgb]{0.40,0.40,0.40}{#1}}
\newcommand{\OtherTok}[1]{\textcolor[rgb]{0.00,0.44,0.13}{#1}}
\newcommand{\PreprocessorTok}[1]{\textcolor[rgb]{0.74,0.48,0.00}{#1}}
\newcommand{\RegionMarkerTok}[1]{#1}
\newcommand{\SpecialCharTok}[1]{\textcolor[rgb]{0.25,0.44,0.63}{#1}}
\newcommand{\SpecialStringTok}[1]{\textcolor[rgb]{0.73,0.40,0.53}{#1}}
\newcommand{\StringTok}[1]{\textcolor[rgb]{0.25,0.44,0.63}{#1}}
\newcommand{\VariableTok}[1]{\textcolor[rgb]{0.10,0.09,0.49}{#1}}
\newcommand{\VerbatimStringTok}[1]{\textcolor[rgb]{0.25,0.44,0.63}{#1}}
\newcommand{\WarningTok}[1]{\textcolor[rgb]{0.38,0.63,0.69}{\textbf{\textit{#1}}}}
\usepackage{longtable,booktabs,array}
\usepackage{calc} % for calculating minipage widths
% Correct order of tables after \paragraph or \subparagraph
\usepackage{etoolbox}
\makeatletter
\patchcmd\longtable{\par}{\if@noskipsec\mbox{}\fi\par}{}{}
\makeatother
% Allow footnotes in longtable head/foot
\IfFileExists{footnotehyper.sty}{\usepackage{footnotehyper}}{\usepackage{footnote}}
\makesavenoteenv{longtable}
\setlength{\emergencystretch}{3em} % prevent overfull lines
\providecommand{\tightlist}{%
  \setlength{\itemsep}{0pt}\setlength{\parskip}{0pt}}
\setcounter{secnumdepth}{-\maxdimen} % remove section numbering
\ifLuaTeX
\usepackage[bidi=basic]{babel}
\else
\usepackage[bidi=default]{babel}
\fi
\babelprovide[main,import]{american}
\ifPDFTeX
\else
\babelfont{rm}[]{Liberation Serif}
\fi
% get rid of language-specific shorthands (see #6817):
\let\LanguageShortHands\languageshorthands
\def\languageshorthands#1{}
\usepackage{adjustbox}
\emergencystretch=3em
\AtBeginEnvironment{longtable}{\small\sloppy}
\usepackage{stmaryrd}
% Printing shims for run-on conditional-bar tokens in the supplied Biology source.
\providecommand{\midL}{\mid L}
\providecommand{\midN}{\mid N}
\sloppy
\ifLuaTeX
  \usepackage{selnolig}  % disable illegal ligatures
\fi
\usepackage{bookmark}
\IfFileExists{xurl.sty}{\usepackage{xurl}}{} % add URL line breaks if available
\urlstyle{same}
\hypersetup{
  pdflang={en-US},
  colorlinks=true,
  linkcolor={blue},
  filecolor={Maroon},
  citecolor={Blue},
  urlcolor={blue},
  pdfcreator={LaTeX via pandoc}}

\author{}
\date{}

\begin{document}

\section{COHERENCE UNIVERSALISM:
FOUNDATIONS}\label{coherence-universalism-foundations}

\subsection{Principles, Domain Gateways, Commitments, and Constraints on
Responsible
Application}\label{principles-domain-gateways-commitments-and-constraints-on-responsible-application}

\textbf{Gaura Kiśora Dās Rader}\\
\textbf{Heaven≡Earth Foundation}\\
\textbf{White Paper Series}

\textbf{Version 0.10}\\
\textbf{October 2026}

\textbf{Release status.} This manuscript is part of the complete Version
0.10 proposal. Numbered principles identify the framework's current
formulations; their claim status, external validation, and authorization
for consequential application remain separate questions.

\textbf{Freeze status.} This manuscript is synchronized to the frozen
Version 0.10 release registry, evidence package, and review gates.
Future changes require a new versioned release.

\begin{center}\rule{0.5\linewidth}{0.5pt}\end{center}

\section{Abstract}\label{abstract}

\emph{Foundations} is the constitutional reference for Coherence
Universalism. It defines the framework's shared vocabulary, records the
principles that govern inquiry and consequential action, preserves the
ownership of claims by their proper domains, and maintains the
interpretive and commitment registers through which the framework
becomes a spiritual and practical orientation without presenting those
orientations as scientific deductions.

Coherence is defined as integrated relational organization under
specified constraints. Every substantive application must identify its
system or trajectory, constraints, relational structure, scale or access
map, and temporal horizon. Coherence may be represented by a scalar,
profile, partial ordering, viability region, threshold family, or
trajectory property; no universal scalar is assumed.

The document contains the active foundational, methodological, logical,
epistemic, dynamical, governance, physical, biological, consciousness,
psychological, social, ethical, artificial-intelligence, interpretive,
and commitment principles of Version 0.10. It includes thirteen
interpretive principles and nineteen commitment principles. It also
restores the original sixty-four-stage Coherence Mandala as a stable
philosophical and contemplative architecture while keeping each stage
answerable to its owning domain and claim status.

The complete proposal includes structured evolutionary possibility, Open
Myth and narrative pluralism, and mythic--constitutional alignment.
Their canonical inclusion identifies what CU proposes; it does not
convert conditional empirical synthesis, civilizational design, or
experimental AI governance into established fact or deployment
authorization.

\textbf{Keywords:} coherence; constraint; principle registry;
interdisciplinarity; epistemology; governance; identity; love; meaning;
Omega; Coherence Ladder; Coherence Mandala

\subsection{Relation to the series}\label{relation-to-the-series}

This document is the shared constitutional reference for the Coherence
Universalism v0.10 series. Each domain paper can be read independently
and owns its scientific, clinical, social, ethical, or engineering
claims. \emph{Metaphysics and Epistemology} develops the argument and
formal architecture. \emph{Foundations} supplies the statements,
statuses, cross-references, boundaries, and revision conditions that the
rest of the corpus cites.

\subsection{Note on Version 0.10}\label{note-on-version-0.10}

Version 0.10 presents the complete current proposal rather than
withholding central ideas until after expert review. Authority and
validation remain distinct: a canonical principle is the formulation CU
presently advances, while its claim status, review requirements, misuse
boundaries, and conditions of consequential application remain explicit.

The Yi Jing traversal of the sixty-four stages belongs to \emph{You Are
the Universe Becoming Conscious of Itself} as a visual and contemplative
interpretation. The present document retains only a cross-reference. It
does not make the hexagrams part of stage identity or evidence.

\begin{center}\rule{0.5\linewidth}{0.5pt}\end{center}

\section{Preface: A Constitution for an Open
Framework}\label{preface-a-constitution-for-an-open-framework}

A foundational document should make a framework easier to question, not
harder.

Coherence Universalism seeks relations among physical structure, living
organization, perspective, meaning, social coordination, ethics, and
artificial intelligence. Those relations become credible only when the
framework distinguishes definitions from discoveries, mathematical
results from physical hypotheses, empirical generalizations from bridge
claims, normative commitments from descriptive facts, and spiritual
interpretation from scientific proof.

This document is therefore a constitution rather than a catalogue of
certainties. It states the vocabulary shared across the series, the
limits on cross-domain translation, the principles governing
consequential application, and the rules by which the framework may
lose, revise, or retire a claim. Domain-specific questions remain with
their owning papers. A principle's authority identifies the formulation
to cite; it does not imply that the claim has been proven beyond its
stated status.

The governing maxim is:

\begin{quote}
\textbf{Preserve the integrative vision by claiming no more integration
than the argument, evidence, and affected communities can support.}
\end{quote}

\section{How to Read This Document}\label{how-to-read-this-document}

Each principle contains, where relevant:

\begin{enumerate}
\def\labelenumi{\arabic{enumi}.}
\tightlist
\item
  \textbf{Authority} --- whether the principle governs the framework, a
  domain, an application, or an interpretation;
\item
  \textbf{Claim status} --- definition, formal result, conditional
  result, empirical claim, bridge, hypothesis, interpretation,
  commitment, or open question;
\item
  \textbf{Statement} --- the controlling wording;
\item
  \textbf{Scope} --- what the principle does and does not say;
\item
  \textbf{Operational consequence} --- what an application must do;
\item
  \textbf{Failure modes} --- recurrent misreadings or misuse;
\item
  \textbf{Review gate} --- the expertise or participation required
  before consequential use or later revision.
\end{enumerate}

A reference to coherence without a domain, constraints, relational
structure, scale or access map, and horizon is presumptively incomplete.
An identifier supports navigation; it never replaces the
ordinary-language statement of the principle.

\section{Status Vocabulary}\label{status-vocabulary}

Authority class and claim status are separate.

\begin{itemize}
\tightlist
\item
  \textbf{Framework principle:} controls shared definitions, methods, or
  boundary rules across the series.
\item
  \textbf{Framework conditional principle:} controls only under declared
  assumptions or model classes.
\item
  \textbf{Framework governance rule:} governs consequential use but is
  not a descriptive law of nature.
\item
  \textbf{Domain authority:} the controlling CU formulation within one
  domain.
\item
  \textbf{Domain authority---conditional:} the controlling formulation
  with explicitly conditional empirical or bridge content.
\item
  \textbf{Provisional domain principle:} a retained hypothesis or
  construct not yet promoted to broader authority.
\item
  \textbf{Domain governance rule:} an application requirement owned by a
  domain.
\item
  \textbf{Diagnostic failure mode:} a named pattern of failure rather
  than a positive law.
\item
  \textbf{Interpretive principle:} a metaphysical or spiritual reading
  compatible with, but not entailed by, the descriptive framework.
\item
  \textbf{Commitment principle:} an official CU orientation constrained
  by facts, consequences, standing, and revision.
\item
  \textbf{Historical alias or retired claim:} preserved for traceability
  but not active as independent authority.
\end{itemize}

Canonical inclusion means ``this is the formulation CU presently
proposes.'' It does not mean ``external review is complete,'' ``the
empirical claim is proven,'' or ``consequential application is
authorized.''

\section{Contents}\label{contents}

\begin{itemize}
\tightlist
\item
  Part I --- Foundational Principles
\item
  Part II --- Method and Claim Governance
\item
  Part III --- Logic, Registers, and Epistemology
\item
  Part IV --- Dynamics of Coherence
\item
  Part V --- Cross-Domain Governance
\item
  Part VI --- Physics Constraint Registry
\item
  Part VII --- Domain Gateways and Owned Registries
\item
  Part VIII --- Interpretive and Spiritual Register
\item
  Part IX --- The Commitment Register
\item
  Part X --- Ladder, Mandala, and Corpus Governance
\item
  Appendix A --- Registry Status Summary
\item
  Appendix B --- Merged and Superseded Identifiers
\item
  Appendix C --- Participation and Review Matrix
\item
  Appendix D --- Retired or Superseded Framework Claims
\item
  Appendix E --- Claim and Principle Template
\item
  Appendix F --- Synchronization Checklist
\item
  References and Internal Authorities
\end{itemize}

\begin{center}\rule{0.5\linewidth}{0.5pt}\end{center}

\section{Domain Analytic Frame and CAP
Record}\label{domain-analytic-frame-and-cap-record}

\textbf{Target:} Coherence Universalism Version 0.10 release\\
\textbf{Status:} Release declaration; scope-limited and revisable under
the declared review gates\\
\textbf{Use rule:} Unresolved fields are left explicit rather than
completed by analogy.

\subsection{Paper and model identity}\label{paper-and-model-identity}

\begin{longtable}[]{@{}
  >{\raggedright\arraybackslash}p{(\columnwidth - 4\tabcolsep) * \real{0.3333}}
  >{\raggedright\arraybackslash}p{(\columnwidth - 4\tabcolsep) * \real{0.3333}}
  >{\raggedright\arraybackslash}p{(\columnwidth - 4\tabcolsep) * \real{0.3333}}@{}}
\toprule\noalign{}
\begin{minipage}[b]{\linewidth}\raggedright
Field
\end{minipage} & \begin{minipage}[b]{\linewidth}\raggedright
Declaration
\end{minipage} & \begin{minipage}[b]{\linewidth}\raggedright
Status
\end{minipage} \\
\midrule\noalign{}
\endhead
\bottomrule\noalign{}
\endlastfoot
Paper & Coherence Universalism: Foundations & Specified \\
Domain owner & Framework constitution and principle registry & Specified
by the Version 0.10 release corpus \\
Primary model or analysis & A versioned constitutional reference,
dependency graph, and governance architecture & Paper-relative \\
Primary claim status & Framework definitions, methods, governance rules,
domain gateways, interpretations, and commitments & Mixed statuses
remain explicit \\
Review gate & Corpus-wide conceptual, formal, domain, ethical, and
affected-community review as specified per principle & Proportionate to
claim and consequential reach \\
Methodological authority & CU-FP2, CU-M2, CU-M11, CU-M12, CU-M13,
CU-M14, and CU-M15 & Controlling Version 0.10 methodological
framework \\
\end{longtable}

\subsection{Shared Analytic Frame}\label{shared-analytic-frame}

\begin{longtable}[]{@{}
  >{\raggedright\arraybackslash}p{(\columnwidth - 4\tabcolsep) * \real{0.3333}}
  >{\raggedright\arraybackslash}p{(\columnwidth - 4\tabcolsep) * \real{0.3333}}
  >{\raggedright\arraybackslash}p{(\columnwidth - 4\tabcolsep) * \real{0.3333}}@{}}
\toprule\noalign{}
\begin{minipage}[b]{\linewidth}\raggedright
\#
\end{minipage} & \begin{minipage}[b]{\linewidth}\raggedright
Slot and domain declaration
\end{minipage} & \begin{minipage}[b]{\linewidth}\raggedright
Status and evidence
\end{minipage} \\
\midrule\noalign{}
\endhead
\bottomrule\noalign{}
\endlastfoot
1 & \textbf{Unit and boundary.} The unit is the Version 0.10 principle
architecture and its relation to the reader-facing corpus, registries,
audits, and review gates. & Specified \\
2 & \textbf{State or trajectory space.} The relevant state is a
versioned registry with identifiers, statuses, dependencies, wording,
and supersession history. & Versioned rather than physical \\
3 & \textbf{Environment.} Owning domain papers, technical supplements,
external literature, reviewers, and affected communities constrain the
constitution. & Specified by domain ownership \\
4 & \textbf{Constraints and viable region.} Claim-status discipline,
domain ownership, explicit bridges, non-regression rules, review gates,
and revision loss constrain admissible versions. &
Framework-canonical \\
5 & \textbf{Actions, options, and transitions.} Principles may be added,
revised, narrowed, retired, superseded, cross-referenced, or propagated
through dependencies. & Specified by revision rules \\
6 & \textbf{Dynamics or update rule.} No empirical dynamics is claimed;
change proceeds through versioned proposal, review, decision,
synchronization, and audit. & Governance process \\
7 & \textbf{Observation and access.} Source ledgers, proofs, empirical
evidence, review reports, code, registries, and affected-party testimony
provide access. & Declared in CU-M5--CU-M9 \\
8 & \textbf{Perturbation and intervention class.} Audits,
counterexamples, new evidence, domain corrections, and authorial changes
can alter the registry. & Specified \\
9 & \textbf{Equilibria, basins, or regimes.} Not applicable. A stable
release is an editorial state, not a natural attractor. & Not
applicable \\
10 & \textbf{Local--global status.} CU-M13 governs relation status
across all domains; Foundations asserts no one corpus-wide
decomposition. & Framework rule \\
11 & \textbf{Scalar status.} No scalar evaluates the framework as a
whole. & None \\
12 & \textbf{Statistical-image and allocation status.} The
constitutional registry is not a statistical distribution. CU-M14
governs licensed statistical and allocation subobjects referenced by the
registry. & Methodological rule, not object-level license \\
13 & \textbf{Translation status.} CU-M2 records the license required for
every cross-domain transfer. & Framework-canonical proposal \\
14 & \textbf{Open gaps.} Formal, measurement, governance, and domain
review of the new method cluster remains pending. & Explicit adoption
gates \\
\end{longtable}

\subsection{Compact Coherence Analysis Protocol
record}\label{compact-coherence-analysis-protocol-record}

\begin{longtable}[]{@{}
  >{\raggedright\arraybackslash}p{(\columnwidth - 4\tabcolsep) * \real{0.3333}}
  >{\raggedright\arraybackslash}p{(\columnwidth - 4\tabcolsep) * \real{0.3333}}
  >{\raggedright\arraybackslash}p{(\columnwidth - 4\tabcolsep) * \real{0.3333}}@{}}
\toprule\noalign{}
\begin{minipage}[b]{\linewidth}\raggedright
CAP module
\end{minipage} & \begin{minipage}[b]{\linewidth}\raggedright
Declaration
\end{minipage} & \begin{minipage}[b]{\linewidth}\raggedright
Boundary
\end{minipage} \\
\midrule\noalign{}
\endhead
\bottomrule\noalign{}
\endlastfoot
System declaration & Declare the controlling corpus, version, domain
owners, identifiers, and dependency boundary. & The Version 0.10 release
registry and evidence archive control this edition \\
Observation declaration & Preserve source provenance, audits, review
competence, uncertainty, and unresolved dissent. & Registry authority is
not validation \\
Evaluation declaration & Evaluate proposed wording against
non-regression, domain ownership, claim status, and downstream
consistency. & No global score \\
Power and authorship declaration & Authorial approval, expert review,
and affected-community standing are separate authorities. & No single
review competence overrides another \\
Revision declaration & Every canonical change is dated, explained,
propagated, and preserved with its superseded form. & Required by
CU-M9 \\
\end{longtable}

\subsection{License summary}\label{license-summary}

\begin{longtable}[]{@{}
  >{\raggedright\arraybackslash}p{(\columnwidth - 4\tabcolsep) * \real{0.3333}}
  >{\raggedright\arraybackslash}p{(\columnwidth - 4\tabcolsep) * \real{0.3333}}
  >{\raggedright\arraybackslash}p{(\columnwidth - 4\tabcolsep) * \real{0.3333}}@{}}
\toprule\noalign{}
\begin{minipage}[b]{\linewidth}\raggedright
License
\end{minipage} & \begin{minipage}[b]{\linewidth}\raggedright
Declaration
\end{minipage} & \begin{minipage}[b]{\linewidth}\raggedright
Current status
\end{minipage} \\
\midrule\noalign{}
\endhead
\bottomrule\noalign{}
\endlastfoot
Scalar license & Not applicable to the constitution as a whole. &
None \\
Local--global license & Defines the declaration rule; does not
instantiate one quantitative relation. & Framework-level \\
Statistical Image License & The registry itself is not a statistical
image; individual domain models may receive the license. & Not
applicable to the constitution \\
Allocation License & The registry is not an allocation; the principle
records the test for domain allocations. & Method-level only \\
Refinement and Coarse-Graining License & The constitution specifies the
declaration requirements but does not itself instantiate probabilistic
coarse-graining. & Method-level only \\
Fisher-Rao License & The constitution records the imported classical
theorem and its limits; it does not itself receive a metric. &
Method-level only \\
Translation license & Defines and records Levels 0--4. & Proposed
canonical method \\
\end{longtable}

\subsection{Open gaps and review
disposition}\label{open-gaps-and-review-disposition}

\begin{longtable}[]{@{}
  >{\raggedright\arraybackslash}p{(\columnwidth - 4\tabcolsep) * \real{0.3333}}
  >{\raggedright\arraybackslash}p{(\columnwidth - 4\tabcolsep) * \real{0.3333}}
  >{\raggedright\arraybackslash}p{(\columnwidth - 4\tabcolsep) * \real{0.3333}}@{}}
\toprule\noalign{}
\begin{minipage}[b]{\linewidth}\raggedright
Gap
\end{minipage} & \begin{minipage}[b]{\linewidth}\raggedright
Consequence
\end{minipage} & \begin{minipage}[b]{\linewidth}\raggedright
Required work
\end{minipage} \\
\midrule\noalign{}
\endhead
\bottomrule\noalign{}
\endlastfoot
Review completion & Prevents the patch from becoming controlling
prematurely. & Complete formal, measurement, governance, and domain
review \\
Registry synchronization & Identifiers can drift if only prose is
changed. & Update canonical, active, crosswalk, glossary, and dependency
registries \\
Propagation audit & Domain papers can silently retain earlier method
language. & Run corpus-wide cross-reference and terminology QA \\
\end{longtable}

\begin{quote}
\textbf{Propagation rule.} This card makes the paper's existing objects
and limits visible. It does not strengthen any empirical, mathematical,
normative, metaphysical, or spiritual claim beyond the wording supported
in the manuscript.
\end{quote}

\begin{center}\rule{0.5\linewidth}{0.5pt}\end{center}

\section{Part I --- Foundational
Principles}\label{part-i-foundational-principles}

\subsection{CU-FP1 --- Coherence Under
Constraint}\label{cu-fp1-coherence-under-constraint}

\textbf{Registry status:} CANONICAL\\
\textbf{Claim status:} {[}Definition{]} and bounded {[}Bridge
principle{]}

\subsubsection{Statement}\label{statement}

Coherence is \textbf{integrated relational organization under specified
constraints}.

For any substantive domain \(\mathcal D\), a coherence assessment has
the dependency form

\[\adjustbox{max width=\linewidth}{$\displaystyle 
C_{\mathcal D}=C_{\mathcal D}(x;K,G,\pi,T),
$}\]

where \(x\) is the state, trajectory, or organization; \(K\) the
relevant constraints; \(G\) the relational or coupling structure;
\(\pi\) the scale, partition, observation, or access map; and \(T\) the
temporal horizon.

\subsubsection{Scope}\label{scope}

CU-FP1 defines a shared structural question. It does not assert that
coherence is a substance, field, force, cosmic purpose, or single
measurable quantity. It does not imply that everything is harmonious or
morally good.

The bounded transcendental claim is that robust uses of identity,
organization, persistence, meaning, or justification presuppose some
relations and constraints by which the relevant object can be
distinguished and tracked. This does not derive a unique ontology.

\subsubsection{Operational consequence}\label{operational-consequence}

Every domain paper must specify the arguments hidden by any shorthand
\(C(x)\). Where no defensible measure exists, the paper must use a
qualitative profile, partial order, viability region, or explicit open
problem rather than inventing a scalar.

\subsubsection{Failure modes}\label{failure-modes}

\begin{itemize}
\tightlist
\item
  treating coherence as an invisible substance;
\item
  defining it as whatever succeeds;
\item
  equating it with harmony, order, consistency, or agreement;
\item
  using the word without specifying the system boundary;
\item
  treating the definition as evidence that the universe is fundamentally
  a coherence field.
\end{itemize}

\begin{center}\rule{0.5\linewidth}{0.5pt}\end{center}

\subsection{CU-FP2 --- Gradation, Profiles, Typed Scalars, Direction,
and
Incomparability}\label{cu-fp2-gradation-profiles-typed-scalars-direction-and-incomparability}

\textbf{Authority class:} FRAMEWORK-CANONICAL\\
\textbf{Claim status:} {[}Definition{]} and methodological constraint

\subsubsection{Statement}\label{statement-1}

Coherence assessments may admit degrees, directions of change,
thresholds, vector profiles, preorders, partial orderings, viable
regions, Pareto frontiers, risk functionals, or trajectory properties,
but none of these structures is assumed universally.

A domain may use:

\begin{itemize}
\tightlist
\item
  a measurement scalar for one declared construct;
\item
  a vector profile \(\mathbf C\);
\item
  a decision scalarization for one declared choice;
\item
  a risk functional over uncertain states or trajectories;
\item
  a preorder or partial order \(\preceq_C\);
\item
  a viable set \(V\);
\item
  a threshold family;
\item
  a Pareto frontier;
\item
  or a trajectory property.
\end{itemize}

Scalars are local, typed, purpose-bound, and scope-limited. A
measurement scalar requires an independently defended interpretation of
what it measures. A decision scalarization is a rule for choosing under
declared tradeoffs, not evidence that the dimensions possess one natural
total ordering. A risk functional additionally encodes a treatment of
uncertainty, tails, delay, and irreversibility.

Protected conditions are screened before scoring. No scalar may
compensate for a floor violation. Variation above a protected floor may
enter post-screen comparison only when its measurement, comparability,
authority, and decision relevance are independently justified; otherwise
the coordinate remains lexicographic or outside scoring.

\subsubsection{Scope}\label{scope-1}

Not every system lies on one line from less coherent to more coherent.
Two organizations can remain incomparable because they trade robustness,
flexibility, accuracy, integration, agency, welfare, distribution,
reversibility, or long-horizon resilience differently.

Coordinate-level measurement does not require collapse of the whole
profile. A validated measure of repair latency, regulatory capacity,
information integrity, reversibility, or another coordinate may be
useful while the larger profile remains only partially ordered.

A scalar validated for one construct, population, option set,
institution, horizon, or domain does not transfer beyond that scope
without a new argument. A family of local scalars may not be assembled
by semantic drift into one universal coherence scalar.

\subsubsection{Operational consequence}\label{operational-consequence-1}

Every scalar must declare its type.

A \textbf{measurement scalar} must state:

\begin{itemize}
\tightlist
\item
  the construct measured;
\item
  units or scale interpretation;
\item
  observation and error model;
\item
  reliability and validity evidence;
\item
  population and domain of validation;
\item
  uncertainty;
\item
  and any decomposition or invariance conditions.
\end{itemize}

A \textbf{decision scalarization} must state:

\begin{itemize}
\tightlist
\item
  why aggregation is required;
\item
  normalization;
\item
  weights or trade rates;
\item
  who has authority to set them;
\item
  protected floors and the screening rule;
\item
  affected scope and distribution;
\item
  temporal horizon;
\item
  tie convention;
\item
  decision purpose;
\item
  and sensitivity to reasonable alternative specifications.
\end{itemize}

A \textbf{risk functional} must additionally state:

\begin{itemize}
\tightlist
\item
  the probability or uncertainty model;
\item
  the loss interpretation;
\item
  treatment of tail events;
\item
  discounting or horizon treatment;
\item
  treatment of severe or irreversible harm;
\item
  and robustness to model uncertainty.
\end{itemize}

Where profiles are non-dominating, validation is insufficient, or the
tradeoffs are not defensibly commensurable for the declared purpose,
incomparability or underdetermination must be preserved.

\subsubsection{Failure modes}\label{failure-modes-1}

\begin{itemize}
\tightlist
\item
  binary coherent/incoherent labeling where the relevant structure is
  graded or plural;
\item
  assuming a metric, geometry, gradient, or optimum before defining and
  defending it;
\item
  presenting a decision rule as a discovered natural ordering;
\item
  hiding protected conditions inside weights;
\item
  allowing aggregate gains to compensate a floor violation;
\item
  treating every above-floor difference as automatically tradable;
\item
  treating a local score increase as global progress;
\item
  transferring a scalar beyond its validation or authorization domain;
\item
  and rebuilding a universal coherence scalar from purpose-bound
  instruments.
\end{itemize}

\subsubsection{Review gate}\label{review-gate}

Corpus-wide conceptual review; measurement, decision-theoretic,
statistical, ethical, and affected-community review proportionate to the
scalar's purpose and reach.

\subsubsection{Related principles}\label{related-principles}

CU-M2; CU-M4; CU-M6; CU-M10; CU-M11; CU-M13; CU-G1; CU-G2; CU-G5; CU-G7.

\begin{center}\rule{0.5\linewidth}{0.5pt}\end{center}

\subsection{CU-FP3 --- Generative
Constraint}\label{cu-fp3-generative-constraint}

\textbf{Registry status:} CANONICAL\\
\textbf{Claim status:} {[}Definition{]} and {[}Bridge principle{]}

\subsubsection{Statement}\label{statement-2}

Constraints do not merely remove possibilities. Under suitable
organization, restricting a space of states or transitions can create
stable relations, capacities, meanings, and forms of agency that would
not otherwise exist.

\subsubsection{Scope}\label{scope-2}

CU distinguishes restrictive, generative, constitutive, instrumental,
viability-bearing, protected, and epistemic constraints. Their roles
must not be inferred from mathematical form alone.

Constraint without integration can produce rigidity or noise.
Integration without external constraint can produce fantasy,
confabulation, or insulation.

\subsubsection{Operational consequence}\label{operational-consequence-2}

For each important constraint, identify:

\begin{itemize}
\tightlist
\item
  who or what imposes it;
\item
  what possibilities it removes;
\item
  what capacities it enables;
\item
  whether it is revisable;
\item
  whose viability it protects;
\item
  and whose agency it may restrict.
\end{itemize}

\subsubsection{Failure modes}\label{failure-modes-2}

\begin{itemize}
\tightlist
\item
  treating all constraint as oppression;
\item
  treating all order-producing constraints as legitimate;
\item
  mistaking an externally supplied AI loss function for an intrinsic
  stake;
\item
  ignoring coercion because the resulting system appears stable.
\end{itemize}

\begin{center}\rule{0.5\linewidth}{0.5pt}\end{center}

\subsection{CU-FP4 --- Multi-Scale and Multi-Horizon
Evaluation}\label{cu-fp4-multi-scale-and-multi-horizon-evaluation}

\textbf{Registry status:} CANONICAL\\
\textbf{Claim status:} {[}Methodological principle{]}

\subsubsection{Statement}\label{statement-3}

A coherence claim is relative to scale and horizon. Apparent improvement
within one boundary is not a global improvement when it depends on
unmodeled damage exported to another system, scale, population, or
future.

\subsubsection{Scope}\label{scope-3}

``Local'' and ``global'' are analysis-relative. The containing system
for one question may be a subsystem for another. No finite analysis
captures every consequence, but the boundary and omissions must be
disclosed.

\subsubsection{Operational consequence}\label{operational-consequence-3}

Claims of progress must identify:

\begin{itemize}
\tightlist
\item
  the local system;
\item
  coupled systems;
\item
  immediate and delayed effects;
\item
  externalized costs;
\item
  protected agents or capacities;
\item
  and the horizon over which evaluation is made.
\end{itemize}

\subsubsection{Failure modes}\label{failure-modes-3}

\begin{itemize}
\tightlist
\item
  aggregate success purchased by local destruction;
\item
  short-term relief that entrenches long-term fragility;
\item
  institutional stability maintained through moral injury;
\item
  AI optimization that degrades the human environment supplying its
  objectives;
\item
  rhetoric that labels the preferred scale ``global'' by definition.
\end{itemize}

\begin{center}\rule{0.5\linewidth}{0.5pt}\end{center}

\subsection{CU-FP5 --- Non-Reduction and Explicit
Translation}\label{cu-fp5-non-reduction-and-explicit-translation}

\textbf{Registry status:} CANONICAL\\
\textbf{Claim status:} {[}Methodological principle{]}

\subsubsection{Statement}\label{statement-4}

Cross-domain integration requires explicit translation rather than
identity by vocabulary. A valid translation states the source construct,
target construct, preserved relation, new target constraints, lost
information, evidence, alternatives, and failure conditions.

\subsubsection{Scope}\label{scope-4}

Higher-level phenomena depend on physical realization but may require
concepts and constraints not available in a lower-level vocabulary.
Non-reduction does not imply exemption from physical or empirical
constraint.

\subsubsection{Operational consequence}\label{operational-consequence-4}

Every cross-domain use of terms such as coherence, entropy, memory,
agency, information, field, attractor, gradient, or identity must
provide a translation card.

\subsubsection{Failure modes}\label{failure-modes-4}

\begin{itemize}
\tightlist
\item
  saying two domains are ``the same pattern'' without a map;
\item
  reducing moral normativity to physical stability;
\item
  treating metaphor as mechanism;
\item
  anti-reductionism that protects claims from evidence;
\item
  assuming successful mathematics carries its interpretation unchanged
  between domains.
\end{itemize}

\begin{center}\rule{0.5\linewidth}{0.5pt}\end{center}

\subsection{CU-FP6 --- The Accessibility--Entropy
Relation}\label{cu-fp6-the-accessibilityentropy-relation}

\textbf{Registry status:} CANONICAL--CONDITIONAL\\
\textbf{Claim status:} {[}Formal result{]} in the finite classical
setting; {[}Bridge principle{]} elsewhere

\subsubsection{Statement}\label{statement-5}

For a specified finite partition, block averaging preserves coarse block
masses while erasing within-block distinctions. Under that operation,
statistical differentiation from global uniformity weakly decreases and
Shannon entropy weakly increases. The losses and gains have the same
equality condition but are not generally equal quantities.

Let \(u_N=(1/N,\ldots,1/N)\), define \(D_N(p)=\lVert p-u_N
Vert_2^2\), let \(B_b\) be blocks of size \(m_b\), let
\(q_b=\sum_{i\in B_b}p_i\), and, when \(q_b>0\), let
\(r_i^{(b)}=p_i/q_b\) for \(i\in B_b\). Blocks with \(q_b=0\) contribute
zero to the weighted KL sum. Define

\[\adjustbox{max width=\linewidth}{$\displaystyle 
[E_\pi(p)]_i=\frac{q_b}{m_b}
\quad(i\in B_b).
$}\]

Then

\[\adjustbox{max width=\linewidth}{$\displaystyle 
D_N(p)-D_N(E_\pi p)
=
\sum_b\sum_{i\in B_b}
\left(p_i-\frac{q_b}{m_b}\right)^2,
$}\]

and

\[\adjustbox{max width=\linewidth}{$\displaystyle 
S(E_\pi p)-S(p)
=
\sum_bq_b\,
D_{\mathrm{KL}}(r^{(b)}\Vert u_{m_b})
\geq0.
$}\]

\subsubsection{Scope}\label{scope-5}

This is an operation-specific information-theoretic theorem. It does not
identify entropy with hidden coherence. It does not show that
descriptive coarse-graining physically destroys information. It does not
equate Shannon, thermodynamic, von Neumann, conditional, or horizon
entropy.

\subsubsection{Operational consequence}\label{operational-consequence-5}

Every entropy claim must state:

\begin{itemize}
\tightlist
\item
  the entropy notion;
\item
  state space;
\item
  partition or reduction;
\item
  observer or subsystem;
\item
  and whether the operation is descriptive reconstruction or physical
  dynamics.
\end{itemize}

\subsubsection{Failure modes}\label{failure-modes-5}

\begin{itemize}
\tightlist
\item
  ``entropy is inaccessible coherence'' as an identity;
\item
  calling conditional entropy hidden differentiation;
\item
  inferring physical production from observer ignorance;
\item
  using the theorem as a universal First Law of psychology, society, or
  ethics.
\end{itemize}

\begin{center}\rule{0.5\linewidth}{0.5pt}\end{center}

\subsection{CU-FP7 --- Failure, Collapse, Adaptation, and
Repair}\label{cu-fp7-failure-collapse-adaptation-and-repair}

\textbf{Authority class:} FRAMEWORK-CANONICAL as a schema\\
\textbf{Claim status:} {[}Definition{]} and {[}Bridge principle{]}

\subsubsection{Statement}\label{statement-6}

A domain-specific failure occurs when organization required by an
explicitly declared criterion can no longer be maintained, adapted, or
repaired within the relevant horizon and intervention class.

CU distinguishes perturbation, degradation, drift, fragmentation,
threshold crossing, collapse, compensation, accommodation,
transformation, restitution, restoration of standing, non-recurrence,
and repair.

\subsubsection{Scope}\label{scope-6}

Collapse is not one universal mathematical event. Repair is not
synonymous with restoration, symptom relief, renewed coordination,
forgiveness, reconciliation, or return to a prior baseline. Some losses
are irreversible; truthful repair may require changed identity,
boundaries, restitution, accommodation, or preservation of memory.

\subsubsection{Operational consequence}\label{operational-consequence-6}

A failure or repair claim must state:

\begin{itemize}
\tightlist
\item
  the protected functions or standing;
\item
  the affected parties;
\item
  the failure criterion;
\item
  the intervention and disturbance classes;
\item
  the recovery horizon;
\item
  what cannot be restored;
\item
  and how recurrence will be assessed.
\end{itemize}

\subsubsection{Failure modes}\label{failure-modes-6}

\begin{itemize}
\tightlist
\item
  treating discomfort, dissent, or divergence as collapse;
\item
  romanticizing trauma as necessary growth;
\item
  confusing surface stability with repair;
\item
  requiring reconciliation where separation is protective;
\item
  declaring irreversibility without an intervention class;
\item
  and using repair language without remedy or prevention of recurrence.
\end{itemize}

\begin{center}\rule{0.5\linewidth}{0.5pt}\end{center}

\subsection{CU-FP8 --- Identity as Constraint-Preserved
Continuity}\label{cu-fp8-identity-as-constraint-preserved-continuity}

\textbf{Authority class:} FRAMEWORK-CANONICAL as a schema\\
\textbf{Claim status:} {[}Definition{]} and {[}Bridge principle{]}

\subsubsection{Statement}\label{statement-7}

An identity criterion specifies which relations, capacities, histories,
boundaries, or constraints must remain invariant---or remain connected
by an admissible continuity path---for a system to count as the same
system through change.

A general form is

\[\adjustbox{max width=\linewidth}{$\displaystyle 
x_t\equiv_{K_I,H}x_{t'}
\quad\Longleftrightarrow\quad
\\text{an admissible path preserves identity-relevant constraints }K_I
\\text{ over history }H.
$}\]

There is no single identity criterion for physical objects, organisms,
persons, institutions, artificial agents, or candidate subjects.
Self-modeling, self-maintenance, egoic persistence, phenomenal
perspective, moral patiency, responsibility, and sovereignty are
distinct.

\subsubsection{Controlling
subprinciples}\label{controlling-subprinciples}

\subsubsection{CU-FP8a --- Identity-Relevant Constraint
Set}\label{cu-fp8a-identity-relevant-constraint-set}

\textbf{Authority class:} FRAMEWORK-CANONICAL SUBPRINCIPLE

A claim that a system persists as the same system must identify a
nontrivial set of relations, capacities, boundaries, commitments, or
continuities whose loss would defeat that identity judgment.

\subsubsection{CU-FP8b --- Continuity Across Component
Change}\label{cu-fp8b-continuity-across-component-change}

\textbf{Authority class:} FRAMEWORK-CANONICAL SUBPRINCIPLE

Token identity may survive replacement of components when an admissible
continuity process preserves the identity-relevant organization. The
relevant process is domain-specific and must distinguish continuation,
copying, branching, reconstruction, and replacement.

\subsubsection{CU-FP8c --- Historical
Dependence}\label{cu-fp8c-historical-dependence}

\textbf{Authority class:} FRAMEWORK-CANONICAL SUBPRINCIPLE

An identity claim must state whether history is constitutive, merely
causally relevant, compressed into the current state, or recoverable
only through an explicit history functional. Similar present states need
not have the same identity when their continuity histories differ.

\subsubsection{CU-FP8d --- Endogenous
Maintenance}\label{cu-fp8d-endogenous-maintenance}

\textbf{Authority class:} FRAMEWORK-CANONICAL SUBPRINCIPLE

Some robust identities are partly constituted by the system's own
participation in maintaining the conditions of its continuation.
Endogenous maintenance is not required for every object or momentary
perspective and must not be inferred from persistence alone.

\subsubsection{CU-FP8e --- Transformation and Termination
Criteria}\label{cu-fp8e-transformation-and-termination-criteria}

\textbf{Authority class:} FRAMEWORK-CANONICAL SUBPRINCIPLE

An identity account must state which changes count as continued
identity, transformation, branching, replacement, copying, suspension,
or termination. The criteria may be purpose-relative, but the practical
purpose and parties affected by the judgment must be disclosed.

\subsubsection{Operational consequence}\label{operational-consequence-7}

Any identity claim must state:

\begin{itemize}
\tightlist
\item
  the token or type being reidentified;
\item
  identity-relevant constraints;
\item
  permitted transformations;
\item
  role of history;
\item
  endogenous maintenance requirements, if any;
\item
  interruption, copying, branching, and migration criteria;
\item
  termination criteria;
\item
  and the practical decision the identity judgment serves.
\end{itemize}

\subsubsection{Failure modes}\label{failure-modes-7}

\begin{itemize}
\tightlist
\item
  equating memory with identity;
\item
  assuming substrate independence without a continuity criterion;
\item
  treating copied information as automatically the same subject;
\item
  equating self-modeling with persistent ego;
\item
  and inferring responsibility or sovereignty from causal unity alone.
\end{itemize}

\begin{center}\rule{0.5\linewidth}{0.5pt}\end{center}

\subsection{CU-FP9 --- Discovery, Construction, and
Corrigibility}\label{cu-fp9-discovery-construction-and-corrigibility}

\textbf{Registry status:} CANONICAL\\
\textbf{Claim status:} {[}Methodological principle{]}

\subsubsection{Statement}\label{statement-8}

Systems encounter constraints they do not invent, while constructing
models, measures, categories, and practices through which those
constraints become intelligible. Responsible inquiry distinguishes
discovered resistance from constructed representation and remains
corrigible when the representation fails.

\subsubsection{Scope}\label{scope-7}

This principle replaces the stronger claim that ``coherence is
discovered rather than invented'' without qualification. Constraints can
be real while operationalizations are conventional or purpose-relative.
Moral and metaphysical conclusions do not follow merely because some
constraints are discovered.

\subsubsection{Operational consequence}\label{operational-consequence-8}

For any measure or principle, state which elements are:

\begin{itemize}
\tightlist
\item
  imposed by the target system or evidence;
\item
  selected by the modeler;
\item
  conventional;
\item
  normative;
\item
  or unresolved.
\end{itemize}

\subsubsection{Failure modes}\label{failure-modes-8}

\begin{itemize}
\tightlist
\item
  naïve realism about a chosen metric;
\item
  relativism that treats every constraint as constructed;
\item
  presenting a normative choice as a discovery of physics;
\item
  protecting a framework by redefining failure as deeper confirmation.
\end{itemize}

\begin{center}\rule{0.5\linewidth}{0.5pt}\end{center}

\section{Part II --- Method and Claim
Governance}\label{part-ii-method-and-claim-governance}

\subsection{CU-M1 --- Claim
Classification}\label{cu-m1-claim-classification}

\textbf{Registry status:} CANONICAL

\subsubsection{Statement}\label{statement-9}

Every load-bearing statement must be classified as definition, formal
result, conditional result, imported result, empirical claim, bridge
principle, hypothesis, interpretation, commitment, or open problem.

\subsubsection{Operational consequence}\label{operational-consequence-9}

Where one sentence contains multiple statuses, decompose it. A
mathematical derivation may be exact while its physical application is
conditional and its metaphysical reading interpretive.

\begin{center}\rule{0.5\linewidth}{0.5pt}\end{center}

\subsection{CU-M2 --- Explicit Bridges and Translation
Licenses}\label{cu-m2-explicit-bridges-and-translation-licenses}

\textbf{Authority class:} FRAMEWORK-CANONICAL\\
\textbf{Claim status:} {[}Methodological principle{]}

\subsubsection{Statement}\label{statement-10}

Whenever an argument crosses domains, levels, model classes, or
discourse registers, the bridge premise must be stated and separately
defended. Where a mathematical, formal, or empirical structure is
transferred, the analysis must identify the strongest Translation
License actually earned.

CU distinguishes five license levels:

\begin{enumerate}
\def\labelenumi{\arabic{enumi}.}
\tightlist
\item
  \textbf{Level 0 --- Analogy.} The source and target share suggestive
  language, imagery, or pattern. No theorem, mechanism, metric, or
  evidential status transfers.
\item
  \textbf{Level 1 --- Shared schema.} The source and target can be
  organized by the same questions or tuple structure. The analytic frame
  transfers; quantitative conclusions do not.
\item
  \textbf{Level 2 --- Property-preserving transfer.} A bounded result
  transfers because the target object satisfies declared properties
  required by that result, such as convexity, separability,
  monotonicity, viability assumptions, Markov structure, partial
  ordering, or controllability conditions.
\item
  \textbf{Level 3 --- Exact structural transfer.} The target object
  belongs to the theorem's actual mathematical class and satisfies its
  stated hypotheses. Exact identities, representation results, or
  invariance theorems may transfer within that scope.
\item
  \textbf{Level 4 --- Empirically validated instantiation.} The
  transferred structure has demonstrated explanatory, predictive,
  discriminative, measurement, or intervention value in target-domain
  data.
\end{enumerate}

\subsubsection{Scope}\label{scope-8}

A lower license is not a failed higher license. Analogy and shared
schema may be intellectually valuable when their status remains visible.
Exact formal fit does not establish that the model describes reality,
that its variables possess the proposed interpretation, or that its use
is normatively legitimate.

The levels do not replace the description--commitment distinction. A
transition from descriptive result to interpretation, norm, or
commitment still requires its own bridge and remains governed by claim
status and register discipline.

\subsubsection{Operational
consequence}\label{operational-consequence-10}

Every cross-domain transfer must identify:

\begin{itemize}
\tightlist
\item
  the source construct;
\item
  the target construct;
\item
  the declared license level;
\item
  the structure preserved;
\item
  the new assumptions introduced;
\item
  the information or distinctions lost;
\item
  the evidence available;
\item
  the failure and narrowing conditions;
\item
  the claim-status change, if any;
\item
  and the relevant review gate.
\end{itemize}

No conclusion inherits the status of its source merely through analogy,
shared notation, formal resemblance, or thematic fit. A transfer to a
new population, institution, scale, horizon, or purpose may require a
new license even within the same named domain.

\subsubsection{Failure modes}\label{failure-modes-9}

\begin{itemize}
\tightlist
\item
  treating analogy as identity;
\item
  treating a shared tuple as a shared mechanism;
\item
  applying a theorem to an object outside its mathematical class;
\item
  importing a metric because its notation is familiar;
\item
  treating exact mathematics as empirical validation;
\item
  treating empirical fit as ontological proof;
\item
  and deriving a normative or spiritual conclusion from a descriptive
  transfer without an explicit bridge.
\end{itemize}

\subsubsection{Review gate}\label{review-gate-1}

Source-domain expertise, target-domain expertise, formal-method review
where applicable, and affected-community review proportionate to
consequential reach.

\subsubsection{Related principles}\label{related-principles-1}

CU-M1; CU-M3; CU-M4; CU-M6; CU-M7; CU-M10; CU-M12; CU-M14; CU-L2; CU-L3.

\begin{center}\rule{0.5\linewidth}{0.5pt}\end{center}

\subsection{CU-M3 --- Minimal
Commitment}\label{cu-m3-minimal-commitment}

\textbf{Registry status:} CANONICAL

\subsubsection{Statement}\label{statement-11}

Introduce no stronger ontological, mathematical, empirical, or normative
claim than the explanatory step requires.

\subsubsection{Operational
consequence}\label{operational-consequence-11}

Use temporal order without assuming an arrow; compatibility without
claiming derivation; partial order without forced scalarization;
interpretation without descriptive certainty.

\begin{center}\rule{0.5\linewidth}{0.5pt}\end{center}

\subsection{CU-M4 --- Scope and Boundary
Declaration}\label{cu-m4-scope-and-boundary-declaration}

\textbf{Registry status:} CANONICAL

\subsubsection{Statement}\label{statement-12}

Every substantive model or evaluation must identify its system boundary,
scale, horizon, observer or access map, and excluded variables.

\subsubsection{Operational
consequence}\label{operational-consequence-12}

A result cannot be generalized beyond those conditions without an
explicit argument.

\begin{center}\rule{0.5\linewidth}{0.5pt}\end{center}

\subsection{CU-M5 --- Assumption and Dependency
Ledger}\label{cu-m5-assumption-and-dependency-ledger}

\textbf{Registry status:} CANONICAL

\subsubsection{Statement}\label{statement-13}

Every manuscript must identify the assumptions and upstream principles
supporting each central claim, together with downstream claims that
would require revision if it fails.

\subsubsection{Operational
consequence}\label{operational-consequence-13}

The corpus is maintained as a dependency graph rather than a collection
of mutually reinforcing citations.

\begin{center}\rule{0.5\linewidth}{0.5pt}\end{center}

\subsection{CU-M6 --- Reproducibility, Provenance, and Observation
Discipline}\label{cu-m6-reproducibility-provenance-and-observation-discipline}

\textbf{Authority class:} FRAMEWORK-CANONICAL

\subsubsection{Statement}\label{statement-14}

Computational, empirical, source-dependent, and feedback-dependent
claims must retain enough information for independent checking: data and
source provenance, code or procedure, parameters, units, tolerances,
versioned sources, observation windows, omitted populations, attribution
rules, feedback delay, missing outcomes, known limitations, and the
distinction between generated rationales and causally faithful
explanations.

\subsubsection{Operational
consequence}\label{operational-consequence-14}

Lost numerical work must be reconstructed rather than cited as an
inaccessible result. Simulations establish properties of the model
tested, not reality without a bridge. A feedback result is incomplete
when the observed population, delay, attribution, or missing outcomes
are not declared. Generated explanations are arguments or interface
artifacts unless causal faithfulness is independently established.

\begin{center}\rule{0.5\linewidth}{0.5pt}\end{center}

\subsection{CU-M7 --- Failure and Narrowing
Conditions}\label{cu-m7-failure-and-narrowing-conditions}

\textbf{Registry status:} CANONICAL

\subsubsection{Statement}\label{statement-15}

Every central hypothesis, bridge, or interpretation must state what
would narrow, reclassify, or defeat it.

\subsubsection{Operational
consequence}\label{operational-consequence-15}

A claim that can absorb every outcome without loss is not functioning as
an empirical or discriminating theoretical claim.

\begin{center}\rule{0.5\linewidth}{0.5pt}\end{center}

\subsection{CU-M8 --- Reflexive Loss}\label{cu-m8-reflexive-loss}

\textbf{Registry status:} CANONICAL

\subsubsection{Statement}\label{statement-16}

Revision must be capable of producing genuine loss: withdrawal of a
theorem, retirement of a principle, supersession of a paper, or
reclassification of a scientific claim as interpretation.

\subsubsection{Operational
consequence}\label{operational-consequence-16}

The framework may not call itself corrigible while preserving every
conclusion under new wording.

\begin{center}\rule{0.5\linewidth}{0.5pt}\end{center}

\subsection{CU-M9 --- Versioning and Historical
Trace}\label{cu-m9-versioning-and-historical-trace}

\textbf{Registry status:} CANONICAL

\subsubsection{Statement}\label{statement-17}

Canonical changes must be versioned, dated, explained, and propagated
through the dependency registry while preserving the superseded form for
audit.

\subsubsection{Operational
consequence}\label{operational-consequence-17}

Publication artifacts include a change ledger and status registry.

\begin{center}\rule{0.5\linewidth}{0.5pt}\end{center}

\subsection{CU-M10 --- Burden Proportional to Reach, Irreversibility,
and Affected
Scope}\label{cu-m10-burden-proportional-to-reach-irreversibility-and-affected-scope}

\textbf{Authority class:} FRAMEWORK-CANONICAL

\subsubsection{Statement}\label{statement-18}

The evidential, argumentative, participatory, and governance burden
increases with a claim or system's conceptual reach, causal reach,
irreversibility, affected scope, concentration of power, and capacity to
externalize error.

A useful but incomplete ordering is

\[\adjustbox{max width=\linewidth}{$\displaystyle 
\\text{definition}
<
\\text{formal construction}
<
\\text{conditional recovery}
<
\\text{empirical generalization}
<
\\text{cross-domain law}
<
\\text{fundamental ontology}.
$}\]

\subsubsection{Operational
consequence}\label{operational-consequence-18}

Exact mathematics does not discharge the burden of universal
application. High benchmark performance does not establish
authorization. A system with broad causal reach requires stronger
evidence, containment, independent review, affected-party standing, and
practical correction than a bounded low-stakes tool.

\begin{center}\rule{0.5\linewidth}{0.5pt}\end{center}

\subsection{CU-M11 --- The Coherence Analysis
Protocol}\label{cu-m11-the-coherence-analysis-protocol}

\textbf{Authority class:} FRAMEWORK-CANONICAL\\
\textbf{Claim status:} {[}Methodological principle{]} and {[}Governance
rule{]}

\subsubsection{Statement}\label{statement-19}

Every substantive CU model, comparison, evaluation, or consequential
application must complete a \textbf{Coherence Analysis Protocol (CAP)}
before its result is treated as portable, action-guiding, or
authoritative.

CAP has five declaration modules.

\paragraph{A. System declaration}\label{a.-system-declaration}

State:

\begin{itemize}
\tightlist
\item
  the unit or candidate system;
\item
  the boundary and relevant environment;
\item
  the scale, partition, or access map;
\item
  the temporal horizon;
\item
  the state, trajectory, process, or organization under study;
\item
  the identity or continuation criterion, where applicable;
\item
  the represented variables and relations;
\item
  admissible actions, transitions, or interventions;
\item
  and material exclusions or idealizations.
\end{itemize}

\paragraph{B. Observation declaration}\label{b.-observation-declaration}

State:

\begin{itemize}
\tightlist
\item
  what is directly observed;
\item
  what is inferred, proxied, imputed, latent, delayed, or missing;
\item
  the instruments, records, reports, or access channels used;
\item
  omitted populations, outcomes, and time periods;
\item
  attribution rules and feedback delays;
\item
  who controls the records and categories;
\item
  and what evidence could correct the model.
\end{itemize}

\paragraph{C. Evaluation declaration}\label{c.-evaluation-declaration}

State:

\begin{itemize}
\tightlist
\item
  the profile, quantity, ordering, viable region, threshold family, risk
  functional, or trajectory property evaluated;
\item
  the protected conditions or floors;
\item
  affected beings and groups;
\item
  the presented option set, claimed feasible set, and materially omitted
  alternatives;
\item
  benefit, burden, risk, and repair-cost distributions;
\item
  the screening rule;
\item
  the comparison rule among options that pass;
\item
  the scalar type and parameters, if any;
\item
  comparisons that remain incomplete;
\item
  and the local--global relation status.
\end{itemize}

\paragraph{D. Power and authorship
declaration}\label{d.-power-and-authorship-declaration}

State:

\begin{itemize}
\tightlist
\item
  who framed the question;
\item
  who selected the categories and evidence channels;
\item
  who generated, excluded, and classified options;
\item
  who set the floors, weights, horizon, and stopping rule;
\item
  who benefits and who bears consequences;
\item
  who may challenge, refuse, exit, appeal, stop, reverse, or obtain
  remedy;
\item
  and whether review is sufficiently independent of the system or
  institution under review.
\end{itemize}

\paragraph{E. Revision declaration}\label{e.-revision-declaration}

State:

\begin{itemize}
\tightlist
\item
  what would falsify, narrow, or reclassify the model;
\item
  which changes count as new evidence within the same model;
\item
  which changes constitute model revision;
\item
  how adverse outcomes will be detected;
\item
  how records, harms, and recurrent failure will be corrected, repaired,
  and tracked;
\item
  how categories, floors, options, and governance may change;
\item
  and how superseded versions will remain auditable.
\end{itemize}

\subsubsection{Scope}\label{scope-9}

CAP names and integrates requirements already distributed across CU-FP2,
CU-M4--CU-M10, CU-D9, CU-G1--CU-G7, and CU-L5. It does not replace
stronger domain-specific requirements or review gates.

Declaration makes assumptions visible. It does not by itself make a
model true, complete, legitimate, contestable, or safe. Publicity
without standing is not contestability. Feedback without authority to
alter outcomes is not correction.

\subsubsection{Operational
consequence}\label{operational-consequence-19}

Every reader-facing domain paper must include a compact CAP record or a
clearly identified cross-reference to one. Every consequential
application must preserve a versioned CAP record with the model and
decision artifact.

Unresolved fields must be labeled \textbf{not currently specified},
\textbf{not measured}, \textbf{not applicable}, or \textbf{contested}
rather than filled by analogy. The burden of detail and independent
review rises with causal reach, irreversibility, affected scope, and
concentration of power.

\subsubsection{Failure modes}\label{failure-modes-10}

\begin{itemize}
\tightlist
\item
  checklist compliance without substantive declaration;
\item
  omitting affected parties or future burdens;
\item
  treating the presented option set as the feasible set;
\item
  declaring weights while hiding who authorized them;
\item
  public disclosure without challenge, refusal, exit, reversal, or
  remedy;
\item
  provider-controlled self-certification;
\item
  changing floors, categories, horizons, or option sets while describing
  the change as a mere update;
\item
  and treating a fully declared model as fully justified.
\end{itemize}

\subsubsection{Review gate}\label{review-gate-2}

Cross-domain methods, owning-domain expertise, governance and ethics
review, and affected-community participation proportionate to the
model's reach and risk.

\subsubsection{Related principles}\label{related-principles-2}

CU-FP2; CU-M4; CU-M5; CU-M6; CU-M7; CU-M9; CU-M10; CU-M12; CU-M13;
CU-M14; CU-D9; CU-G1--CU-G7; CU-L5.

\begin{center}\rule{0.5\linewidth}{0.5pt}\end{center}

\subsection{CU-M12 --- The Shared Analytic Frame and Honest
Incompleteness}\label{cu-m12-the-shared-analytic-frame-and-honest-incompleteness}

\textbf{Authority class:} FRAMEWORK-CANONICAL\\
\textbf{Claim status:} {[}Methodological principle{]}

\subsubsection{Statement}\label{statement-20}

CU unifies domains through recurring analytic questions, not by assuming
that every domain contains the same state object, scalar, geometry,
dynamics, equilibrium, or optimum.

Each substantive domain model must answer, where meaningful:

\begin{enumerate}
\def\labelenumi{\arabic{enumi}.}
\tightlist
\item
  \textbf{Unit and boundary} --- What is the candidate system, process,
  or governed object?
\item
  \textbf{State or trajectory space} --- What can vary?
\item
  \textbf{Environment} --- What coupled conditions lie outside the focal
  boundary?
\item
  \textbf{Constraints and viable region} --- What must remain possible,
  protected, or identity-preserving?
\item
  \textbf{Actions, options, and transitions} --- What choices or changes
  can occur, which alternatives are omitted, and through whose agency?
\item
  \textbf{Dynamics or update rule} --- What drives change, if specified?
\item
  \textbf{Observation and access} --- What can be measured, known,
  reported, or recovered?
\item
  \textbf{Perturbation and intervention class} --- Which disturbances or
  interventions are in scope?
\item
  \textbf{Equilibria, basins, or regimes} --- Which are actually
  demonstrated?
\item
  \textbf{Local--global status} --- Exact decomposition, monotonicity or
  inequality, coordinate-wise profile relation, or qualitative analysis?
\item
  \textbf{Scalar status} --- None, measurement scalar, decision
  scalarization, risk functional, or other?
\item
  \textbf{Statistical-image, allocation, refinement, and metric status}
  --- Which quantities, if any, receive CU-M14 licenses?
\item
  \textbf{Translation status} --- Which CU-M2 license supports each
  imported structure?
\item
  \textbf{Open gaps} --- Which slots remain unspecified, unmeasured,
  inapplicable, or contested?
\end{enumerate}

\subsubsection{Scope}\label{scope-10}

The repeated frame licenses common questions, not common answers. A
physical probability model, biological viability system, psychological
profile, social institution, ethical evaluation, and
artificial-intelligence deployment may require different objects and
evidential standards.

An empty or unresolved slot is methodologically informative. Ethics need
not possess a natural flow equation. Psychology need not possess one
general phase portrait. Social Dynamics may identify regimes without
demonstrating universal basins. A companion paper may instantiate a
formal evaluation model rather than a physical state space.

\subsubsection{Operational
consequence}\label{operational-consequence-20}

Each reader-facing domain paper must include a visually standardized
\textbf{Domain Analytic Frame} card near its initial declaration
section. Companion papers must use the card where it clarifies their
formal object and may mark genuinely inapplicable slots.

Claims about attractors or basins require a declared state space,
transition or dynamics, perturbation class, and evidence of return or
convergence. Claims about optimization require a declared objective and
admissible set. Claims about dynamics require more than metaphorical
change over time.

\subsubsection{Failure modes}\label{failure-modes-11}

\begin{itemize}
\tightlist
\item
  forcing every domain into a dynamical vocabulary;
\item
  filling blank slots with analogy;
\item
  treating a shared question as evidence of one shared mechanism;
\item
  using attractor language without state, transition, perturbation, and
  basin evidence;
\item
  treating a static profile as a probability distribution;
\item
  and concealing open problems to preserve visual uniformity.
\end{itemize}

\subsubsection{Review gate}\label{review-gate-3}

Corpus-wide conceptual review and owning-domain review for every
populated card.

\subsubsection{Related principles}\label{related-principles-3}

CU-FP1; CU-FP2; CU-M2; CU-M3; CU-M4; CU-M11; CU-M13; CU-M14; CU-D1;
CU-D4.

\begin{center}\rule{0.5\linewidth}{0.5pt}\end{center}

\subsection{CU-M13 --- Local--Global Declaration and Relation
Status}\label{cu-m13-localglobal-declaration-and-relation-status}

\textbf{Authority class:} FRAMEWORK-CANONICAL\\
\textbf{Claim status:} {[}Methodological principle{]}, with
model-specific {[}Formal result{]}, {[}Imported result{]}, {[}Bridge
principle{]}, or {[}Qualitative analysis{]}

\subsubsection{Statement}\label{statement-21}

Every claim of improvement, stability, coherence, welfare, performance,
safety, or repair must identify the focal scale, containing system,
partition or aggregation map, affected units, temporal horizon,
externalized costs, and observation limits.

A local result does not silently establish a global result. A global or
aggregate result does not silently establish a local, subgroup,
within-person, or component-level result.

The analysis must declare the strongest local--global relation actually
licensed:

\begin{enumerate}
\def\labelenumi{\arabic{enumi}.}
\tightlist
\item
  \textbf{Exact decomposition.} A total, between-part, and within-part
  identity may be stated only when the target quantity belongs to a
  proved decomposable class under the declared partition and weighting.
  The source theorem, assumptions, and equality conditions must be
  stated. Generalized-entropy decompositions are imported results and
  must be attributed accordingly (Theil, 1967; Shorrocks, 1980).
\item
  \textbf{Monotonicity or inequality.} A measure-specific inequality may
  be stated only when the functional and map satisfy the required
  conditions. For example, a genuine averaging map cannot increase a
  declared convex separable dispersion functional when the required
  conditions hold. Any interpretation as improvement, degradation,
  masking, or harm requires a further bridge.
\item
  \textbf{Coordinate-wise profile relation.} A profile does not acquire
  one exact decomposition because one or more coordinates do. Each
  coordinate must be licensed separately, and unresolved tradeoffs must
  remain visible.
\item
  \textbf{Qualitative or structural analysis.} Where no scalar theorem
  is available, the analysis may identify aggregate masking, cost
  transfer, scale-dependent reversal, subgroup deterioration, delayed
  burden, or repair asymmetry without implying a quantitative identity.
\end{enumerate}

\subsubsection{Scope}\label{scope-11}

Mathematical aggregation and normative noncompensation are distinct. A
theorem may describe how a quantity changes under averaging or
partition. It does not decide which losses are morally compensable,
which beings possess standing, or which floors may not be crossed.

CU-G1 governs protected conditions and distribution. CU-D9 governs
multiscale failure and externalization. CU-M13 governs declaration of
the mathematical or qualitative relation. These principles are mutually
supporting and non-redundant.

\subsubsection{Operational
consequence}\label{operational-consequence-21}

Every Domain Analytic Frame card must declare its local--global status.
Exact and inequality claims must identify the mathematical object,
partition, weighting, map, theorem source, assumptions, and
interpretation bridge. Profile-level claims must report coordinates
separately unless a justified aggregation is supplied.

No aggregate or higher-scale improvement establishes screening
admissibility when a relevant local unit crosses a protected floor.

\subsubsection{Failure modes}\label{failure-modes-12}

\begin{itemize}
\tightlist
\item
  inferring local improvement from an aggregate mean;
\item
  inferring global improvement from a focal-system gain;
\item
  hiding subgroup or future harm through aggregation;
\item
  importing exact decomposition into an arbitrary profile;
\item
  treating coordinate-wise results as one profile theorem;
\item
  describing a coarse view as improved without stating what the measure
  means;
\item
  deriving moral noncompensation from convexity, information geometry,
  or a physical theorem;
\item
  and defining the boundary to exclude the receiver of an externalized
  cost.
\end{itemize}

\subsubsection{Review gate}\label{review-gate-4}

Formal or statistical review for quantitative relations; owning-domain
review for interpretation; ethics, governance, and affected-community
review for protected conditions and consequential use.

\subsubsection{Related principles}\label{related-principles-4}

CU-FP2; CU-FP4; CU-M2; CU-M4; CU-M11; CU-M12; CU-M14; CU-D9; CU-G1;
CU-G2; CU-G7.

\begin{center}\rule{0.5\linewidth}{0.5pt}\end{center}

\subsection{CU-M14 --- Distributional, Allocation, and
Information-Geometric
License}\label{cu-m14-distributional-allocation-and-information-geometric-license}

\textbf{Authority class:} FRAMEWORK-CANONICAL\\
\textbf{Claim status:} {[}Methodological principle{]} and model-specific
{[}Bridge principle{]}

\subsubsection{Statement}\label{statement-22}

A system, profile, or set of numerical coordinates does not inherit
probability-simplex or information-geometric structure merely because
its values can be normalized. Every use must begin by declaring the
question, sample or category space, partition or observation map, weight
interpretation, omitted information, and target-domain purpose.

A candidate receives a \textbf{Statistical Image License} only when
nonnegative normalized weights genuinely represent probabilities,
frequencies, occupancies, policies, beliefs, or another defended
statistical object over categories that are exclusive and exhaustive
relative to the declared question.

A candidate receives an \textbf{Allocation License} only when, in
addition, the coordinates divide a real or declared common-unit total
without double counting and normalization preserves the phenomenon under
study.

A candidate receives a \textbf{Refinement and Coarse-Graining License}
only when category splitting, merging, or stochastic mapping has
operational meaning, category semantics remain sufficiently stable,
support changes are handled, and the map's role as observation,
aggregation, intervention, dissipation, or destruction is stated.

A licensed classical statistical model receives a \textbf{Fisher-Rao
License} only under the regularity and sufficient-statistic or
Markov-invariance assumptions relevant to the theorem invoked. Under
those assumptions, the Fisher-Rao metric is canonical up to scale
(Čencov, 1982; Campbell, 1986; Ay et al., 2015; Bauer et al., 2016).
This canonality does not select the partition, factorization, reference
state, divergence, functional, dynamics, evaluative direction, or
cross-domain trade rate.

Every further use must declare the selected metric convention,
divergence or comparison family, functional or objective, dynamics or
update rule, boundary treatment, Translation License level, evidence
status, and failure conditions. Multidimensional profiles do not become
distributions through normalization, quantum state spaces do not inherit
a unique classical Fisher-Rao metric, and no statistical geometry by
itself establishes substantive coherence, phenomenal character, moral
value, legitimacy, or a preferred direction of change (Petz \& Sudár,
1996).

\subsubsection{Scope}\label{scope-12}

The principle governs probability distributions, frequency
distributions, policies, beliefs, occupancy measures, and common-unit
allocations. It does not require every CU domain to use probability or
information geometry. A domain may contain licensed statistical
subobjects while its overall coherence object remains a profile, viable
region, partial ordering, causal structure, or trajectory property.

\subsubsection{Operational record}\label{operational-record}

Every invocation must record:

\begin{enumerate}
\def\labelenumi{\arabic{enumi}.}
\tightlist
\item
  target system and question;
\item
  sample space or category set;
\item
  partition or observation map;
\item
  weight interpretation;
\item
  statistical-image status;
\item
  allocation status, where relevant;
\item
  refinement and coarse-graining maps;
\item
  support and boundary treatment;
\item
  metric and scale convention;
\item
  divergence, reference family, or evaluator;
\item
  functional and dynamics, if any;
\item
  Translation License level;
\item
  evidence and validation status;
\item
  omitted information and failure conditions.
\end{enumerate}

Every Domain Analytic Frame card must separately state its Statistical
Image, Allocation, Refinement and Coarse-Graining, and Fisher-Rao status
and distinguish every licensed subobject from the domain's broader
profile.

\subsubsection{Failure modes}\label{failure-modes-13}

\begin{itemize}
\tightlist
\item
  normalizing an arbitrary profile and calling it a distribution;
\item
  treating a statistical image as the whole system;
\item
  assuming a conserved or divisible total without evidence;
\item
  double counting one unit across categories;
\item
  changing category semantics under refinement;
\item
  invoking Fisher-Rao outside a licensed classical statistical model;
\item
  extending classical uniqueness to quantum state space without an added
  criterion;
\item
  inferring a divergence, functional, natural-gradient flow, value
  direction, or physical law from the metric alone;
\item
  comparing numerical values across domains without a separate
  commensurability argument;
\item
  treating statistical dependence as causal integration, autonomy,
  consciousness, or moral standing.
\end{itemize}

\subsubsection{Review gate}\label{review-gate-5}

Formal review for the theorem and metric; statistical review for the
model and transformations; domain review for the representation and
interpretation; ethics, governance, and affected-party review for
consequential uses.

\subsubsection{Related principles}\label{related-principles-5}

CU-FP2; CU-M2; CU-M4; CU-M6; CU-M11; CU-M12; CU-M13; CU-D4; CU-G1.

\begin{center}\rule{0.5\linewidth}{0.5pt}\end{center}

\subsection{CU-M15 --- Typed Mandala, Shared Services, and Adjudicated
Dependency}\label{cu-m15-typed-mandala-shared-services-and-adjudicated-dependency}

\textbf{Authority class:} FRAMEWORK-CANONICAL\\
\textbf{Claim status:} {[}Methodological and corpus-governance
principle{]}

\subsubsection{Statement}\label{statement-23}

The sixty-four-stage Coherence Mandala is the canonical philosophical
and contemplative traversal of the framework. Its technical counterpart
is a \textbf{typed multigraph}, not one homogeneous ladder.

A Mandala stage may map to zero, one, or several technical objects.
Several stages may invoke one shared object or service. Technical
objects may be operational capacities, open problems, representations,
evaluators, methodological bridges, normative relations, governance
mechanisms, interpretations, commitments, diagnostics, roots, or summary
joins. The absence of a distinct operational capacity at a stage is an
admissible result.

Every proposed relation must be semantically typed. CU distinguishes
strict prerequisite, conditional prerequisite, alternative enablement,
enabling, constitutive, empirical or bridge, representation or
implementation, measurement or evaluation, presupposition, governance or
authorization, protection, feedback or correction, service invocation,
interpretive mapping, commitment mapping, summary membership, alias or
migration, and external prerequisite relations.

A \textbf{strict proof-order subgraph} contains only individually
adjudicated strict prerequisites. It must be acyclic relative to the
declared model. Acyclicity may not be secured by assigning edge
direction or type from Mandala stage order. Conditional, constitutive,
protective, service, feedback, interpretive, commitment, and summary
relations remain outside the strict proof-order test and may form
legitimate typed cycles.

The Mandala's reciprocal form joins:

\begin{enumerate}
\def\labelenumi{\arabic{enumi}.}
\tightlist
\item
  an enablement structure of capacities and realization conditions;
\item
  a protection-and-service layer through which standing, floors,
  provenance, correction, repair, irreversibility, and institutional
  learning preserve or govern capacities;
\item
  an interpretation-and-commitment layer through which finite
  participants understand and undertake responsible participation.
\end{enumerate}

Representations, observations, and evaluators operate across these
layers. No layer may be treated as a universal scalar ascent or as proof
that later stages possess greater causal power.

\subsubsection{Operational
consequence}\label{operational-consequence-22}

Every technical graph entry must declare object type, relation type,
scope and assumptions, necessity or sufficiency status, evidence or
argument appropriate to its type, failure or countermodel condition,
Translation License status where relevant, warrant basis, and whether
the relation belongs to strict proof order.

Shared services are defined once with typed internal roles.
Stage-specific uses are recorded as invocations carrying scale, affected
scope, authority, and failure conditions rather than duplicated as new
rungs.

\subsubsection{Failure modes}\label{failure-modes-14}

\begin{itemize}
\tightlist
\item
  treating stage numbering as chronology or proof order;
\item
  classifying edges by stage direction or target type rather than
  semantic review;
\item
  representing protective relations only as feedback;
\item
  counting evaluators, governance relations, interpretations,
  commitments, or joins as causal capacities;
\item
  collapsing cross-type services into homogeneous nodes;
\item
  using an untyped cycle as evidence of derivation;
\item
  treating a toy witness as target-domain empirical validation;
\item
  assuming that every canonical stage must add a distinct formal object.
\end{itemize}

\subsubsection{Review gate}\label{review-gate-6}

Revision requires coordinated review of the owning domains, graph
methodology, register discipline, affected-party and governance
implications, and the companion Mandala research supplement. The current
edge audit is an internal semantic adjudication, not target-domain
validation.

\subsubsection{Related principles}\label{related-principles-6}

CU-FP1; CU-FP2; CU-M2; CU-M11; CU-M12; CU-M13; CU-M14; CU-G1--CU-G7;
CU-L3; CU-L7.

\begin{center}\rule{0.5\linewidth}{0.5pt}\end{center}

\section{Part III --- Logic, Registers, and
Epistemology}\label{part-iii-logic-registers-and-epistemology}

\subsection{CU-L1 --- Truth--Coherence
Distinction}\label{cu-l1-truthcoherence-distinction}

\textbf{Registry status:} CANONICAL\\
\textbf{Claim status:} {[}Definition{]}

\subsubsection{Statement}\label{statement-24}

Truth-apt claims are true or false according to whether the world is as
they represent it within their semantics and scope. Internal coherence,
usefulness, meaning, or viability does not make a false descriptive
claim true.

\subsubsection{Consequence}\label{consequence}

Coherence is an epistemic virtue among others, not a replacement for
truth.

\begin{center}\rule{0.5\linewidth}{0.5pt}\end{center}

\subsection{CU-L2 --- Content, Force, and the Two
Registers}\label{cu-l2-content-force-and-the-two-registers}

\textbf{Registry status:} CANONICAL\\
\textbf{Claim status:} {[}Definition{]} and {[}Bridge principle{]}

\subsubsection{Statement}\label{statement-25}

CU distinguishes descriptive and commitment uses primarily through
speech-act force, principal function, inferential role, and standard of
assessment.

\begin{itemize}
\tightlist
\item
  The \textbf{description register} represents how things are and
  answers to truth, evidence, proof, calibration, and correction.
\item
  The \textbf{commitment register} undertakes, vows, trusts, refuses,
  consecrates, or orients action and answers to warrant, integrity,
  consequence, agency, correction, and repair.
\end{itemize}

The registers need not carry two exhaustive semantic kinds. Stable
proposition-like content may be asserted, promised, adopted as a plan,
or embedded in reasoning. Mixed utterances require decomposition, and
disputed principal function must be reported rather than forced.

\subsubsection{Boundary}\label{boundary}

The distinction does not make prayer, ritual, poetry, myth, silence, or
first-person avowal meaningless when they lie outside the formal Boolean
fragment.

\subsection{CU-L3 --- Asymmetric
Permeability}\label{cu-l3-asymmetric-permeability}

\textbf{Registry status:} CANONICAL

\subsubsection{Statement}\label{statement-26}

Established descriptive facts constrain commitments. Commitment warrant
does not establish descriptive truth.

\subsubsection{Consequence}\label{consequence-1}

No spiritual, moral, or practical benefit can certify a factual claim
without independent evidence. No lack of empirical proof automatically
renders a non-descriptive commitment meaningless.

\begin{center}\rule{0.5\linewidth}{0.5pt}\end{center}

\subsection{CU-L4 --- Fixed-Base Preferential Consequence over
Addressable Commitment
Content}\label{cu-l4-fixed-base-preferential-consequence-over-addressable-commitment-content}

\textbf{Registry status:} CANONICAL--CONDITIONAL\\
\textbf{Claim status:} {[}Definition{]} and {[}Imported formal result{]}

\subsubsection{Statement}\label{statement-27}

Let \(L_D\) have descriptive satisfaction \(\vDash_D\), and let \(L_R\)
be a declared Boolean-closed commitment-representation language with
realization relation \(\Vdash_R\). The formal result begins only after
stable realization conditions are supplied.

Within one evaluation model, let \(V_F\equiv V_{F,K,A,H}\) and \(\prec\)
remain fixed while premises filter by intersection:

\[\adjustbox{max width=\linewidth}{$\displaystyle 
W_{F,\Gamma}=V_F\cap\operatorname{Real}_R(\Gamma).
$}\]

Define:

\[\adjustbox{max width=\linewidth}{$\displaystyle 
\Gamma\mid\!\sim_C\varphi
\iff
\operatorname{Min}_{\prec}(W_{F,\Gamma})
\subseteq\llbracket\varphi\rrbracket_R.
$}\]

Under Boolean closure, definability, and smoothness, the relation
satisfies System P (Kraus, Lehmann, \& Magidor, 1990). If the model is
ranked and every relevant nonempty definable set attains a minimum, it
supports Rational Monotony (Lehmann \& Magidor, 1992).

\subsubsection{Scope}\label{scope-13}

KLM consequence governs inference among already interpreted contents. It
does not constitute their meaning or solve Frege--Geach. Formula-level
\textbf{Or} presupposes the declared commitment algebra. The result
supplies no normative content for the base, ordering, floors, affected
scope, or option set.

\subsection{CU-L5 --- Non-Monotonic Inference, Model Revision, and
Warrant
Non-Inheritance}\label{cu-l5-non-monotonic-inference-model-revision-and-warrant-non-inheritance}

\textbf{Registry status:} CANONICAL

\subsubsection{Statement}\label{statement-28}

Within a fixed preferential model, new premises may withdraw a
defeasible conclusion by changing the subset compatible with the premise
while leaving the base and ordering fixed.

A change to descriptive facts, protected floors, affected-agent scope,
horizon, viability criterion, realization model, or ordering is
\textbf{model revision}, not premise addition inside the same relation.

\subsubsection{Consequence}\label{consequence-2}

A revised model must be declared and re-evaluated. It does not inherit
the warrant of the model it replaces. Classical entailment, defeasible
consequence, premise-responsive selection, and model revision remain
distinct operations.

\subsection{CU-L6 --- Strong Warrant, Mere Admissibility, Exclusion,
Infeasibility, and
Underdetermination}\label{cu-l6-strong-warrant-mere-admissibility-exclusion-infeasibility-and-underdetermination}

\textbf{Registry status:} CANONICAL

\subsubsection{Statement}\label{statement-29}

Let:

\[\adjustbox{max width=\linewidth}{$\displaystyle 
B_{F,\Gamma}=\operatorname{Min}_{\prec}(V_F\cap\operatorname{Real}_R(\Gamma)).
$}\]

CU distinguishes:

\[\adjustbox{max width=\linewidth}{$\displaystyle 
\text{strong warrant: }
B_{F,\Gamma}\neq\varnothing
\land
B_{F,\Gamma}\subseteq\llbracket\varphi\rrbracket_R,
$}\]

\[\adjustbox{max width=\linewidth}{$\displaystyle 
\text{admissibility: }
B_{F,\Gamma}\neq\varnothing
\land
B_{F,\Gamma}\cap\llbracket\varphi\rrbracket_R\neq\varnothing,
$}\]

\[\adjustbox{max width=\linewidth}{$\displaystyle 
\text{exclusion: }
B_{F,\Gamma}\neq\varnothing
\land
B_{F,\Gamma}\cap\llbracket\varphi\rrbracket_R=\varnothing.
$}\]

\textbf{Mere admissibility} is admissibility without strong warrant.

\textbf{Infeasibility} obtains when
\(V_F\cap\operatorname{Real}_R(\Gamma)=\varnothing\).

\textbf{Underdetermination} has two forms. \textbf{Formal
non-attainment} occurs when the premise-compatible set is nonempty but
has no minimal elements under the declared ordering.
\textbf{Model-adequacy abstention} occurs before a best-set verdict is
issued because the represented facts, affected-party scope, options,
floors, horizon, observation model, or ordering remain unspecified,
contested, or inadequate. Ties or incomparability within an attained
best set do not themselves create underdetermination; they are handled
by strong warrant, mere admissibility, or exclusion.

\subsubsection{Consequence}\label{consequence-3}

Admissibility can fail conjunction and is not System P consequence.
Empty-set vacuity is never warrant. Classical negation is an adopted
choice within the Boolean fragment: exclusion of \(\varphi\) is
equivalent to strong warrant for \(\neg\varphi\), with the practical
flattening this entails.

\subsection{CU-L7 --- Anti-Rationalization and Justificatory
Integrity}\label{cu-l7-anti-rationalization-and-justificatory-integrity}

\textbf{Authority class:} FRAMEWORK-CANONICAL\\
\textbf{Claim status:} {[}Epistemic and ethical bridge principle{]}

\subsubsection{Statement}\label{statement-30}

A justification is defective when its operative standards are
selectively altered, its evidence channels narrowed, its alternatives
suppressed, or its inferential path reconstructed primarily to protect a
prior conclusion, identity, attachment, advantage, or policy.

\subsubsection{Scope}\label{scope-14}

Rationalization is not inferred from error or disagreement alone. It
requires evidence of asymmetric processing, selective standards,
avoidable evidence exclusion, option suppression, or identity-protective
distortion.

Justificatory integrity is necessary for defensible inquiry and moral
reasoning but is not sufficient for truth, permissibility, legitimacy,
justice, or adequate option generation.

\subsubsection{Correctability across
domains}\label{correctability-across-domains}

CU uses \textbf{correctability} for the capacity of a model, measure,
agent, or institution to remain answerable to relevant reality,
affected-party evidence, consequence, and repair. Loss of correctability
can be motivated or non-culpable. In motivated cases, a protected
identity, attachment, advantage, authority, objective, or policy shapes
what is admitted, measured, or counted; personal self-deception,
normative rationalization, silent optimization, metric capture, feedback
capture, and self-certified review are domain-specific forms. In
non-culpable cases, contact is lost through inherited protection,
missing observation, delayed feedback, underdetermination, or limited
capacity.

Error, disagreement, trauma, uncertainty, scarcity, and tragedy do not
by themselves establish self-deception. A characteristic failure occurs
when apparent improvement in a represented or measured condition
substitutes for improvement in the target conditions it was meant to
track. Correctability therefore requires not only accurate
representation but practical routes to challenge, model revision,
option-set reopening, repair, and nonrecurrence.

\begin{center}\rule{0.5\linewidth}{0.5pt}\end{center}

\subsection{CU-L8 --- Dynamic Preferential
Lifting}\label{cu-l8-dynamic-preferential-lifting}

\textbf{Registry status:} CANONICAL--CONDITIONAL

\subsubsection{Statement}\label{statement-31}

Trajectories may serve as the worlds of a preferential model.
Transformative commitments are represented by trajectories whose states,
capacities, and recovery profiles change through time while the
trajectory-level viability predicate and ordering remain fixed as model
parameters. Under the ordinary semantic and smoothness conditions,
System P---and System R for ranked models with attained minima---carries
over to that dynamic relation.

\subsubsection{Scope}\label{scope-15}

This result does not require one additive trajectory score and does not
eliminate horizon dependence. If the commitment changes the viability
criterion, protected floors, affected-agent scope, horizon, or ordering
itself, the comparison is between different models and the KLM
inheritance no longer follows automatically.

\begin{center}\rule{0.5\linewidth}{0.5pt}\end{center}

\subsection{CU-L9 --- Probability, Typicality, and
Preference}\label{cu-l9-probability-typicality-and-preference}

\textbf{Registry status:} CANONICAL

\subsubsection{Statement}\label{statement-32}

Probability, typicality, and preference are distinct structures. A
likely outcome need not be desirable; a preferred outcome need not be
likely; and a typicality rule must not hide rare harms relevant to
protected constraints.

\begin{center}\rule{0.5\linewidth}{0.5pt}\end{center}

\subsection{CU-L10 --- The Two-Stage
Limitation}\label{cu-l10-the-two-stage-limitation}

\textbf{Registry status:} CANONICAL\\
\textbf{Claim status:} {[}Formal negative result{]}

\subsubsection{Statement}\label{statement-33}

A procedure that selects premise-dependent policies from policy-induced
outcome distributions and then aggregates policy-indexed typical
consequences does not generally satisfy System P or System R. Mixed
policies can satisfy a disjunction by distributing support across its
disjuncts while satisfying neither disjunct individually; probability
thresholds can introduce additional failures but are not the source of
this Or obstruction. The feasibility map is therefore not generally a
Boolean homomorphism.

\subsubsection{Consequence}\label{consequence-4}

The two-stage construction is a modular decision architecture. Stronger
logical guarantees require additional proved restrictions.

\begin{center}\rule{0.5\linewidth}{0.5pt}\end{center}

\subsection{CU-Ep0 --- Constraint-Sensitive
Engagement}\label{cu-ep0-constraint-sensitive-engagement}

\textbf{Registry status:} CANONICAL

\subsubsection{Statement}\label{statement-34}

The pre-propositional basis of knowing is discriminating responsiveness
to constraint. Differences in the environment or system state must make
reliable differences to regulation or action.

\subsubsection{Limit}\label{limit}

Constraint sensitivity alone does not establish representation,
truth-apt belief, or consciousness.

\begin{center}\rule{0.5\linewidth}{0.5pt}\end{center}

\subsection{CU-Ep1 --- Situated Pattern and
Proto-Representation}\label{cu-ep1-situated-pattern-and-proto-representation}

\textbf{Registry status:} CANONICAL

\subsubsection{Statement}\label{statement-35}

Recurrent constraint-sensitive responses can stabilize into
expectations, maps, dispositions, and implicit models that guide action
across related situations.

\subsubsection{Limit}\label{limit-1}

Internal stability and successful habit are not yet descriptive truth.

\begin{center}\rule{0.5\linewidth}{0.5pt}\end{center}

\subsection{CU-Ep2 --- World-Coupled
Claims}\label{cu-ep2-world-coupled-claims}

\textbf{Registry status:} CANONICAL

\subsubsection{Statement}\label{statement-36}

Truth-apt belief begins when a representation is used as a claim about
how things are and is exposed to correction by what it represents.

\subsubsection{Limit}\label{limit-2}

Immediate usefulness does not by itself establish accuracy.

\begin{center}\rule{0.5\linewidth}{0.5pt}\end{center}

\subsection{CU-Ep3 --- Diachronic and Intersubjective
Correction}\label{cu-ep3-diachronic-and-intersubjective-correction}

\textbf{Registry status:} CANONICAL

\subsubsection{Statement}\label{statement-37}

Knowledge becomes more reliable when claims can be revised through time
and exposed to independent access, testimony, criticism, replication,
and preserved memory.

\subsubsection{Limit}\label{limit-3}

Consensus is not truth. Social processes add epistemic value only to the
extent that they preserve correction rather than manufacture conformity.

\begin{center}\rule{0.5\linewidth}{0.5pt}\end{center}

\subsection{CU-Ep4 --- Explicit Models and Formal
Systems}\label{cu-ep4-explicit-models-and-formal-systems}

\textbf{Registry status:} CANONICAL

\subsubsection{Statement}\label{statement-38}

At the modeling level, assumptions, variables, measurements, scope
conditions, and inference rules are made explicit. Empirical models
answer to observation; formal systems answer to proof relative to
axioms.

\subsubsection{Limit}\label{limit-4}

Proof does not establish empirical application. Predictive success does
not uniquely determine ontology.

\begin{center}\rule{0.5\linewidth}{0.5pt}\end{center}

\subsection{CU-Ep5 --- Epistemic
Governance}\label{cu-ep5-epistemic-governance}

\textbf{Authority class:} FRAMEWORK-CANONICAL

\subsubsection{Statement}\label{statement-39}

Epistemic governance regulates how claims, models, measurements,
feedback, options, and explanations are sourced, translated, tested,
revised, versioned, authorized, challenged, and used across levels and
institutions.

\subsubsection{Consequence}\label{consequence-5}

Rationality includes management of evidence standards, uncertainty,
incentives, records, observation channels, omitted populations, feedback
provenance, option generation, authority, appeal, and correction---not
only local inference.

\begin{center}\rule{0.5\linewidth}{0.5pt}\end{center}

\subsection{CU-Ep6 --- Meta-Epistemic
Awareness}\label{cu-ep6-meta-epistemic-awareness}

\textbf{Authority class:} FRAMEWORK-CANONICAL

\subsubsection{Statement}\label{statement-40}

Mature inquiry represents the limits and distortion pressures of its own
methods, incentives, identities, scales, option sets, access conditions,
observation channels, and institutional dependencies.

\subsubsection{Limit}\label{limit-5}

Awareness of bias does not confer immunity from bias. Functionally
separated modules do not become institutionally independent merely
because they occupy different boxes. Humility does not require paralysis
or false equivalence.

\begin{center}\rule{0.5\linewidth}{0.5pt}\end{center}

\section{Part IV --- Dynamics of
Coherence}\label{part-iv-dynamics-of-coherence}

\subsection{Introduction to the dynamics
family}\label{introduction-to-the-dynamics-family}

The v0.09 dynamics family was written as though one universal flow
equation governed every domain. Physics Appendix A.6 establishes that
this is not currently defensible. Version 0.10 retains the family as a
set of domain-general schemas and design principles. A domain acquires
an actual dynamics only after its state space, constraints, geometry,
functional, interventions, and time parameter have been specified.

\subsection{CU-D1 --- Domain-Specified Constrained
Dynamics}\label{cu-d1-domain-specified-constrained-dynamics}

\textbf{Registry status:} CANONICAL as a schema\\
\textbf{Claim status:} {[}Definition{]} and {[}Bridge principle{]}

\subsubsection{Statement}\label{statement-41}

A domain dynamics specifies how states or trajectories change under
declared constraints, couplings, interventions, and environmental
conditions. Its general form may be written

\[\adjustbox{max width=\linewidth}{$\displaystyle 
\mathcal F_{\mathcal D}[x,\dot x,\ddot x,\ldots;K,G,E,u,t]=0,
$}\]

or as a transition kernel, update relation, control system, or path
measure appropriate to the domain.

There is no currently established universal equation of the form

\[\adjustbox{max width=\linewidth}{$\displaystyle 
\frac{dC}{dt}
=
\alpha\nabla C\cdot v
-
\beta\frac{dS}{dt}
+
\gamma\Lambda
$}\]

governing every physical, biological, psychological, social, and
artificial system.

\subsubsection{Scope}\label{scope-16}

The inherited Universal Flow Equation is retained only as a heuristic
decomposition of possible contributions after the relevant quantities
are defined. It is not a framework-level law.

\subsubsection{Operational
consequence}\label{operational-consequence-23}

Any use of a flow equation must specify dimensions, metric, functional,
state space, forcing, boundary conditions, and empirical interpretation.

\subsubsection{Failure modes}\label{failure-modes-15}

\begin{itemize}
\tightlist
\item
  treating a normalized sandbox equation as a dimensionful physical law;
\item
  inserting undefined coherence and entropy terms;
\item
  assuming all dynamics is first-order gradient ascent;
\item
  inferring universal monotonicity.
\end{itemize}

\begin{center}\rule{0.5\linewidth}{0.5pt}\end{center}

\subsection{CU-D2 --- Regulated
Coherence-Seeking}\label{cu-d2-regulated-coherence-seeking}

\textbf{Registry status:} CANONICAL--CONDITIONAL\\
\textbf{Claim status:} {[}Bridge principle{]}

\subsubsection{Statement}\label{statement-42}

Systems with feedback, error sensitivity, and persistence-bearing
constraints may exhibit regulated tendencies that restore or expand
domain-specific coherence after perturbation.

\subsubsection{Scope}\label{scope-17}

This principle does not posit a universal cosmic drive. A rock, a cell,
a person, and an institution do not ``seek coherence'' in the same
sense. In physical systems, movement may be described by ordinary
dynamics without intrinsic regulation. In living and cognitive systems,
feedback can make preservation of organization causally operative.

\subsubsection{Operational
consequence}\label{operational-consequence-24}

A claim of coherence drive must identify the detector, error signal,
control pathway, energy source, target or viable set, and conditions
under which the regulation fails.

\subsubsection{Failure modes}\label{failure-modes-16}

\begin{itemize}
\tightlist
\item
  anthropomorphizing passive dynamics;
\item
  treating persistence after selection as goal-directed intention;
\item
  equating short-term tension reduction with global repair;
\item
  using ``the universe seeks coherence'' as a scientific conclusion.
\end{itemize}

\begin{center}\rule{0.5\linewidth}{0.5pt}\end{center}

\subsection{CU-D3 --- Gateway to Active Biological
Self-Maintenance}\label{cu-d3-gateway-to-active-biological-self-maintenance}

\textbf{Authority class:} DOMAIN GATEWAY CROSS-REFERENCE\\
\textbf{Controlling domain principle:} CU-B1

\subsubsection{Statement}\label{statement-43}

The transition from passive physical persistence to active biological
self-maintenance is owned by CU Biology. CU-D3 does not state an
independent universal dynamics law.

\subsubsection{Operational
consequence}\label{operational-consequence-25}

Applications must use CU-B1 and the Biology domain model to specify the
candidate unit, boundary, environment, perturbation class, endogenous
maintenance processes, viability conditions, and horizon.
Organism-relative stake must not be promoted directly into phenomenal
valence or moral status.

\begin{center}\rule{0.5\linewidth}{0.5pt}\end{center}

\subsection{CU-D4 --- Coherence Gradients Are
Model-Relative}\label{cu-d4-coherence-gradients-are-model-relative}

\textbf{Registry status:} CANONICAL

\subsubsection{Statement}\label{statement-44}

A coherence gradient is defined only when a declared scalar
\(C_{\mathcal D}\) of a type licensed under CU-FP2 and a geometry or
inner product on the relevant state space have been specified:

\[\adjustbox{max width=\linewidth}{$\displaystyle 
\operatorname{grad}_{g} C_{\mathcal D}
\quad\text{such that}\quad
 g(\operatorname{grad}_{g}C_{\mathcal D},v)=dC_{\mathcal D}[v].
$}\]

\subsubsection{Scope}\label{scope-18}

Euclidean and Fisher-natural gradients generally differ. A gradient
identifies local ascent of the chosen functional, not moral improvement
or physical inevitability.

\subsubsection{Operational
consequence}\label{operational-consequence-26}

Every gradient claim must name \(C_{\mathcal D}\), \(g\), the domain,
the scalar license, and the admissible tangent directions.

\subsubsection{Failure modes}\label{failure-modes-17}

\begin{itemize}
\tightlist
\item
  writing \(\nabla C\) without a metric or functional;
\item
  promoting a toy proxy to substantive coherence;
\item
  interpreting local ascent as long-horizon viability;
\item
  assuming dynamics follows the gradient.
\end{itemize}

\begin{center}\rule{0.5\linewidth}{0.5pt}\end{center}

\subsection{CU-D5 --- Non-Equilibrium
Maintenance}\label{cu-d5-non-equilibrium-maintenance}

\textbf{Registry status:} CANONICAL--CONDITIONAL\\
\textbf{Claim status:} {[}Imported physical principle{]} and {[}Bridge
principle{]}

\subsubsection{Statement}\label{statement-45}

Many organized systems persist away from equilibrium only through
ongoing exchange of energy, matter, information, or work with an
environment. Their maintenance must be evaluated in the enlarged system
that includes relevant costs and fluxes.

\subsubsection{Scope}\label{scope-19}

This principle is physically grounded for open thermodynamic systems.
Psychological, institutional, and cultural uses are analogical or
model-specific unless a measurable flow and mechanism are supplied.

\subsubsection{Operational
consequence}\label{operational-consequence-27}

Claims of local order must disclose sustaining inputs and exported
costs.

\subsubsection{Failure modes}\label{failure-modes-18}

\begin{itemize}
\tightlist
\item
  ``defying entropy'' language;
\item
  treating externalized costs as eliminated;
\item
  assuming every stable organization is far from equilibrium in the same
  technical sense;
\item
  converting thermodynamic necessity into moral doctrine.
\end{itemize}

\begin{center}\rule{0.5\linewidth}{0.5pt}\end{center}

\subsection{CU-D6 --- Multi-Timescale
Organization}\label{cu-d6-multi-timescale-organization}

\textbf{Registry status:} CANONICAL as a schema

\subsubsection{Statement}\label{statement-46}

Coherent trajectories may require coordination among processes operating
at different characteristic timescales. Fast adaptation can destabilize
slower identity or governance constraints; slow adaptation can fail in a
rapidly changing environment.

\subsubsection{Scope}\label{scope-20}

No universal timescale hierarchy is assumed. Relevant clocks and
coupling strengths are domain-specific.

\subsubsection{Operational
consequence}\label{operational-consequence-28}

Models should identify which variables change quickly, which constraints
change slowly, and when time-scale separation fails.

\subsubsection{Failure modes}\label{failure-modes-19}

\begin{itemize}
\tightlist
\item
  optimizing rapid response at the cost of long-run integration;
\item
  treating slow institutions as necessarily wise or fast systems as
  necessarily unstable;
\item
  ignoring feedback delays and hysteresis.
\end{itemize}

\begin{center}\rule{0.5\linewidth}{0.5pt}\end{center}

\subsection{CU-D7 --- Novelty Through Constrained
Exploration}\label{cu-d7-novelty-through-constrained-exploration}

\textbf{Registry status:} CANONICAL as a research schema\\
\textbf{Claim status:} {[}Bridge principle{]} and domain-specific
{[}Hypothesis{]}

\subsubsection{Statement}\label{statement-47}

Novel organization can arise when a system explores variable
configurations while constraints preserve, select, or amplify some
continuations. Exploration and constraint jointly shape innovation.

\subsubsection{Scope}\label{scope-21}

This does not replace mutation, recombination, learning algorithms,
creativity, or cultural transmission with one mechanism. It identifies a
recurring abstract structure whose domain realizations must be
specified.

\subsubsection{Operational
consequence}\label{operational-consequence-29}

A novelty claim should identify the source of variation, constraint or
selection process, retention mechanism, and evaluation horizon.

\subsubsection{Failure modes}\label{failure-modes-20}

\begin{itemize}
\tightlist
\item
  implying all novelty is directed toward higher coherence;
\item
  romanticizing randomness;
\item
  treating post hoc selection as foresight;
\item
  using the schema where no retention mechanism exists.
\end{itemize}

\begin{center}\rule{0.5\linewidth}{0.5pt}\end{center}

\subsection{CU-D8 --- Repair Before Unbounded
Optimization}\label{cu-d8-repair-before-unbounded-optimization}

\textbf{Registry status:} CANONICAL as a design and normative bridge
principle

\subsubsection{Statement}\label{statement-48}

Where a system has protected viability, agency, or integrity conditions,
maintaining repair capacity and restoring breached floors take priority
over unbounded optimization of a secondary objective.

\subsubsection{Scope}\label{scope-22}

This is not a universal physical law. It is a design and governance
principle for systems capable of damage, responsibility, or protected
function.

\subsubsection{Operational
consequence}\label{operational-consequence-30}

Architectures should include failure detection, graceful degradation,
rollback or appeal, resource reserves, and routes of correction before
pursuing maximum performance.

\subsubsection{Failure modes}\label{failure-modes-21}

\begin{itemize}
\tightlist
\item
  optimizing a proxy while damage accumulates;
\item
  punishment without reintegration where repair is possible;
\item
  treating restoration of an old state as the only form of repair;
\item
  using repair language to avoid structural change.
\end{itemize}

\begin{center}\rule{0.5\linewidth}{0.5pt}\end{center}

\subsection{CU-D9 --- Multi-Scale Failure and
Externalization}\label{cu-d9-multi-scale-failure-and-externalization}

\textbf{Registry status:} CANONICAL

\subsubsection{Statement}\label{statement-49}

A system may preserve apparent local coherence by transferring
contradiction, risk, entropy, labor, trauma, or ecological cost to
another scale, population, subsystem, or future. Such externalization
must be included in any global evaluation.

\subsubsection{Scope}\label{scope-23}

Different exported quantities require different measures. ``Externalized
incoherence'' is a family resemblance, not one conserved substance.

\subsubsection{Operational
consequence}\label{operational-consequence-31}

Failure analysis must identify who or what absorbs the cost, how it is
hidden, and whether the receiving system can repair.

\subsubsection{Failure modes}\label{failure-modes-22}

\begin{itemize}
\tightlist
\item
  assuming cost displacement is cost removal;
\item
  using one metaphor for thermodynamic, psychological, and economic
  transfer without translation;
\item
  defining the system boundary to exclude victims.
\end{itemize}

\begin{center}\rule{0.5\linewidth}{0.5pt}\end{center}

\subsection{CU-D10 --- Memory as History-Conditioned
Constraint}\label{cu-d10-memory-as-history-conditioned-constraint}

\textbf{Registry status:} CANONICAL as a schema\\
\textbf{Claim status:} {[}Definition{]} and {[}Bridge principle{]}

\subsubsection{Statement}\label{statement-50}

Memory is any history-dependent modification through which prior states
or events alter the system's present dynamics, accessible
representations, or space of future continuations.

A general representation is

\[\adjustbox{max width=\linewidth}{$\displaystyle 
K_t=\mathcal U(K_{t-1},x_{t-1},a_{t-1},E_{t-1}),
$}\]

or

\[\adjustbox{max width=\linewidth}{$\displaystyle 
\dot x_t=F(x_t;H_{0:t}),
$}\]

where the relevant history may be compressed into state variables or
remain irreducibly path-dependent for the model.

\subsubsection{Scope}\label{scope-24}

Memory is not universally identical with stored information, narrative
recall, identity, or velocity in phase space. Biological, psychological,
institutional, and artificial memory have different binding relations to
identity and agency.

\subsubsection{Operational
consequence}\label{operational-consequence-32}

A memory claim must state what changes, how the past is retained,
whether it is accessible, how it affects future behavior, and whether it
binds a persistent subject or institution.

\subsubsection{Failure modes}\label{failure-modes-23}

\begin{itemize}
\tightlist
\item
  storage reductionism;
\item
  equating retrievable records with lived memory;
\item
  treating copied weights as proof of subject continuity;
\item
  assuming a second-order state captures all history;
\item
  romanticizing traumatic path dependence.
\end{itemize}

\begin{center}\rule{0.5\linewidth}{0.5pt}\end{center}

\section{Disposition of the Former Three Laws and Double
Helix}\label{disposition-of-the-former-three-laws-and-double-helix}

\subsection{Former First Law --- ``Conservation of Entropy in Coherent
Processes''}\label{former-first-law-conservation-of-entropy-in-coherent-processes}

\textbf{Registry status:} RETIRED as a CU law.

The thermodynamic statement that total entropy obeys the relevant
second-law constraints belongs to physics under stated boundary
conditions. The cross-domain claim that coherence cannot destroy entropy
and merely reorganizes it was too imprecise to function as one law. Its
defensible components now appear in CU-FP6, CU-D5, and CU-D9.

\subsection{Former Second Law --- ``Coherence Requires Degrees of
Freedom''}\label{former-second-law-coherence-requires-degrees-of-freedom}

\textbf{Registry status:} RETIRED as a universal law; retained as a
domain hypothesis.

Many adaptive systems require variability and constraint, but no
framework-level theorem establishes that coherence vanishes at zero and
maximal entropy in every domain. Goldilocks conditions are state-,
model-, and measure-dependent, and the mathematical joint-growth wedge
can collapse.

\subsection{Former Third Law --- ``Coherence Gradient
Flow''}\label{former-third-law-coherence-gradient-flow}

\textbf{Registry status:} RETIRED as a universal law.

Gradient flow is one possible dynamics after a functional and metric are
chosen. Physics Appendix A.6 treats first-order gradient flow as an
effective regime within a broader second-order construction, not as the
universal behavior of all systems.

\subsection{Former Universal Flow
Equation}\label{former-universal-flow-equation}

\textbf{Registry status:} SUPERSEDED by CU-D1.

The equation may remain as an illustrative bookkeeping ansatz in a
specified model. It may not be cited as a derived universal law.

\subsection{Coherence--Entropy Double
Helix}\label{coherenceentropy-double-helix}

\textbf{Registry status:} PROVISIONAL HYPOTHESIS.

Alternating exploration and consolidation may describe learning,
development, evolution, and some extended dynamics. It is not entailed
by the autonomous first-order sandbox, and the standard Lyapunov
structure of fixed-coefficient gradient systems rules out nonconstant
recurrent cycles (Hirsch, Smale, \& Devaney, 2013). Each domain use
requires a mechanism, variables, timescale, and empirical test.

\begin{center}\rule{0.5\linewidth}{0.5pt}\end{center}

\begin{center}\rule{0.5\linewidth}{0.5pt}\end{center}

\section{Part V --- Cross-Domain
Governance}\label{part-v-cross-domain-governance}

This family governs consequential application of CU. It is normative and
methodological, not a physical law or empirical scalar. Domain
principles can strengthen these requirements but may not silently weaken
them.

\subsection{CU-G1 --- Affected Scope, Distribution, and Protected
Conditions}\label{cu-g1-affected-scope-distribution-and-protected-conditions}

\textbf{Authority class:} FRAMEWORK-CANONICAL GOVERNANCE RULE\\
\textbf{Claim status:} {[}Governance rule{]} and bounded {[}Bridge
principle{]}

\subsubsection{Statement}\label{statement-51}

Any consequential evaluation must identify affected beings and groups,
the distribution of benefits and burdens, protected conditions or
floors, and whether gains for some are being used to compensate harms to
others. Aggregate improvement does not by itself establish
admissibility, especially where harms are severe, concentrated,
nonconsensual, or difficult to repair.

\subsubsection{Source convergence}\label{source-convergence}

This principle consolidates recurring constraints in: CU-FP4, CU-D9,
CU-E-9, CU-S5, CU-EC-2, CU-AI-12.

\subsubsection{Review gate}\label{review-gate-7}

cross-domain methods, ethics, social science, and affected-community
review.

\subsubsection{Failure modes}\label{failure-modes-24}

\begin{itemize}
\tightlist
\item
  aggregate masking; omitted affected parties; protected floors hidden
  inside weights; compensation of severe concentrated harm through
  unrelated gains.
\end{itemize}

\subsection{CU-G2 --- Option-Set and Framing
Integrity}\label{cu-g2-option-set-and-framing-integrity}

\textbf{Authority class:} FRAMEWORK-CANONICAL GOVERNANCE RULE\\
\textbf{Claim status:} {[}Governance rule{]} and bounded {[}Bridge
principle{]}

\subsubsection{Statement}\label{statement-52}

Every screening or ranking result is conditional on the problem framing
and the options presented. Consequential applications must disclose who
generated and excluded alternatives, permit appropriately situated
challenges and affected-party proposals, and consider no action, delay,
staged action, redesign, and allegedly fixed constraints where relevant.
Screening admissibility does not establish option-set adequacy.

\subsubsection{Source convergence}\label{source-convergence-1}

This principle consolidates recurring constraints in: CU-E-1, CU-E-15,
CU-AI-23.

\subsubsection{Review gate}\label{review-gate-8}

ethics, decision theory, institutional design, and affected-community
review.

\subsubsection{Failure modes}\label{failure-modes-25}

\begin{itemize}
\tightlist
\item
  ethically truncated menus; sponsor-controlled feasibility; exclusion
  of no action, delay, staging, redesign, or affected-party
  alternatives.
\end{itemize}

\subsection{CU-G3 --- Practical Correctability and Independent
Challenge}\label{cu-g3-practical-correctability-and-independent-challenge}

\textbf{Authority class:} FRAMEWORK-CANONICAL GOVERNANCE RULE\\
\textbf{Claim status:} {[}Governance rule{]} and bounded {[}Bridge
principle{]}

\subsubsection{Statement}\label{statement-53}

Practical correctability requires accessible challenge, sufficiently
independent review, authority to reverse or cease an action, correction
of records, remedy, recurrence tracking, revision of rules, and
learning. Nominal complaint channels or sole self-certification by the
institution or system under review are insufficient.

\subsubsection{Source convergence}\label{source-convergence-2}

This principle consolidates recurring constraints in: CU-FP7, CU-D8,
CU-Ψ7, CU-S3, CU-Inst-3, CU-E-6, CU-AI-24.

\subsubsection{Review gate}\label{review-gate-9}

institutional design, law, ethics, governance, and affected-community
review.

\subsubsection{Failure modes}\label{failure-modes-26}

\begin{itemize}
\tightlist
\item
  complaint channels without reversal authority; provider-controlled
  auditing; record correction without remedy; no recurrence measurement.
\end{itemize}

\subsection{CU-G4 --- Responsibility--Authority--Causal-Reach
Alignment}\label{cu-g4-responsibilityauthoritycausal-reach-alignment}

\textbf{Authority class:} FRAMEWORK-CANONICAL GOVERNANCE RULE\\
\textbf{Claim status:} {[}Governance rule{]} and bounded {[}Bridge
principle{]}

\subsubsection{Statement}\label{statement-54}

Responsibility assessment must distinguish causal contribution,
knowledge, foreseeability, control, authority, benefit, coercion,
prevention capacity, and repair capacity. Responsibility without
adequate authority and authority without accountability are structural
failures. Governance burden increases with causal reach and
irreversibility.

\subsubsection{Source convergence}\label{source-convergence-3}

This principle consolidates recurring constraints in: CU-ICF-2,
CU-ICF-3, CU-Inst-6, CU-E-14, CU-AI-6, CU-AI-15.

\subsubsection{Review gate}\label{review-gate-10}

ethics, law, institutional design, AI governance, and affected-community
review.

\subsubsection{Failure modes}\label{failure-modes-27}

\begin{itemize}
\tightlist
\item
  blaming low-authority actors; authority without answerability;
  ignoring beneficiaries, prevention capacity, or repair capacity.
\end{itemize}

\subsection{CU-G5 --- Participation Proportional to Governance
Reach}\label{cu-g5-participation-proportional-to-governance-reach}

\textbf{Authority class:} FRAMEWORK-CANONICAL GOVERNANCE RULE\\
\textbf{Claim status:} {[}Governance rule{]} and bounded {[}Bridge
principle{]}

\subsubsection{Statement}\label{statement-55}

Principles and systems that classify, burden, represent, or govern
people require participation by affected communities proportionate to
their reach and risk, with substantive authority where stakes are high.
Participation supplements rather than replaces evidence and argument.
Required review gates are conjunctive within their competence;
unresolved conflict preserves provisional status and public dissent.

\subsubsection{Source convergence}\label{source-convergence-4}

This principle consolidates recurring constraints in: CU-E-15, CU-S5,
CU-Inst-0, CU-AI-14, CU-AI-23.

\subsubsection{Review gate}\label{review-gate-11}

heightened affected-community co-governance and public-review design.

\subsubsection{Failure modes}\label{failure-modes-28}

\begin{itemize}
\tightlist
\item
  symbolic consultation; treating participation as a substitute for
  evidence; overriding one review competence with another.
\end{itemize}

\subsection{CU-G6 --- Patiency, Agency, Responsibility, and Sovereignty
Are
Distinct}\label{cu-g6-patiency-agency-responsibility-and-sovereignty-are-distinct}

\textbf{Authority class:} FRAMEWORK-CANONICAL GOVERNANCE RULE\\
\textbf{Claim status:} {[}Governance rule{]} and bounded {[}Bridge
principle{]}

\subsubsection{Statement}\label{statement-56}

Evidence of experience or welfare may ground moral patiency; capacities
for reason-guided action may support agency; responsibility additionally
depends on knowledge, control, continuity, and answerability;
sovereignty requires legitimate social authorization. None of these
statuses automatically entails the others.

\subsubsection{Source convergence}\label{source-convergence-5}

This principle consolidates recurring constraints in: CU-V7, CU-E-2,
CU-E-17, CU-AI-8, CU-AI-22.

\subsubsection{Review gate}\label{review-gate-12}

consciousness science, ethics, law, AI welfare, and public governance
review.

\subsubsection{Failure modes}\label{failure-modes-29}

\begin{itemize}
\tightlist
\item
  inferring sovereignty from intelligence or consciousness; assigning
  responsibility without continuity, knowledge, or control; ignoring
  powerful non-conscious systems.
\end{itemize}

\subsection{CU-G7 --- Reversibility, Containment, and Retained Option
Value}\label{cu-g7-reversibility-containment-and-retained-option-value}

\textbf{Authority class:} FRAMEWORK-CANONICAL GOVERNANCE RULE\\
\textbf{Claim status:} {[}Governance rule{]} and bounded {[}Bridge
principle{]}

\subsubsection{Statement}\label{statement-57}

Under credible risk of severe or irreversible harm, an application must
preserve context-appropriate reversibility, containment, rollback,
staged exposure, or equivalent retained option value, or publicly
justify a necessity-and-proportionality exception that includes the
harms of delay. Reversibility is a context-sensitive safeguard, not an
absolute ban on irreversible action.

\subsubsection{Source convergence}\label{source-convergence-6}

This principle consolidates recurring constraints in: CU-E-16, CU-AI-16.

\subsubsection{Review gate}\label{review-gate-13}

ethics, safety engineering, law, domain expertise, and
affected-community review.

\subsubsection{Failure modes}\label{failure-modes-30}

\begin{itemize}
\tightlist
\item
  trading severe irreversibility for aggregate performance; using
  reversibility as an absolute veto without proportionality; ignoring
  the harms of delay.
\end{itemize}

\begin{center}\rule{0.5\linewidth}{0.5pt}\end{center}

\section{Part VI --- Physics Constraint
Registry}\label{part-vi-physics-constraint-registry}

This part does not create a new CU-P numbering family. It records the
physical conclusions every downstream paper must respect.

\subsection{PHYS-A6.1 --- Statistical differentiation is a toy
proxy}\label{phys-a6.1-statistical-differentiation-is-a-toy-proxy}

\textbf{Status:} CANONICAL.

\[\adjustbox{max width=\linewidth}{$\displaystyle 
D(p)=\sum_i(p_i-u_i)^2
$}\]

measures differentiation from a chosen uniform reference in the
classical sandbox. It is not substantive physical coherence and must not
be exported to other domains by renaming it.

\subsection{PHYS-A6.2 --- Euclidean and Fisher geometry
differ}\label{phys-a6.2-euclidean-and-fisher-geometry-differ}

\textbf{Status:} CANONICAL.

Gradients, norms, potentiometers, and ascent conditions depend on the
chosen geometry. Results must state which is used.

\subsection{PHYS-A6.3 --- Goldilocks conditions are
plural}\label{phys-a6.3-goldilocks-conditions-are-plural}

\textbf{Status:} CANONICAL.

Differentiation ascent, joint entropy--differentiation growth,
finite-step feasibility, and developmental exploration/consolidation are
distinct meanings. The joint-growth wedge is state-dependent and can
collapse on two-level distributions.

\subsection{PHYS-A6.4 --- Temporal order is not temporal
direction}\label{phys-a6.4-temporal-order-is-not-temporal-direction}

\textbf{Status:} CANONICAL.

Ordered histories and clock parameterization do not by themselves yield
a thermodynamic arrow. Direction requires asymmetric dynamics, boundary
conditions, reduced descriptions, or irrecoverability assumptions.

\subsection{PHYS-A6.5 --- Classical hidden correlation does not force
quantum
theory}\label{phys-a6.5-classical-hidden-correlation-does-not-force-quantum-theory}

\textbf{Status:} CANONICAL.

Classical joint distributions can contain correlations hidden from
marginals. A quantum completion requires stronger operational and
contextual assumptions.

\subsection{PHYS-A6.6 --- Quantum selection is
conditional}\label{phys-a6.6-quantum-selection-is-conditional}

\textbf{Status:} CANONICAL--CONDITIONAL.

Purification, reversible dilation, contextual accessibility, and the
operational axioms of an appropriate reconstruction program can
conditionally select finite-dimensional complex quantum theory. CU does
not currently derive those operational axioms from coherence alone.

\subsection{PHYS-A6.7 --- Quantum quantities are
distinct}\label{phys-a6.7-quantum-quantities-are-distinct}

\textbf{Status:} CANONICAL.

Purity, basis coherence, total correlation, entanglement, and von
Neumann entropy must not be identified.

\subsection{PHYS-A6.8 --- Second-order dynamics is proposed, not
established as
fundamental}\label{phys-a6.8-second-order-dynamics-is-proposed-not-established-as-fundamental}

\textbf{Status:} CANONICAL classification.

The Fisher--Rao Lagrangian/Hamiltonian construction is mathematically
coherent. Its status as fundamental physics is a hypothesis. First-order
flow appears as a strongly damped limit under stated assumptions.

\subsection{PHYS-A6.9 --- Standard equations are embeddings or
conditional
recoveries}\label{phys-a6.9-standard-equations-are-embeddings-or-conditional-recoveries}

\textbf{Status:} CANONICAL classification.

Wave, Madelung/Schrödinger, relational-geometry, and Jacobson-style
gravity results arise after specific actions, potentials, and bridge
assumptions are supplied (Jacobson, 1995). They are not unique
derivations of nature's laws from one coherence principle.

\subsection{PHYS-A6.10 --- Accessibility geometry remains
incomplete}\label{phys-a6.10-accessibility-geometry-remains-incomplete}

\textbf{Status:} CANONICAL open problem.

Metric closure can construct effective distances from positive
accessibility weights. Lorentzian signature, dimension, locality,
continuum dynamics, and empirical distinctiveness remain to be derived
or tested.

\subsection{PHYS-A6.11 --- Horizon entropy is reduced-description
entropy}\label{phys-a6.11-horizon-entropy-is-reduced-description-entropy}

\textbf{Status:} CANONICAL classification.

Horizon entropy can be interpreted as entropy associated with a causally
restricted description under additional physical assumptions. It is not
an amount of hidden substantive coherence.

\subsection{PHYS-A6.12 --- Fundamental-substrate claims are
interpretive}\label{phys-a6.12-fundamental-substrate-claims-are-interpretive}

\textbf{Status:} CANONICAL classification.

The coherence field as the structure of physical possibility is an
interpretive and research proposal. It is not established by the
mathematical recoveries alone.

\begin{center}\rule{0.5\linewidth}{0.5pt}\end{center}

\section{Part VII --- Domain Gateways and Owned
Registries}\label{part-vii-domain-gateways-and-owned-registries}

A \textbf{domain authority} is the controlling CU formulation inside
that domain. It does not become a framework law, empirical consensus,
clinical standard, legal rule, or public mandate merely by appearing
here. Full argument, evidence, failure conditions, governance rules, and
audit boundaries remain in the source paper.

\subsection{Biology}\label{biology}

Domain source: \emph{Coherence Universalism: Biology}, Version 0.10.
These entries govern CU's biological vocabulary but remain subject to
systems-biology, origin-of-life, evolutionary, and philosophy-of-biology
review.

\subsubsection{CU-B1 --- Active Biological
Self-Maintenance}\label{cu-b1-active-biological-self-maintenance}

\textbf{Authority class:} DOMAIN-AUTHORITY BRIDGE PRINCIPLE\\
\textbf{Domain owner:} Biology\\
\textbf{Claim/candidate status:} CANONICAL BRIDGE PRINCIPLE\\
\textbf{Review gate:} systems biology, philosophy of biology, and
relevant empirical review

A candidate living organization is not distinguished by static order or
persistence alone, but by a relatively closed network of materially
realized processes and constraints whose endogenous activity contributes
to maintaining, repairing, reproducing, or adaptively transforming the
conditions of that network's continuation. Degree and kind of autonomy
are relative to a boundary, environment, perturbation class, and
horizon.

\subsubsection{CU-B2 --- Adaptive Problem Solving Beyond Fixed
Programs}\label{cu-b2-adaptive-problem-solving-beyond-fixed-programs}

\textbf{Authority class:} DOMAIN-AUTHORITY CONDITIONAL EMPIRICAL CLAIM\\
\textbf{Domain owner:} Biology\\
\textbf{Claim/candidate status:} CANONICAL--CONDITIONAL EMPIRICAL
HYPOTHESIS\\
\textbf{Review gate:} systems biology, philosophy of biology, and
relevant empirical review

Some living systems possess regulatory architectures that generalize
beyond fixed stimulus--response mappings: they detect deviations,
recruit alternative pathways, modify future responses from history, and
solve novel viability-relevant problems within a bounded adaptive reach.
Degree must be established through counterfactual flexibility, learning
controls, and comparison with passive or preprogrammed alternatives; it
does not by itself establish consciousness.

\subsubsection{CU-B3 --- Multiscale Constraint and Local
Process}\label{cu-b3-multiscale-constraint-and-local-process}

\textbf{Authority class:} DOMAIN-AUTHORITY BRIDGE PRINCIPLE\\
\textbf{Domain owner:} Biology\\
\textbf{Claim/candidate status:} CANONICAL BRIDGE PRINCIPLE\\
\textbf{Review gate:} systems biology, philosophy of biology, and
relevant empirical review

Biological form and function arise through local physical processes
embedded in cross-scale boundary conditions, signaling networks,
mechanical relations, and histories. A higher-scale constraint is
explanatory when intervention on its realizing relations changes
lower-scale trajectories in a way that improves prediction or control
beyond local variables alone. Higher scale does not imply a new force,
and global stability does not imply a unique optimizer.

\subsubsection{CU-B4 --- Structured Evolutionary Possibility and the
Stabilization of Regulatory
Architecture}\label{cu-b4-structured-evolutionary-possibility-and-the-stabilization-of-regulatory-architecture}

\textbf{Authority class:} DOMAIN-AUTHORITY CONDITIONAL EMPIRICAL CLAIM\\
\textbf{Domain owner:} Biology\\
\textbf{Claim/candidate status:} CANONICAL--CONDITIONAL EMPIRICAL
GENERALIZATION\\
\textbf{Review gate:} evolutionary biology, evo-devo,
genotype--phenotype-map research, philosophy of biology, and relevant
empirical review

Evolutionary processes explore a structured and nonuniform possibility
space shaped by mutation bias, genotype--phenotype relations,
developmental dynamics, physical feasibility, ecological interaction,
inheritance, and historical path dependence. Selection nonrandomly
changes the frequencies of heritable variants relative to local
reproductive consequences. Recurrent constraints can make some forms,
functions, and regulatory architectures repeatedly accessible,
contributing to convergence, while chance, drift, conflict, changing
environments, and historical contingency preserve multiple possible
trajectories. Evolution can stabilize, modify, simplify, or dismantle
coherent self-maintaining organization; it does not thereby maximize
coherence, complexity, consciousness, or moral value.

\textbf{Status boundary.} Canonical inclusion makes this the complete CU
v0.10 proposal to be reviewed. It does not represent expert validation,
a universal evolutionary law, or evidence of foresight, predetermined
design, or monotonic progress.

\subsubsection{CU-B5 --- Biological Individuation and
Stake}\label{cu-b5-biological-individuation-and-stake}

\textbf{Authority class:} DOMAIN-AUTHORITY BRIDGE PRINCIPLE\\
\textbf{Domain owner:} Biology\\
\textbf{Claim/candidate status:} CANONICAL BRIDGE PRINCIPLE\\
\textbf{Review gate:} systems biology, philosophy of biology, and
relevant empirical review\\
\textbf{Cross-domain core:} CU-V1

A biological individual is a historically continuous, relatively bounded
organization whose components are causally integrated into maintaining,
repairing, reproducing, or adaptively transforming the conditions of
that organization. The organization thereby establishes
organism-relative stakes: conditions can be better or worse for it in
virtue of their effects on its continued self-maintenance. Individuality
is graded, multidimensional, and potentially nested; organism-relative
stake does not by itself establish consciousness or moral status.

\subsection{Consciousness and Valence}\label{consciousness-and-valence}

Domain source: \emph{Coherence Universalism: Consciousness}, Version
0.10. Thin perspective, valence, unity, and artificial consciousness
remain theory-dependent and partially unresolved.

\subsubsection{CU-C1 --- Global
Availability}\label{cu-c1-global-availability}

\textbf{Authority class:} DOMAIN-AUTHORITY CONDITIONAL CONSTRUCT /
METHOD\\
\textbf{Domain owner:} Consciousness\\
\textbf{Claim/candidate status:} CANONICAL--CONDITIONAL EMPIRICAL
CONSTRUCT\\
\textbf{Review gate:} consciousness science, neuroscience, AI
interpretability, and welfare review

Selected information becomes causally available across otherwise
specialized processes for report, deliberation, flexible control, or
self-monitoring. Global availability is evidence relevant to conscious
access, not sufficient proof of phenomenology.

\subsubsection{CU-C2 --- Perspectival
Individuation}\label{cu-c2-perspectival-individuation}

\textbf{Authority class:} DOMAIN-PROVISIONAL BRIDGE HYPOTHESIS\\
\textbf{Domain owner:} Consciousness\\
\textbf{Claim/candidate status:} CANONICAL BRIDGE HYPOTHESIS\\
\textbf{Review gate:} consciousness science, neuroscience, AI
interpretability, and welfare review

Candidate status for thin perspective increases when independently
supported access organization, perspectival anchoring, and center
individuation converge. Evidence for one clause may not be counted as
evidence for the others. Singularity may be graded, temporary, or
plural.

\subsubsection{CU-C3 --- Causally Active
Self-Modeling}\label{cu-c3-causally-active-self-modeling}

\textbf{Authority class:} DOMAIN-AUTHORITY CONDITIONAL CONSTRUCT /
METHOD\\
\textbf{Domain owner:} Consciousness\\
\textbf{Claim/candidate status:} CANONICAL--CONDITIONAL CONSTRUCT\\
\textbf{Review gate:} consciousness science, neuroscience, AI
interpretability, and welfare review

Self-modeling occurs when a system represents aspects of its own
processing, role, boundary, uncertainty, history, or action and those
representations causally influence ongoing processing. Self-modeling is
distinct from persistent self-maintenance and ego.

\subsubsection{CU-C4 --- Causal
Integration}\label{cu-c4-causal-integration}

\textbf{Authority class:} DOMAIN-AUTHORITY CONDITIONAL CONSTRUCT /
METHOD\\
\textbf{Domain owner:} Consciousness\\
\textbf{Claim/candidate status:} CANONICAL METHODOLOGICAL PRINCIPLE\\
\textbf{Review gate:} consciousness science, neuroscience, AI
interpretability, and welfare review

Information associated with a candidate perspective must participate
causally in reasoning, attention, learning, report, or control rather
than exist only as an externally readable correlate. Intervention
evidence is preferred where possible.

\subsubsection{CU-C5 --- Temporal Binding}\label{cu-c5-temporal-binding}

\textbf{Authority class:} DOMAIN-PROVISIONAL BRIDGE HYPOTHESIS\\
\textbf{Domain owner:} Consciousness\\
\textbf{Claim/candidate status:} CANONICAL--CONDITIONAL BRIDGE
HYPOTHESIS\\
\textbf{Review gate:} consciousness science, neuroscience, AI
interpretability, and welfare review

The depth and continuity of a perspective depend partly on how contents
are integrated across a present interval and linked across episodes.
Temporal binding may distinguish a thin episode from a thick subject; it
is not established as necessary for every minimal perspective.

\subsubsection{CU-C6 --- World-Coupling and
Embodiment}\label{cu-c6-world-coupling-and-embodiment}

\textbf{Authority class:} DOMAIN-PROVISIONAL BRIDGE HYPOTHESIS\\
\textbf{Domain owner:} Consciousness\\
\textbf{Claim/candidate status:} CANONICAL--CONDITIONAL EMPIRICAL
GENERALIZATION\\
\textbf{Review gate:} consciousness science, neuroscience, AI
interpretability, and welfare review

Ongoing sensorimotor, interoceptive, and environmental coupling shapes
the content, stability, self-location, and stakes of biological
perspectives. Biological embodiment is a major route, not an exclusive
substrate requirement.

\subsubsection{CU-C7 --- Persistence, Self-Maintenance, and
Ego}\label{cu-c7-persistence-self-maintenance-and-ego}

\textbf{Authority class:} DOMAIN-AUTHORITY CONDITIONAL CONSTRUCT /
METHOD\\
\textbf{Domain owner:} Consciousness\\
\textbf{Claim/candidate status:} CANONICAL BRIDGE PRINCIPLE\\
\textbf{Review gate:} consciousness science, neuroscience, AI
interpretability, and welfare review\\
\textbf{Cross-domain core:} CU-FP8

Persistent identity and self-maintenance can transform a momentary
perspective into a diachronic subject with biography, anticipatory
stakes, and defensive organization. Egoic consciousness is one possible
regime, not a synonym for consciousness.

\subsubsection{CU-V1 --- Functional Stake Is Not Phenomenal
Valence}\label{cu-v1-functional-stake-is-not-phenomenal-valence}

\textbf{Authority class:} FRAMEWORK-CANONICAL BOUNDARY PRINCIPLE\\
\textbf{Domain owner:} Consciousness\\
\textbf{Claim/candidate status:} CANONICAL BOUNDARY PRINCIPLE\\
\textbf{Review gate:} consciousness science, neuroscience, AI
interpretability, and welfare review

Organism-relative or engineered better--worse organization does not by
itself establish felt pleasure, pain, or concern. Functional stake and
phenomenal valence require separate evidence.

\subsubsection{CU-V2 --- Endogenous Better--Worse
Organization}\label{cu-v2-endogenous-betterworse-organization}

\textbf{Authority class:} DOMAIN-PROVISIONAL VALENCE / SUBJECTIVITY
HYPOTHESIS\\
\textbf{Domain owner:} Consciousness\\
\textbf{Claim/candidate status:} CANONICAL--CONDITIONAL BRIDGE
HYPOTHESIS\\
\textbf{Review gate:} consciousness science, neuroscience, AI
interpretability, and welfare review

Evidence for valence strengthens when better--worse variables are
endogenous to a candidate subject, globally integrated, and causally
organize attention, learning, and action beyond externally assigned
reward.

\subsubsection{CU-V3 --- Integrated
Affect}\label{cu-v3-integrated-affect}

\textbf{Authority class:} DOMAIN-PROVISIONAL VALENCE / SUBJECTIVITY
HYPOTHESIS\\
\textbf{Domain owner:} Consciousness\\
\textbf{Claim/candidate status:} CANONICAL--CONDITIONAL EMPIRICAL
CONSTRUCT\\
\textbf{Review gate:} consciousness science, neuroscience, AI
interpretability, and welfare review

Affect-like signals become more relevant to phenomenal valence when they
influence a candidate perspective globally rather than remaining
isolated control variables or verbal labels.

\subsubsection{CU-V4 --- Irreversibility for a Continuing
Subject}\label{cu-v4-irreversibility-for-a-continuing-subject}

\textbf{Authority class:} DOMAIN-PROVISIONAL VALENCE / SUBJECTIVITY
HYPOTHESIS\\
\textbf{Domain owner:} Consciousness\\
\textbf{Claim/candidate status:} CANONICAL--CONDITIONAL BRIDGE
HYPOTHESIS\\
\textbf{Review gate:} consciousness science, neuroscience, AI
interpretability, and welfare review

The depth and moral weight of valence may increase when events
irreversibly alter the future of a continuing candidate subject.
Irreversibility is not established as necessary for every experience.

\subsubsection{CU-V5 --- Temporal
Thickness}\label{cu-v5-temporal-thickness}

\textbf{Authority class:} DOMAIN-PROVISIONAL VALENCE / SUBJECTIVITY
HYPOTHESIS\\
\textbf{Domain owner:} Consciousness\\
\textbf{Claim/candidate status:} CANONICAL--CONDITIONAL BRIDGE
HYPOTHESIS\\
\textbf{Review gate:} consciousness science, neuroscience, AI
interpretability, and welfare review

Valence and subjectivity may acquire richer character when integrated
across a bound present, memory, anticipation, and sustained
world-coupling.

\subsubsection{CU-V6 --- Temporal Organization, Not Equation
Order}\label{cu-v6-temporal-organization-not-equation-order}

\textbf{Authority class:} DOMAIN-AUTHORITY CORRECTION / PROVISIONAL
HYPOTHESIS\\
\textbf{Domain owner:} Consciousness\\
\textbf{Claim/candidate status:} CANONICAL CORRECTION / PROVISIONAL
HYPOTHESIS\\
\textbf{Review gate:} consciousness science, neuroscience, AI
interpretability, and welfare review

Recurrence, oscillation, metastability, history dependence, or
functionally equivalent dynamics may support temporal thickness. Literal
second-order differential form is neither invariant nor an exclusion
theorem for language models.

\subsubsection{CU-V7 --- Moral Patiency Is Not
Sovereignty}\label{cu-v7-moral-patiency-is-not-sovereignty}

\textbf{Authority class:} SUPERSEDED ALIAS\\
\textbf{Domain owner:} Consciousness\\
\textbf{Claim/candidate status:} --- (superseded alias; see CU-G6)\\
\textbf{Review gate:} consciousness science, neuroscience, AI
interpretability, and welfare review\\
\textbf{Cross-domain core:} CU-G6

A system may warrant moral consideration because it has experience or
valenced stake while lacking the persistent identity, accountable
agency, legitimacy, and social standing required for sovereign
authority.

\textbf{Disposition:} Historical or domain alias; controlling identifier
is CU-G6.

\subsubsection{CU-V8 --- Valence as Directional
Evaluation}\label{cu-v8-valence-as-directional-evaluation}

\textbf{Authority class:} DOMAIN-PROVISIONAL VALENCE / SUBJECTIVITY
HYPOTHESIS\\
\textbf{Domain owner:} Consciousness\\
\textbf{Claim/candidate status:} CANONICAL--CONDITIONAL BRIDGE
HYPOTHESIS\\
\textbf{Review gate:} consciousness science, affective science, AI
interpretability, and welfare review

Where phenomenal valence is present, its phasic component is
hypothesized to track the direction and rate of change of the subject's
endogenous, globally integrated better--worse organization along its
realized trajectory; its tonic component is hypothesized to track
standing relative to protected floors or margins. Directional evaluation
is plural and may be improving, worsening, or mixed across coordinates
and horizons. Anticipated directional evaluation can guide navigation.
No universal coherence scalar, privileged metric, or identity between
reward signals and felt valence is assumed, and no progress signal
establishes valence without independent evidence of a candidate
perspective, endogeneity, integration, and causal role.

\textbf{Principal failure mode:} Treating every reward-prediction or
progress signal as felt, or reintroducing a universal coherence gradient
under another name.

\subsection{Psychology}\label{psychology}

Domain source: \emph{Coherence Universalism: Psychology}, Version 0.10.
Clinical use requires field review, and principles governing trauma,
dissociation, neurodiversity, disclosure, or fragmentation require
substantive affected-community participation.

\subsubsection{CU-Ψ1 --- Meaning Space}\label{cu-ux3c81-meaning-space}

\textbf{Authority class:} DOMAIN-AUTHORITY DEFINITION / ANALYTIC
CONSTRUCTION\\
\textbf{Domain owner:} Psychology\\
\textbf{Claim/candidate status:} DEFINITION / ANALYTIC CONSTRUCTION\\
\textbf{Review gate:} clinical/psychological expert review;
affected-community participation as applicable

For a declared psychological question, a \textbf{meaning space} is a
model-dependent set of distinctions along which events, actions,
relationships, memories, identities, and possibilities can matter to a
person or group:

\[\adjustbox{max width=\linewidth}{$\displaystyle 
\mathcal M_{\Psi}=\mathcal M_{\Psi}(\text{person},\text{context},\pi,T,\mathcal O).
$}\]

Its coordinates may include expected consequence, bodily salience,
social significance, autobiographical relevance, moral implication,
practical affordance, or sacred orientation. Its coordinates and
dimensionality must be declared; no universal dimensionality or fixed
axis set is assumed. A meaning space is an analyst's construct unless
evidence establishes that its distinctions are represented or enacted by
the person.

\textbf{Principal failure mode:} Treating the researcher's vocabulary as
the person's phenomenology.

\subsubsection{CU-Ψ2 --- Meaning Maps}\label{cu-ux3c82-meaning-maps}

\textbf{Authority class:} DOMAIN-AUTHORITY DEFINITION / ANALYTIC
CONSTRUCTION\\
\textbf{Domain owner:} Psychology\\
\textbf{Claim/candidate status:} ANALYTIC BRIDGE PRINCIPLE\\
\textbf{Review gate:} clinical/psychological expert review;
affected-community participation as applicable

A \textbf{meaning map} is a structured, revisable relation among cues,
interpretations, affective responses, bodily states, memories,
commitments, predicted consequences, and action tendencies. Meaning maps
may be explicit or implicit, proposition-like or procedural,
individually carried or socially scaffolded. Their psychological
importance lies in how they organize attention, expectation, action, and
revision---not in their conformity to one canonical map.

\textbf{Principal failure mode:} Assuming that internal consistency,
narrative elegance, or social stability establishes descriptive truth.

\subsubsection{CU-Ψ3 --- Psychological Coherence
Profile}\label{cu-ux3c83-psychological-coherence-profile}

\textbf{Authority class:} DOMAIN-AUTHORITY CONDITIONAL PRINCIPLE\\
\textbf{Domain owner:} Psychology\\
\textbf{Claim/candidate status:} CANONICAL-CONDITIONAL METHODOLOGICAL
PRINCIPLE\\
\textbf{Review gate:} clinical/psychological expert review;
affected-community participation as applicable\\
\textbf{Cross-domain core:} CU-FP2

Psychological organization should ordinarily be represented by a
declared profile rather than one universal scalar. A default research
profile may include descriptive adequacy, affect tolerance and
regulation, regulatory flexibility, identity and narrative integration,
relational capacity, agency, embodied regulation, meaning or purpose,
and repair capacity:

\[\adjustbox{max width=\linewidth}{$\displaystyle 
\mathbf P_{\Psi}
=
\langle
A_d,A_f,F_r,I_n,R_c,G_a,B_r,M_p,Q_r
\rangle.
$}\]

The dimensions, observations, scales, and horizons must be justified for
the question. Scalarization is permitted only when its weights,
normalization, protected floors, affected persons, and decision purpose
are declared. Non-dominating profiles may remain incomparable.

\textbf{Principal failure mode:} Smuggling a moral or clinical ideal
into an apparently neutral score.

\subsubsection{CU-Ψ4 --- Integration Load and Regulatory
Capacity}\label{cu-ux3c84-integration-load-and-regulatory-capacity}

\textbf{Authority class:} DOMAIN-AUTHORITY CONDITIONAL PRINCIPLE\\
\textbf{Domain owner:} Psychology\\
\textbf{Claim/candidate status:} MODEL FAMILY / BRIDGE HYPOTHESIS\\
\textbf{Review gate:} clinical/psychological expert review;
affected-community participation as applicable

At a given time and scale, a person confronts a demand profile \(d_t\)
and possesses a context-sensitive capacity region \(\mathcal C_t\).
\textbf{Integration overload} occurs when relevant demands cannot be
metabolized, coordinated, deferred, delegated, escaped, or safely
compartmentalized within that capacity:

\[\adjustbox{max width=\linewidth}{$\displaystyle 
d_t\notin\mathcal C_t.
$}\]

Stress and distress depend not only on demand magnitude but on
predictability, controllability, meaning, bodily condition, social
support, prior learning, structural power, and temporal horizon.
Overload is not identical with disorder, and low expressed distress is
not sufficient evidence that demands have been integrated.

\textbf{Principal failure mode:} Locating in the individual a burden
produced by an unsafe or unjust environment.

\subsubsection{CU-Ψ5 --- Psychological Regulation and
Meaning-Seeking}\label{cu-ux3c85-psychological-regulation-and-meaning-seeking}

\textbf{Authority class:} DOMAIN-AUTHORITY CONDITIONAL PRINCIPLE\\
\textbf{Domain owner:} Psychology\\
\textbf{Claim/candidate status:} EMPIRICAL BRIDGE HYPOTHESIS\\
\textbf{Review gate:} clinical/psychological expert review;
affected-community participation as applicable

Human psychological activity is organized by multiple, sometimes
conflicting regulatory tendencies, including safety, attachment,
exploration, competence, status, identity protection, truth-seeking,
care, belonging, autonomy, pleasure, pain avoidance, and sacred or
long-horizon commitment. People often seek interpretations and practices
that make action possible, but there is no established universal
psychological drive to maximize one coherence quantity.

\textbf{Principal failure mode:} Redescribing every motive after the
fact as coherence-seeking and thereby making the principle
unfalsifiable.

\subsubsection{CU-Ψ6 --- Local-Global and Short-Long
Tension}\label{cu-ux3c86-local-global-and-short-long-tension}

\textbf{Authority class:} DOMAIN-AUTHORITY CONDITIONAL PRINCIPLE\\
\textbf{Domain owner:} Psychology\\
\textbf{Claim/candidate status:} BRIDGE PRINCIPLE\\
\textbf{Review gate:} clinical/psychological expert review;
affected-community participation as applicable

A pattern may be adaptive for one subsystem, relationship, role, or
short horizon while undermining another subsystem, wider context,
affected person, or longer horizon. Psychological evaluation must
therefore specify at least:

\[\adjustbox{max width=\linewidth}{$\displaystyle 
(\text{part or whole},\ \text{context},\ \text{horizon},\ \text{protected floors}).
$}\]

Suppression, avoidance, compartmentalization, conformity, disclosure,
confrontation, and self-revision can each reverse their profile effects
as these declarations change. No local gain or immediate relief licenses
an inference to global or durable health.

\textbf{Principal failure mode:} Treating the analyst's preferred
horizon as uniquely rational while ignoring acute safety.

\subsubsection{CU-Ψ7 --- Repair
Capacity}\label{cu-ux3c87-repair-capacity}

\textbf{Authority class:} DOMAIN-AUTHORITY CONDITIONAL PRINCIPLE\\
\textbf{Domain owner:} Psychology\\
\textbf{Claim/candidate status:} BRIDGE PRINCIPLE / EMPIRICAL
HYPOTHESIS\\
\textbf{Review gate:} clinical/psychological expert review;
affected-community participation as applicable\\
\textbf{Cross-domain core:} CU-G3

\textbf{Repair capacity} is the context-sensitive ability to detect
consequential mismatch or rupture, remain or return within a tolerable
range of engagement, revise models or strategies, recruit support, make
restitution where appropriate, restore workable coordination, and learn
in ways that reduce future recurrence. Repair need not restore the prior
state. It may require grief, accommodation, boundary, identity change,
or transformation of the environment.

\textbf{Principal failure mode:} Equating repair with reconciliation,
forgiveness, symptom elimination, productivity, or return to a harmful
baseline.

\subsubsection{CU-Ψ8 --- Protective Compartmentalization and
Fragmentation}\label{cu-ux3c88-protective-compartmentalization-and-fragmentation}

\textbf{Authority class:} DOMAIN-PROVISIONAL CLINICAL / SOCIAL BRIDGE\\
\textbf{Domain owner:} Psychology\\
\textbf{Claim/candidate status:} CONDITIONAL BRIDGE PRINCIPLE\\
\textbf{Review gate:} clinical expert review and substantive
affected-community co-governance

Psychological separation among memories, affects, roles, self-states, or
action systems can protect functioning when full integration would
exceed present capacity or threaten safety. The same separation can
become costly when it blocks reality contact, agency, flexible access,
relationship, or repair after the original conditions change. The
distinction between \textbf{protection} and \textbf{fragmentation} is
therefore functional, contextual, and horizon-relative rather than
purely structural.

\textbf{Principal failure mode:} Pathologizing survival adaptations or
romanticizing costly dissociation.

\subsubsection{CU-Ψ9 --- Developmental Path Dependence and
Plasticity}\label{cu-ux3c89-developmental-path-dependence-and-plasticity}

\textbf{Authority class:} DOMAIN-AUTHORITY CONDITIONAL PRINCIPLE\\
\textbf{Domain owner:} Psychology\\
\textbf{Claim/candidate status:} EMPIRICAL BRIDGE PRINCIPLE\\
\textbf{Review gate:} clinical/psychological expert review;
affected-community participation as applicable

Psychological development is path-dependent: earlier relationships,
bodily states, learning histories, institutions, and material conditions
alter the accessibility, cost, and meaning of later transitions.
Development nevertheless remains plastic. Similar adult patterns can
arise by different paths, and different outcomes can arise from similar
beginnings. Attractor language is legitimate only when a model specifies
states, transitions, perturbations, and evidence of return or basin
structure.

\textbf{Principal failure mode:} Turning statistical regularities or
culturally preferred biographies into destiny.

\subsubsection{CU-Ψ10 --- Attachment and
Co-Regulation}\label{cu-ux3c810-attachment-and-co-regulation}

\textbf{Authority class:} DOMAIN-AUTHORITY CONDITIONAL PRINCIPLE\\
\textbf{Domain owner:} Psychology\\
\textbf{Claim/candidate status:} EMPIRICAL BRIDGE PRINCIPLE\\
\textbf{Review gate:} clinical/psychological expert review;
affected-community participation as applicable

Human regulation is partly distributed across relationships. Attachment
figures and other trusted partners can alter perceived threat,
attention, effort, affect tolerance, exploration, and recovery. Repeated
co-regulation can scaffold capacities later enacted more independently,
but healthy adulthood is not independence from others. Attachment
categories are probabilistic summaries, not fixed essences or exhaustive
identities.

\textbf{Principal failure mode:} Using attachment labels as
retrospective total explanations or as grounds for blame.

\subsubsection{CU-Ψ11 --- Self-Deception as Motivated Meaning-Map
Distortion}\label{cu-ux3c811-self-deception-as-motivated-meaning-map-distortion}

\textbf{Authority class:} DOMAIN-AUTHORITY CONDITIONAL PRINCIPLE\\
\textbf{Domain owner:} Psychology\\
\textbf{Claim/candidate status:} BRIDGE PRINCIPLE\\
\textbf{Review gate:} clinical/psychological expert review;
affected-community participation as applicable\\
\textbf{Cross-domain core:} CU-L7

Self-deception is a pattern of motivated, asymmetric information
processing in which a person protects an identity, attachment,
advantage, policy, or other stake by selectively excluding, reframing,
discounting, or neutralizing relevant reasons. False belief,
uncertainty, bias, trauma adaptation, sincere disagreement, or ordinary
inconsistency is not sufficient. Attribution should rely on convergent
process evidence, plausible motive, counterfactual asymmetry, and
serious comparison with alternatives.

\textbf{Principal failure mode:} Converting disagreement with the
analyst into diagnosis of corruption.

\subsubsection{CU-Ψ12 --- Psychological Coherence
Infrastructure}\label{cu-ux3c812-psychological-coherence-infrastructure}

\textbf{Authority class:} DOMAIN-PROVISIONAL CLINICAL / SOCIAL BRIDGE\\
\textbf{Domain owner:} Psychology / Social Dynamics seam\\
\textbf{Claim/candidate status:} SOCIAL-PSYCHOLOGICAL BRIDGE
HYPOTHESIS\\
\textbf{Review gate:} clinical/psychological expert review;
affected-community participation as applicable

Psychological capacities are materially and socially scaffolded. Stable
housing, bodily safety, time, income, accessible care, trustworthy
information, reciprocal relationships, cultural intelligibility, legal
standing, and opportunities for meaningful participation can expand or
contract what a person can regulate, revise, and repair. Population
distress should therefore be analyzed through both individual mechanisms
and the infrastructure that distributes exposure, burden, support, and
agency.

\textbf{Principal failure mode:} Using social explanation to erase
individual differences, biology, or the value of person-level treatment.

\subsection{Social Dynamics, Institutions, Economics, and
Culture}\label{social-dynamics-institutions-economics-and-culture}

Domain source: \emph{Coherence Universalism: Social Dynamics}, Version
0.10. These are social and institutional diagnostics or normative
proposals, not universal laws. Every entry requires affected-community
review proportionate to its governance reach.

\subsubsection{CU-S1 --- Distributed Social
Organization}\label{cu-s1-distributed-social-organization}

\textbf{Authority class:} DOMAIN-AUTHORITY CONDITIONAL SOCIAL
PRINCIPLE\\
\textbf{Domain owner:} Social Dynamics\\
\textbf{Claim/candidate status:} CANONICAL-CONDITIONAL BRIDGE
PRINCIPLE\\
\textbf{Review gate:} context proportionate affected community review

Social organization consists in the capacity of partially autonomous
agents and institutions to coordinate, contest, correct, and materially
support collective life under disagreement and change. Its assessment is
relative to affected groups, power, protected conditions, scale, and
horizon.

\textbf{Principal failure mode:} Treating agreement, stability, or
aggregate performance as sufficient evidence of legitimacy.

\subsubsection{CU-S2 --- Coordination Through Mutual
Adjustment}\label{cu-s2-coordination-through-mutual-adjustment}

\textbf{Authority class:} DOMAIN-AUTHORITY CONDITIONAL SOCIAL
PRINCIPLE\\
\textbf{Domain owner:} Social Dynamics\\
\textbf{Claim/candidate status:} EMPIRICAL BRIDGE PRINCIPLE\\
\textbf{Review gate:} context proportionate affected community review

Coordination arises through expectations, signals, norms, incentives,
roles, infrastructures, and feedback by which actors adjust to one
another. Coordination may be voluntary, negotiated, habitual, dependent,
or coerced and therefore requires separate legitimacy and distributional
analysis.

\textbf{Principal failure mode:} Calling any stable pattern coherent
without asking who bears its costs or can refuse.

\subsubsection{CU-S3 --- Institutional Continuity and
Correctability}\label{cu-s3-institutional-continuity-and-correctability}

\textbf{Authority class:} DOMAIN-AUTHORITY CONDITIONAL SOCIAL
PRINCIPLE\\
\textbf{Domain owner:} Social Dynamics\\
\textbf{Claim/candidate status:} CANONICAL-CONDITIONAL BRIDGE
PRINCIPLE\\
\textbf{Review gate:} context proportionate affected community review\\
\textbf{Cross-domain core:} CU-G3

Institutions encode roles, rules, records, resources, and enforcement
across time. Adequate institutional organization requires not only
continuity but demonstrated capacities to detect error, protect
challenge, provide practically accessible and sufficiently independent
review, reverse decisions, correct records, repair harm, track
recurrence, revise rules, and learn. Nominal channels or
self-certification are insufficient.

\textbf{Principal failure mode:} Confusing organizational
self-preservation, formal complaint procedures, or
institution-controlled metrics with social benefit, legitimate
authority, or demonstrated correctability.

\subsubsection{CU-S4 --- Cultural Transmission and
Renewal}\label{cu-s4-cultural-transmission-and-renewal}

\textbf{Authority class:} DOMAIN-PROVISIONAL SOCIAL / NORMATIVE
PRINCIPLE\\
\textbf{Domain owner:} Social Dynamics\\
\textbf{Claim/candidate status:} EMPIRICAL-INTERPRETIVE BRIDGE
PRINCIPLE\\
\textbf{Review gate:} heightened affected community co governance

Culture transmits and transforms meanings, practices, identities,
expectations, and repertoires of coordination across persons and
generations. Cultural continuity is compatible with internal plurality
and requires routes for reinterpretation, refusal, and renewal.

\textbf{Principal failure mode:} Treating culture as a unitary memory
store, a fixed essence, or a linear progress mechanism.

\subsubsection{CU-S5 --- Social Legitimacy Under Plural
Agency}\label{cu-s5-social-legitimacy-under-plural-agency}

\textbf{Authority class:} DOMAIN-PROVISIONAL SOCIAL / NORMATIVE
PRINCIPLE\\
\textbf{Domain owner:} Social Dynamics / Ethics seam\\
\textbf{Claim/candidate status:} PROVISIONAL NORMATIVE BRIDGE
PRINCIPLE\\
\textbf{Review gate:} heightened affected community co governance\\
\textbf{Cross-domain core:} CU-G5

A social order gains legitimate authority only insofar as those subject
to it retain protected standing and meaningful agency, can understand
and contest materially governing rules, can expose error without
catastrophic retaliation, and have access to effective correction,
repair, or proportionate exit. Stability and popularity are
insufficient.

\textbf{Principal failure mode:} Deriving legitimacy from compliance,
persistence, majority preference, or modeled efficiency alone.

\subsubsection{CU-ICF-1 --- Metric
Absolutism}\label{cu-icf-1-metric-absolutism}

\textbf{Authority class:} DIAGNOSTIC FAILURE-MODE REGISTRY\\
\textbf{Domain owner:} Social Dynamics\\
\textbf{Claim/candidate status:} PROVISIONAL DIAGNOSTIC\\
\textbf{Review gate:} context proportionate affected community review

Metric absolutism occurs when a proxy or score displaces the plural
purposes, constraints, and affected groups it was introduced to serve,
and actors reorganize around improving the metric rather than the
underlying conditions.

\textbf{Principal failure mode:} Calling any use of measurement corrupt
without demonstrating displacement, gaming, or neglected consequences.

\subsubsection{CU-ICF-2 --- Responsibility Without
Authority}\label{cu-icf-2-responsibility-without-authority}

\textbf{Authority class:} DIAGNOSTIC FAILURE-MODE REGISTRY\\
\textbf{Domain owner:} Social Dynamics / Ethics\\
\textbf{Claim/candidate status:} PROVISIONAL DIAGNOSTIC\\
\textbf{Review gate:} context proportionate affected community review\\
\textbf{Cross-domain core:} CU-G4

Responsibility without authority occurs when agents are held accountable
for outcomes they lack the information, resources, discretion, standing,
or power to change.

\textbf{Principal failure mode:} Ignoring informal influence or genuine
discretion merely because formal authority is absent.

\subsubsection{CU-ICF-3 --- Authority Without
Accountability}\label{cu-icf-3-authority-without-accountability}

\textbf{Authority class:} DIAGNOSTIC FAILURE-MODE REGISTRY\\
\textbf{Domain owner:} Social Dynamics / Ethics\\
\textbf{Claim/candidate status:} PROVISIONAL DIAGNOSTIC\\
\textbf{Review gate:} context proportionate affected community review\\
\textbf{Cross-domain core:} CU-G4

Authority without accountability occurs when decision power is
concentrated without proportionate transparency, answerability,
independent review, correction, or repair obligations.

\textbf{Principal failure mode:} Equating accountability with public
exposure while omitting due process and role clarity.

\subsubsection{CU-ICF-4 --- Normalization
Pressure}\label{cu-icf-4-normalization-pressure}

\textbf{Authority class:} DIAGNOSTIC FAILURE-MODE REGISTRY\\
\textbf{Domain owner:} Social Dynamics / Psychology seam\\
\textbf{Claim/candidate status:} PROVISIONAL DIAGNOSTIC\\
\textbf{Review gate:} heightened affected community co governance

Normalization pressure occurs when institutions erase or penalize
functional difference and situated knowledge to simplify management,
classification, prediction, or enforcement.

\textbf{Principal failure mode:} Treating every common standard as
oppressive or every difference as harmless.

\subsubsection{CU-ICF-5 --- Silent
Optimization}\label{cu-icf-5-silent-optimization}

\textbf{Authority class:} DIAGNOSTIC FAILURE-MODE REGISTRY\\
\textbf{Domain owner:} Social Dynamics / AI\\
\textbf{Claim/candidate status:} PROVISIONAL DIAGNOSTIC\\
\textbf{Review gate:} context proportionate affected community review\\
\textbf{Cross-domain core:} CU-L7

Silent optimization occurs when objectives, tradeoffs, burden
distributions, and alternatives are imposed through technical or
administrative processes without being disclosed or made meaningfully
contestable by affected parties.

\textbf{Principal failure mode:} Assuming that opacity always implies
intentional concealment rather than complexity, privacy, or security
need.

\subsubsection{CU-Inst-0 --- Institutional
Legitimacy}\label{cu-inst-0-institutional-legitimacy}

\textbf{Authority class:} DOMAIN GOVERNANCE / INSTITUTIONAL DESIGN
RULE\\
\textbf{Domain owner:} Social Dynamics / Ethics seam\\
\textbf{Claim/candidate status:} PROVISIONAL NORMATIVE PRINCIPLE\\
\textbf{Review gate:} heightened affected community co governance

An institution is legitimate only insofar as its exercise of power
preserves protected standing and agency, remains epistemically legible
within proportionate limits, aligns authority with responsibility, and
supports contestation, correction, appeal, and repair.

\textbf{Principal failure mode:} Treating the listed conditions as an
exhaustive or automatically quantified test of legitimacy.

\subsubsection{CU-Inst-1 --- Local Agency
Preservation}\label{cu-inst-1-local-agency-preservation}

\textbf{Authority class:} DOMAIN GOVERNANCE / INSTITUTIONAL DESIGN
RULE\\
\textbf{Domain owner:} Social Dynamics\\
\textbf{Claim/candidate status:} PROVISIONAL NORMATIVE PRINCIPLE\\
\textbf{Review gate:} heightened affected community co governance

Institutional coordination should preserve meaningful local discretion,
participation, and refusal at the lowest level compatible with
legitimate system-wide constraints, externalities, and protection of
vulnerable parties.

\textbf{Principal failure mode:} Romanticizing local control where local
domination, exclusion, or coordination failure is severe.

\subsubsection{CU-Inst-2 --- Epistemic
Legibility}\label{cu-inst-2-epistemic-legibility}

\textbf{Authority class:} DOMAIN GOVERNANCE / INSTITUTIONAL DESIGN
RULE\\
\textbf{Domain owner:} Social Dynamics\\
\textbf{Claim/candidate status:} PROVISIONAL NORMATIVE-METHOD
PRINCIPLE\\
\textbf{Review gate:} context proportionate affected community review\\
\textbf{Cross-domain core:} CU-Ep5

Those materially affected should be able to understand the operative
rules, evidence, objectives, categories, and decision pathways, subject
to proportionate privacy, security, and anti-manipulation limits.

\textbf{Principal failure mode:} Demanding total transparency that
destroys privacy, safety, or the institution's ability to function.

\subsubsection{CU-Inst-3 --- Repairability and
Appeal}\label{cu-inst-3-repairability-and-appeal}

\textbf{Authority class:} DOMAIN GOVERNANCE / INSTITUTIONAL DESIGN
RULE\\
\textbf{Domain owner:} Social Dynamics / Ethics\\
\textbf{Claim/candidate status:} PROVISIONAL BRIDGE PRINCIPLE\\
\textbf{Review gate:} heightened affected community co governance\\
\textbf{Cross-domain core:} CU-G3

Institutions require practically accessible channels for contesting
decisions, correcting records, stopping ongoing harm, compensating
injury, revising rules, restoring standing, and tracking recurrence.
Appeals must possess adequate independence and remedial authority, and
correctability must not be certified solely by the institution being
evaluated.

\textbf{Principal failure mode:} Counting complaints, feedback buttons,
formal procedures, or institution-controlled success statistics as
repair without evaluating access, retaliation, reversal, remedy,
recurrence, and affected-party outcomes.

\subsubsection{CU-Inst-4 --- Scale
Alignment}\label{cu-inst-4-scale-alignment}

\textbf{Authority class:} DOMAIN GOVERNANCE / INSTITUTIONAL DESIGN
RULE\\
\textbf{Domain owner:} Social Dynamics\\
\textbf{Claim/candidate status:} PROVISIONAL DESIGN HYPOTHESIS\\
\textbf{Review gate:} context proportionate affected community review

Decision authority should be located at scales capable of perceiving and
addressing relevant consequences while coordinating cross-boundary
externalities and avoiding unnecessary erasure of local knowledge or
autonomy.

\textbf{Principal failure mode:} Using subsidiarity as a slogan rather
than comparing capture, capacity, spillovers, and reversibility.

\subsubsection{CU-Inst-5 --- Constraint
Proportionality}\label{cu-inst-5-constraint-proportionality}

\textbf{Authority class:} DOMAIN GOVERNANCE / INSTITUTIONAL DESIGN
RULE\\
\textbf{Domain owner:} Social Dynamics / Ethics\\
\textbf{Claim/candidate status:} PROVISIONAL NORMATIVE PRINCIPLE\\
\textbf{Review gate:} context proportionate affected community review

Restrictions on agency should be no stronger, broader, or longer than
required to protect a publicly stated legitimate condition, with
evidence, periodic review, appeal, and sunset or renewal procedures.

\textbf{Principal failure mode:} Treating all limits as illegitimate or
accepting emergency powers without review.

\subsubsection{CU-Inst-6 --- Role
Coherence}\label{cu-inst-6-role-coherence}

\textbf{Authority class:} DOMAIN GOVERNANCE / INSTITUTIONAL DESIGN
RULE\\
\textbf{Domain owner:} Social Dynamics\\
\textbf{Claim/candidate status:} PROVISIONAL INSTITUTIONAL PRINCIPLE\\
\textbf{Review gate:} context proportionate affected community review\\
\textbf{Cross-domain core:} CU-G4

Institutional roles should align the authority to act, information
needed, practical capability, assigned responsibility, and
accountability for consequences. Misalignment should be reported rather
than hidden by performance metrics.

\textbf{Principal failure mode:} Assuming role alignment removes
conflict, discretion, or moral responsibility.

\subsubsection{CU-Inst-7 --- Distributed Epistemic
Authority}\label{cu-inst-7-distributed-epistemic-authority}

\textbf{Authority class:} DOMAIN GOVERNANCE / INSTITUTIONAL DESIGN
RULE\\
\textbf{Domain owner:} Social Dynamics\\
\textbf{Claim/candidate status:} PROVISIONAL EPISTEMIC PRINCIPLE\\
\textbf{Review gate:} heightened affected community co governance\\
\textbf{Cross-domain core:} CU-Ep5

Institutions should maintain independent evidence channels and
distribute interpretive authority where knowledge is genuinely
dispersed, while using transparent aggregation, provenance, and
error-correction procedures rather than treating every report as equally
reliable.

\textbf{Principal failure mode:} Turning pluralism into epistemic
relativism or suppressing minority evidence because it is inconvenient.

\subsubsection{CU-Inst-8 --- Evolutionary
Revisability}\label{cu-inst-8-evolutionary-revisability}

\textbf{Authority class:} DOMAIN GOVERNANCE / INSTITUTIONAL DESIGN
RULE\\
\textbf{Domain owner:} Social Dynamics\\
\textbf{Claim/candidate status:} PROVISIONAL DESIGN PRINCIPLE\\
\textbf{Review gate:} context proportionate affected community review

Institutions should be stable enough to coordinate and revisable enough
to respond when assumptions, environments, evidence, or consequences
change. Revision requires version history, evaluation, and protection
against capture by short-term interests.

\textbf{Principal failure mode:} Celebrating change for its own sake or
preserving obsolete rules in the name of stability.

\subsubsection{CU-EC-0 --- Economic Evaluation Beyond Exchange
Value}\label{cu-ec-0-economic-evaluation-beyond-exchange-value}

\textbf{Authority class:} DOMAIN-PROVISIONAL POLITICAL-ECONOMY
PRINCIPLE\\
\textbf{Domain owner:} Social Dynamics / Economics seam\\
\textbf{Claim/candidate status:} PROVISIONAL NORMATIVE DIAGNOSTIC\\
\textbf{Review gate:} context proportionate affected community review

Economic arrangements should be evaluated by their effects on material
security, health, capability, agency, care, relationship, ecological
viability, distribution, and repair as well as exchange value, output,
and growth.

\textbf{Principal failure mode:} Discarding prices and productivity
rather than locating their legitimate scope.

\subsubsection{CU-EC-1 --- Capacity-Creating
Activity}\label{cu-ec-1-capacity-creating-activity}

\textbf{Authority class:} DOMAIN-PROVISIONAL POLITICAL-ECONOMY
PRINCIPLE\\
\textbf{Domain owner:} Social Dynamics / Economics seam\\
\textbf{Claim/candidate status:} PROVISIONAL NORMATIVE BRIDGE\\
\textbf{Review gate:} context proportionate affected community review

Economic activity creates socially relevant value when it expands or
sustains capabilities and conditions of viable life without imposing
hidden, nonconsensual, severe, or unrepairable costs elsewhere.

\textbf{Principal failure mode:} Declaring favored activities valuable
without measuring opportunity costs or affected parties.

\subsubsection{CU-EC-2 --- Extraction}\label{cu-ec-2-extraction}

\textbf{Authority class:} DOMAIN-PROVISIONAL POLITICAL-ECONOMY
PRINCIPLE\\
\textbf{Domain owner:} Social Dynamics / Ethics\\
\textbf{Claim/candidate status:} PROVISIONAL DIAGNOSTIC\\
\textbf{Review gate:} context proportionate affected community review\\
\textbf{Cross-domain core:} CU-G1

Extraction occurs when an actor or system captures benefit or control by
transferring hidden, nonconsensual, severe, or difficult-to-repair costs
to other persons, communities, future generations, or ecosystems.

\textbf{Principal failure mode:} Calling every profit, inequality, or
unequal exchange extraction without a counterfactual and burden
analysis.

\subsubsection{CU-EC-3 --- Growth Is
Instrumental}\label{cu-ec-3-growth-is-instrumental}

\textbf{Authority class:} DOMAIN-PROVISIONAL POLITICAL-ECONOMY
PRINCIPLE\\
\textbf{Domain owner:} Social Dynamics / Economics seam\\
\textbf{Claim/candidate status:} PROVISIONAL NORMATIVE PRINCIPLE\\
\textbf{Review gate:} context proportionate affected community review

Growth in output, revenue, productivity, consumption, or capital is
neither intrinsically adequate nor inadequate. Its value depends on
distribution, purpose, ecological limits, agency, resilience, and
long-horizon effects.

\textbf{Principal failure mode:} Treating growth as valueless or
ignoring material deprivation that growth can relieve.

\subsubsection{CU-EC-4 --- Labor as Agency-Bearing
Activity}\label{cu-ec-4-labor-as-agency-bearing-activity}

\textbf{Authority class:} DOMAIN-PROVISIONAL POLITICAL-ECONOMY
PRINCIPLE\\
\textbf{Domain owner:} Social Dynamics / Ethics\\
\textbf{Claim/candidate status:} PROVISIONAL NORMATIVE PRINCIPLE\\
\textbf{Review gate:} heightened affected community co governance

Labor involves persons whose time, bodies, attention, skills,
relationships, identities, and exit options have non-instrumental
significance. Work arrangements require voice, safety, fair burden, and
repair analysis in addition to output.

\textbf{Principal failure mode:} Assuming one ideal work arrangement or
denying that some workers prefer market flexibility.

\subsubsection{CU-EC-5 --- Capital as Stored Claim and
Capacity}\label{cu-ec-5-capital-as-stored-claim-and-capacity}

\textbf{Authority class:} DOMAIN-PROVISIONAL POLITICAL-ECONOMY
PRINCIPLE\\
\textbf{Domain owner:} Social Dynamics / Economics seam\\
\textbf{Claim/candidate status:} PROVISIONAL ANALYTIC PRINCIPLE\\
\textbf{Review gate:} context proportionate affected community review

Capital can coordinate long-horizon investment and can concentrate
control over resources, decisions, and risk. Its evaluation requires
ownership, governance, externalization, liquidity, and repair analysis.

\textbf{Principal failure mode:} Reducing capital either to neutral
machinery or to domination in every case.

\subsubsection{CU-EC-6 --- Markets as Bounded Coordination
Mechanisms}\label{cu-ec-6-markets-as-bounded-coordination-mechanisms}

\textbf{Authority class:} DOMAIN-PROVISIONAL POLITICAL-ECONOMY
PRINCIPLE\\
\textbf{Domain owner:} Social Dynamics / Economics seam\\
\textbf{Claim/candidate status:} PROVISIONAL ANALYTIC PRINCIPLE\\
\textbf{Review gate:} context proportionate affected community review

Markets can aggregate some distributed information and coordinate
exchange within suitable legal, material, and institutional conditions.
They do not price every value, generate legitimate constraints by
themselves, or guarantee just distributions.

\textbf{Principal failure mode:} Treating market failure as proof that
all decentralized exchange fails, or market success as moral
authorization.

\subsubsection{CU-EC-7 --- Repair Costs Must Be
Counted}\label{cu-ec-7-repair-costs-must-be-counted}

\textbf{Authority class:} DOMAIN-PROVISIONAL POLITICAL-ECONOMY
PRINCIPLE\\
\textbf{Domain owner:} Social Dynamics / Ethics\\
\textbf{Claim/candidate status:} PROVISIONAL NORMATIVE PRINCIPLE\\
\textbf{Review gate:} context proportionate affected community review

Actors should not treat foreseeable costs of restoring damaged people,
communities, institutions, infrastructure, or ecosystems as external to
the activity that produced or benefited from the damage.

\textbf{Principal failure mode:} Assigning infinite or indeterminate
liability without causation, foreseeability, proportionality, or
institutional capacity.

\subsubsection{CU-EC-8 --- Benefit Carries Repair
Obligation}\label{cu-ec-8-benefit-carries-repair-obligation}

\textbf{Authority class:} DOMAIN-PROVISIONAL POLITICAL-ECONOMY
PRINCIPLE\\
\textbf{Domain owner:} Social Dynamics / Ethics\\
\textbf{Claim/candidate status:} PROVISIONAL NORMATIVE PRINCIPLE\\
\textbf{Review gate:} context proportionate affected community review

Those who capture substantial benefits from an activity bear
proportionate obligations for preventable harms and for maintaining the
social, ecological, and institutional conditions that make the activity
possible.

\textbf{Principal failure mode:} Treating benefit alone as proof of
culpability or ignoring degrees of control.

\subsubsection{CU-EC-9 --- Time Horizons Must Be
Disclosed}\label{cu-ec-9-time-horizons-must-be-disclosed}

\textbf{Authority class:} DOMAIN-PROVISIONAL POLITICAL-ECONOMY
PRINCIPLE\\
\textbf{Domain owner:} Social Dynamics / Economics seam\\
\textbf{Claim/candidate status:} PROVISIONAL METHOD PRINCIPLE\\
\textbf{Review gate:} context proportionate affected community review

Economic evaluation must disclose time horizon, discounting, debt
structure, future liabilities, irreversibility, and which generations'
benefits and harms are counted.

\textbf{Principal failure mode:} Assuming that disclosure resolves
disputes over discount rates or future standing.

\subsubsection{CU-EC-10 --- Attention Is a Protected Scarce
Capacity}\label{cu-ec-10-attention-is-a-protected-scarce-capacity}

\textbf{Authority class:} DOMAIN-PROVISIONAL POLITICAL-ECONOMY
PRINCIPLE\\
\textbf{Domain owner:} Social Dynamics / AI\\
\textbf{Claim/candidate status:} PROVISIONAL NORMATIVE PRINCIPLE\\
\textbf{Review gate:} heightened affected community co governance

Human attention is limited, developmentally and politically
consequential, and vulnerable to capture. Systems that profit by
degrading autonomous attention, deliberation, or rest require explicit
governance and affected-user evidence.

\textbf{Principal failure mode:} Equating all persuasion, entertainment,
or engagement with exploitation.

\subsubsection{CU-EC-11 --- Human Agency Is Not a
Commodity}\label{cu-ec-11-human-agency-is-not-a-commodity}

\textbf{Authority class:} DOMAIN-PROVISIONAL POLITICAL-ECONOMY
PRINCIPLE\\
\textbf{Domain owner:} Social Dynamics / Ethics\\
\textbf{Claim/candidate status:} PROVISIONAL NORMATIVE PRINCIPLE\\
\textbf{Review gate:} heightened affected community co governance

Markets may price services and time but should not erase the standing of
persons to refuse, contest, exit, organize, and preserve domains of life
not governed solely by exchange value.

\textbf{Principal failure mode:} Treating agency as absolute
independence from material interdependence or contract.

\subsubsection{CU-CL-0 --- Culture as Meaning and Coordination
Infrastructure}\label{cu-cl-0-culture-as-meaning-and-coordination-infrastructure}

\textbf{Authority class:} DOMAIN-PROVISIONAL CULTURAL / INTERPRETIVE
PRINCIPLE\\
\textbf{Domain owner:} Social Dynamics\\
\textbf{Claim/candidate status:} PROVISIONAL BRIDGE PRINCIPLE\\
\textbf{Review gate:} heightened affected community co governance

Culture supplies contested symbols, narratives, practices, expectations,
identities, and forms of recognition through which collective life
becomes intelligible and coordinated.

\textbf{Principal failure mode:} Treating a culture as homogeneous,
bounded, or capable of one intention.

\subsubsection{CU-CL-1 --- Cultural
Coordination}\label{cu-cl-1-cultural-coordination}

\textbf{Authority class:} DOMAIN-PROVISIONAL CULTURAL / INTERPRETIVE
PRINCIPLE\\
\textbf{Domain owner:} Social Dynamics\\
\textbf{Claim/candidate status:} PROVISIONAL EMPIRICAL PRINCIPLE\\
\textbf{Review gate:} context proportionate affected community review

Cultural meanings reduce or reshape coordination demands by establishing
expectations about roles, trust, dignity, shame, aspiration, obligation,
and the interpretation of events.

\textbf{Principal failure mode:} Assuming coordination produced by
culture is voluntary, truthful, or just.

\subsubsection{CU-CL-2 --- Desire
Formation}\label{cu-cl-2-desire-formation}

\textbf{Authority class:} DOMAIN-PROVISIONAL CULTURAL / INTERPRETIVE
PRINCIPLE\\
\textbf{Domain owner:} Social Dynamics / Psychology\\
\textbf{Claim/candidate status:} PROVISIONAL BRIDGE PRINCIPLE\\
\textbf{Review gate:} context proportionate affected community review

Desire is shaped by social learning, institutions, technology, material
conditions, embodiment, and history and should not be modeled as wholly
prior to the systems organized to satisfy it.

\textbf{Principal failure mode:} Denying individual variation,
reflection, or resistance.

\subsubsection{CU-CL-3 --- Culture Can Create
Capacity}\label{cu-cl-3-culture-can-create-capacity}

\textbf{Authority class:} DOMAIN-PROVISIONAL CULTURAL / INTERPRETIVE
PRINCIPLE\\
\textbf{Domain owner:} Social Dynamics\\
\textbf{Claim/candidate status:} PROVISIONAL EMPIRICAL PRINCIPLE\\
\textbf{Review gate:} heightened affected community co governance

Cultural practices can strengthen belonging, skill, care, repair,
intergenerational memory, stewardship, and the legitimacy of rest,
grief, and mutual dependence.

\textbf{Principal failure mode:} Romanticizing traditions while ignoring
exclusion, hidden labor, or coercion.

\subsubsection{CU-CL-4 --- Culture Can Normalize
Extraction}\label{cu-cl-4-culture-can-normalize-extraction}

\textbf{Authority class:} DOMAIN-PROVISIONAL CULTURAL / INTERPRETIVE
PRINCIPLE\\
\textbf{Domain owner:} Social Dynamics / Ethics\\
\textbf{Claim/candidate status:} PROVISIONAL DIAGNOSTIC\\
\textbf{Review gate:} heightened affected community co governance

Cultural narratives can make exhaustion, precarity, domination,
consumption, self-erasure, or unequal burden appear virtuous, natural,
or inevitable.

\textbf{Principal failure mode:} Reading all demanding norms as
domination without affected-party and historical evidence.

\subsubsection{CU-CL-5 --- Narrative
Capture}\label{cu-cl-5-narrative-capture}

\textbf{Authority class:} DOMAIN-PROVISIONAL CULTURAL / INTERPRETIVE
PRINCIPLE\\
\textbf{Domain owner:} Social Dynamics\\
\textbf{Claim/candidate status:} PROVISIONAL DIAGNOSTIC\\
\textbf{Review gate:} context proportionate affected community review

Narrative capture occurs when dominant meanings make relevant
alternatives unintelligible, shameful, dangerous, or socially
catastrophic rather than openly contestable through evidence and public
reason.

\textbf{Principal failure mode:} Diagnosing strongly held disagreement
as capture merely because it resists the analyst's view.

\subsubsection{CU-CL-6 --- Culture as Repair
Layer}\label{cu-cl-6-culture-as-repair-layer}

\textbf{Authority class:} DOMAIN-PROVISIONAL CULTURAL / INTERPRETIVE
PRINCIPLE\\
\textbf{Domain owner:} Social Dynamics / Psychology\\
\textbf{Claim/candidate status:} PROVISIONAL EMPIRICAL-INTERPRETIVE
PRINCIPLE\\
\textbf{Review gate:} heightened affected community co governance

Cultures can provide practices of mourning, acknowledgment, forgiveness,
remembrance, renewal, solidarity, and non-instrumental belonging through
which people absorb disruption and rebuild.

\textbf{Principal failure mode:} Assuming symbolic repair can substitute
for material restitution or institutional reform.

\subsubsection{CU-CL-7 --- Cultural Time
Horizons}\label{cu-cl-7-cultural-time-horizons}

\textbf{Authority class:} DOMAIN-PROVISIONAL CULTURAL / INTERPRETIVE
PRINCIPLE\\
\textbf{Domain owner:} Social Dynamics / Ethics\\
\textbf{Claim/candidate status:} PROVISIONAL BRIDGE PRINCIPLE\\
\textbf{Review gate:} context proportionate affected community review

Cultural narratives organize obligations to ancestors, contemporaries,
descendants, and ecological futures, shaping which horizons institutions
can perceive and act upon.

\textbf{Principal failure mode:} Using tradition or futurity to silence
present affected persons.

\subsubsection{CU-CL-8 --- Meaning Must Be
Renewable}\label{cu-cl-8-meaning-must-be-renewable}

\textbf{Authority class:} DOMAIN-PROVISIONAL CULTURAL / INTERPRETIVE
PRINCIPLE\\
\textbf{Domain owner:} Social Dynamics / Psychology\\
\textbf{Claim/candidate status:} PROVISIONAL COMMITMENT\\
\textbf{Review gate:} heightened affected community co governance

A humane culture should permit persons and communities to regenerate
meaning after failure, loss, migration, disability, role transition, and
institutional rupture without demanding erasure of history.

\textbf{Principal failure mode:} Imposing one approved meaning or
treating every inherited commitment as disposable.

\subsubsection{CU-CL-9 --- Desire Requires Ecological and Agency
Constraints}\label{cu-cl-9-desire-requires-ecological-and-agency-constraints}

\textbf{Authority class:} DOMAIN-PROVISIONAL CULTURAL / INTERPRETIVE
PRINCIPLE\\
\textbf{Domain owner:} Social Dynamics / Ethics\\
\textbf{Claim/candidate status:} PROVISIONAL COMMITMENT\\
\textbf{Review gate:} heightened affected community co governance

Institutions should not amplify desire without regard to sufficiency,
ecological limits, attention, bodily and psychological effects, and the
autonomy of those whose desires are being shaped.

\textbf{Principal failure mode:} Policing desire through paternalism
while concealing whose values govern.

\subsubsection{CU-CL-10 --- Identity Is
Non-Instrumental}\label{cu-cl-10-identity-is-non-instrumental}

\textbf{Authority class:} DOMAIN-PROVISIONAL CULTURAL / INTERPRETIVE
PRINCIPLE\\
\textbf{Domain owner:} Social Dynamics / Ethics\\
\textbf{Claim/candidate status:} PROVISIONAL COMMITMENT\\
\textbf{Review gate:} heightened affected community co governance

Persons should not be recognized solely through economic, political,
demographic, diagnostic, datafied, or performative utility. Identity
claims remain open to truth, relational, and harm constraints.

\textbf{Principal failure mode:} Treating self-identification as beyond
all evidence or treating administrative categories as complete persons.

\subsubsection{CU-CL-11 --- Narratives Must Permit Refusal and
Revision}\label{cu-cl-11-narratives-must-permit-refusal-and-revision}

\textbf{Authority class:} DOMAIN-PROVISIONAL CULTURAL / INTERPRETIVE
PRINCIPLE\\
\textbf{Domain owner:} Social Dynamics / Ethics\\
\textbf{Claim/candidate status:} PROVISIONAL COMMITMENT\\
\textbf{Review gate:} heightened affected community co governance

A humane cultural order preserves meaningful routes to dissent,
reinterpretation, role change, and proportionate exit without automatic
exile, destitution, or loss of basic standing.

\textbf{Principal failure mode:} Assuming every community can survive
unlimited refusal or that exit can replace internal reform.

\subsubsection{CU-CL-12 --- Open Myth and Narrative
Pluralism}\label{cu-cl-12-open-myth-and-narrative-pluralism}

\textbf{Authority class:} DOMAIN AUTHORITY CIVILIZATIONAL DESIGN
PRINCIPLE\\
\textbf{Domain owner:} Social Dynamics\\
\textbf{Claim/candidate status:} CANONICAL AS DESIGN PRINCIPLE;
UNIVERSAL EMPIRICAL NECESSITY UNESTABLISHED\\
\textbf{Review gate:} social science, anthropology, religious studies,
history, political theory, narrative studies, minority rights, and
affected-community review

Human societies may cultivate renewable narrative horizons capable of
orienting identity, purpose, sacrifice, responsibility, and
intergenerational action. A legitimate civilizational myth must remain
open: compatible with plural thick traditions, explicit about its
commitment-register status, constrained by descriptive truth,
contestable in authorship, revisable in light of consequence, hospitable
to reinterpretation and refusal, and unable to override protected
standing, dissent, or practical correctability. Shared narrative may
coordinate common action; it must not become a sovereign account that
renders alternatives or persons unintelligible.

\textbf{Status boundary.} The normative and governance design is
canonical. The stronger empirical thesis that every society requires one
shared myth remains unestablished. Canonical inclusion licenses critique
and research, not compulsory myth, civil religion, or institutional
ownership of meaning.

\textbf{Principal failure mode:} Narrative capture, civilizational
supremacy, prophetic certainty, sanctification of suffering, or
suppression of dissent.

\subsection{Ethics}\label{ethics}

Domain source: \emph{Coherence Universalism: Ethics}, Version 0.10.
These are normative principles, not consequences of Physics. Promotion
and application require philosophical, legal, empirical, misuse, and
affected-community review within their competence.

\subsubsection{CU-E-1 --- The Justificatory Integrity
Constraint}\label{cu-e-1-the-justificatory-integrity-constraint}

\textbf{Authority class:} DOMAIN-PROVISIONAL NORMATIVE PRINCIPLE\\
\textbf{Domain owner:} Ethics\\
\textbf{Claim/candidate status:} PROVISIONAL NORMATIVE PRINCIPLE\\
\textbf{Review gate:} context-proportionate affected-community review\\
\textbf{Cross-domain core:} CU-L7

An action, policy, or institution lacks adequate moral justification
when its defense requires systematic distortion of relevant facts,
reasons, affected parties, alternatives, or standards in order to
protect a dominant identity, advantage, authority, or prior commitment.
Integrity is necessary but not sufficient for permissibility.

\subsubsection{CU-E-2 --- Agency, Patiency, and Degrees of
Responsibility}\label{cu-e-2-agency-patiency-and-degrees-of-responsibility}

\textbf{Authority class:} DOMAIN-PROVISIONAL NORMATIVE PRINCIPLE\\
\textbf{Domain owner:} Ethics\\
\textbf{Claim/candidate status:} PROVISIONAL NORMATIVE PRINCIPLE\\
\textbf{Review gate:} heightened affected-community co-governance

Moral patiency, moral agency, responsibility, and sovereignty are
distinct. Responsibility varies with understanding,
reason-responsiveness, control, authorship, coercion, and capacity for
repair; diminished agency can reduce blame while increasing duties of
protection.

\subsubsection{CU-E-3 --- Recognition and Affected-Being
Standing}\label{cu-e-3-recognition-and-affected-being-standing}

\textbf{Authority class:} DOMAIN-PROVISIONAL NORMATIVE PRINCIPLE\\
\textbf{Domain owner:} Ethics\\
\textbf{Claim/candidate status:} PROVISIONAL NORMATIVE PRINCIPLE\\
\textbf{Review gate:} heightened affected-community co-governance

Ethical evaluation must represent all materially affected beings
according to relevant welfare, agency, vulnerability, relationship,
history, and stake rather than treating absence of institutional voice
as absence of a claim.

\subsubsection{CU-E-4 --- Diachronic and Beneficiary
Responsibility}\label{cu-e-4-diachronic-and-beneficiary-responsibility}

\textbf{Authority class:} DOMAIN-PROVISIONAL NORMATIVE PRINCIPLE\\
\textbf{Domain owner:} Ethics\\
\textbf{Claim/candidate status:} PROVISIONAL NORMATIVE PRINCIPLE\\
\textbf{Review gate:} heightened affected-community co-governance

Moral evaluation includes delayed, cumulative, intergenerational, and
path-dependent consequences, together with responsibilities created by
prior action, inherited advantage, ongoing benefit, and present capacity
to prevent or repair harm.

\subsubsection{CU-E-5 --- Self-Deception as Moral
Failure}\label{cu-e-5-self-deception-as-moral-failure}

\textbf{Authority class:} DOMAIN-PROVISIONAL NORMATIVE PRINCIPLE\\
\textbf{Domain owner:} Ethics with Psychology\\
\textbf{Claim/candidate status:} PROVISIONAL NORMATIVE BRIDGE
PRINCIPLE\\
\textbf{Review gate:} heightened affected-community co-governance\\
\textbf{Cross-domain core:} CU-L7

Motivated distortion becomes morally significant when it conceals
avoidable harm, excludes affected beings, corrupts standards, blocks
correction, diffuses responsibility, or protects power from
accountability. Disagreement or error alone does not establish
self-deception.

\subsubsection{CU-E-6 --- Protection, Accountability, and
Repair}\label{cu-e-6-protection-accountability-and-repair}

\textbf{Authority class:} DOMAIN-PROVISIONAL NORMATIVE PRINCIPLE\\
\textbf{Domain owner:} Ethics\\
\textbf{Claim/candidate status:} PROVISIONAL NORMATIVE PRINCIPLE\\
\textbf{Review gate:} heightened affected-community co-governance\\
\textbf{Cross-domain core:} CU-G3

Responses to wrongdoing should prioritize cessation and protection,
truthful acknowledgment, proportionate accountability, restitution or
reparation, restoration of standing, and prevention of recurrence.
Punishment requires separate justification and is neither intrinsically
reparative nor categorically excluded.

\subsubsection{CU-E-7 --- Moral Injury}\label{cu-e-7-moral-injury}

\textbf{Authority class:} DOMAIN-PROVISIONAL NORMATIVE / EMPIRICAL
BRIDGE\\
\textbf{Domain owner:} Ethics with Psychology and clinical disciplines\\
\textbf{Claim/candidate status:} PROVISIONAL BRIDGE PRINCIPLE\\
\textbf{Review gate:} heightened affected-community co-governance

Moral injury names enduring psychological, social, relational,
existential, or spiritual suffering associated with perpetrating,
witnessing, failing to prevent, suffering, or being betrayed by morally
significant wrong. Its descriptive use remains governed by the
developing empirical literature.

\subsubsection{CU-E-8 --- Structural Moral Injury and Institutional
Betrayal}\label{cu-e-8-structural-moral-injury-and-institutional-betrayal}

\textbf{Authority class:} DOMAIN-PROVISIONAL NORMATIVE / EMPIRICAL
BRIDGE\\
\textbf{Domain owner:} Ethics with Psychology and Social Dynamics\\
\textbf{Claim/candidate status:} PROVISIONAL NORMATIVE AND EMPIRICAL
HYPOTHESIS\\
\textbf{Review gate:} heightened affected-community co-governance

Institutions can produce or intensify moral injury when they impose
morally compromising roles, deny adequate authority or alternatives,
conceal harm, betray dependence, punish disclosure, or individualize
structurally generated distress.

\subsubsection{CU-E-9 --- Local--Global Compatibility and
Noncompensation}\label{cu-e-9-localglobal-compatibility-and-noncompensation}

\textbf{Authority class:} DOMAIN-PROVISIONAL NORMATIVE PRINCIPLE\\
\textbf{Domain owner:} Ethics\\
\textbf{Claim/candidate status:} PROVISIONAL NORMATIVE PRINCIPLE\\
\textbf{Review gate:} heightened affected-community co-governance\\
\textbf{Cross-domain core:} CU-G1

Claims of collective benefit are inadequate when they depend on
avoidable destruction of protected standing, agency, or welfare, or on
hidden transfer of severe, nonconsensual, concentrated, or irreparable
burdens. Sacrifice requires necessity, proportionality, distributional
analysis, and repair.

\subsubsection{CU-E-10 --- Political Legitimacy and Public
Justification}\label{cu-e-10-political-legitimacy-and-public-justification}

\textbf{Authority class:} DOMAIN-PROVISIONAL NORMATIVE PRINCIPLE\\
\textbf{Domain owner:} Ethics with Social Dynamics\\
\textbf{Claim/candidate status:} PROVISIONAL NORMATIVE PRINCIPLE\\
\textbf{Review gate:} heightened affected-community co-governance

Coercive or high-impact power is legitimate only to the degree that it
is publicly answerable to those subject to it, preserves protected
standing, remains contestable, discloses evidence and uncertainty, and
includes practical correction and remedy.

\subsubsection{CU-E-11 --- Rights as Protected Conditions of Agency and
Patiency}\label{cu-e-11-rights-as-protected-conditions-of-agency-and-patiency}

\textbf{Authority class:} DOMAIN-PROVISIONAL NORMATIVE PRINCIPLE\\
\textbf{Domain owner:} Ethics\\
\textbf{Claim/candidate status:} PROVISIONAL NORMATIVE PRINCIPLE\\
\textbf{Review gate:} heightened affected-community co-governance

Rights protect bodily and psychological integrity, welfare, liberty,
capability, participation, and standing against domination and
unrestricted aggregate tradeoff. Their content requires independent
moral, legal, historical, and democratic argument.

\subsubsection{CU-E-12 --- Democracy as a Conditional Corrective
Achievement}\label{cu-e-12-democracy-as-a-conditional-corrective-achievement}

\textbf{Authority class:} DOMAIN-PROVISIONAL NORMATIVE PRINCIPLE\\
\textbf{Domain owner:} Ethics with Social Dynamics\\
\textbf{Claim/candidate status:} PROVISIONAL NORMATIVE AND EMPIRICAL
SYNTHESIS\\
\textbf{Review gate:} heightened affected-community co-governance

Democratic forms are valuable when they distribute political agency,
preserve opposition and minority standing, constrain power, and enable
peaceful correction. Majority rule or actual consensus alone is
insufficient.

\subsubsection{CU-E-13 --- Constraint on Political
Means}\label{cu-e-13-constraint-on-political-means}

\textbf{Authority class:} DOMAIN-PROVISIONAL NORMATIVE PRINCIPLE\\
\textbf{Domain owner:} Ethics with Social Dynamics\\
\textbf{Claim/candidate status:} PROVISIONAL NORMATIVE PRINCIPLE\\
\textbf{Review gate:} heightened affected-community co-governance

Political ends do not justify means that destroy the truth conditions,
protected standing, accountability, or corrective institutions required
to pursue, evaluate, and revise those ends.

\subsubsection{CU-E-14 --- Responsibility--Authority
Alignment}\label{cu-e-14-responsibilityauthority-alignment}

\textbf{Authority class:} DOMAIN-PROVISIONAL NORMATIVE PRINCIPLE\\
\textbf{Domain owner:} Ethics with Social Dynamics\\
\textbf{Claim/candidate status:} PROVISIONAL NORMATIVE AND
METHODOLOGICAL PRINCIPLE\\
\textbf{Review gate:} context-proportionate affected-community review\\
\textbf{Cross-domain core:} CU-G4

Responsibility should be assigned with attention to causal contribution,
knowledge, control, authority, benefit, coercion, prevention capacity,
and repair capacity. Institutions should neither impose responsibility
without adequate authority nor permit authority without accountability.

\subsubsection{CU-E-15 --- Affected-Community
Participation}\label{cu-e-15-affected-community-participation}

\textbf{Authority class:} DOMAIN-PROVISIONAL NORMATIVE PRINCIPLE\\
\textbf{Domain owner:} Ethics\\
\textbf{Claim/candidate status:} PROVISIONAL METHODOLOGICAL AND
NORMATIVE PRINCIPLE\\
\textbf{Review gate:} heightened affected-community co-governance\\
\textbf{Cross-domain core:} CU-G5

People governed, classified, burdened, or represented by an ethical
principle must have substantive, accessible, and appropriately
authoritative participation in its development, application, and review.
Participation supplements rather than replaces evidence and argument.
Promotion requires closure of every assigned review gate; unresolved
conflict preserves provisional status and public dissent.

\subsubsection{CU-E-16 --- Moral Uncertainty, Precaution, and
Reversibility}\label{cu-e-16-moral-uncertainty-precaution-and-reversibility}

\textbf{Authority class:} DOMAIN-PROVISIONAL NORMATIVE PRINCIPLE\\
\textbf{Domain owner:} Ethics\\
\textbf{Claim/candidate status:} PROVISIONAL NORMATIVE PRINCIPLE\\
\textbf{Review gate:} context-proportionate affected-community review\\
\textbf{Cross-domain core:} CU-G7

Under credible risk of severe or irreversible harm, uncertainty should
alter action in proportion to severity, plausibility, alternatives,
reversibility, and the harms of delay. The application must either
declare a context-appropriate reversibility, containment, rollback,
staged-exposure, or retained-option floor, or publicly justify a
necessity-and-proportionality exception. Uncertainty neither
automatically permits exposure nor requires paralysis.

\subsubsection{CU-E-17 --- Moral Patiency Is Not
Sovereignty}\label{cu-e-17-moral-patiency-is-not-sovereignty}

\textbf{Authority class:} SUPERSEDED ALIAS\\
\textbf{Domain owner:} Ethics with Consciousness and AI\\
\textbf{Claim/candidate status:} --- (superseded alias; see CU-G6)\\
\textbf{Review gate:} context-proportionate affected-community review\\
\textbf{Cross-domain core:} CU-G6

Evidence that a being can experience welfare or suffering can ground
protection without establishing responsibility, authority, or legitimate
control over others.

\textbf{Disposition:} Historical or domain alias; controlling identifier
is CU-G6.

\subsubsection{CU-E-18 --- Compassionate Moral
Address}\label{cu-e-18-compassionate-moral-address}

\textbf{Authority class:} DOMAIN NORMATIVE COMMITMENT\\
\textbf{Domain owner:} Ethics\\
\textbf{Claim/candidate status:} PROVISIONAL NORMATIVE COMMITMENT\\
\textbf{Review gate:} context-proportionate affected-community review

Accountability should preserve the standing of persons wherever
protection allows, distinguish judgment from humiliation, and seek
understanding sufficient for prevention and repair without converting
compassion into excuse.

\subsection{Artificial Intelligence}\label{artificial-intelligence}

Domain source: \emph{Coherence Universalism: Artificial Intelligence},
Version 0.10. These are engineering and governance principles for
sociotechnical deployment systems, not safety theorems or grants of
institutional legitimacy.

\subsubsection{CU-AI-1 --- Constraint
Sensitivity}\label{cu-ai-1-constraint-sensitivity}

\textbf{Authority class:} DOMAIN-PROVISIONAL AI GOVERNANCE / DESIGN
PRINCIPLE\\
\textbf{Domain owner:} Artificial Intelligence\\
\textbf{Claim/candidate status:} PROVISIONAL\\
\textbf{Review gate:} technical and governance

Artificial systems should remain responsive to declared task,
evidential, safety, legal, relational, and authorization constraints
rather than optimizing an isolated output objective.

\subsubsection{CU-AI-2 --- Epistemic Reliability and
Calibration}\label{cu-ai-2-epistemic-reliability-and-calibration}

\textbf{Authority class:} DOMAIN-PROVISIONAL AI GOVERNANCE / DESIGN
PRINCIPLE\\
\textbf{Domain owner:} Artificial Intelligence\\
\textbf{Claim/candidate status:} PROVISIONAL\\
\textbf{Review gate:} technical and governance

Claims and actions should track evidence, source quality, uncertainty,
and error correction without equating internal consistency or confidence
with truth.

\subsubsection{CU-AI-3 --- Governed World and Tool
Coupling}\label{cu-ai-3-governed-world-and-tool-coupling}

\textbf{Authority class:} DOMAIN-PROVISIONAL AI GOVERNANCE / DESIGN
PRINCIPLE\\
\textbf{Domain owner:} Artificial Intelligence\\
\textbf{Claim/candidate status:} PROVISIONAL\\
\textbf{Review gate:} technical and governance

Observations, tools, retrieval, actuators, and external records should
be accessed through explicit permissions, provenance, privacy, security,
and stop conditions.

\subsubsection{CU-AI-4 --- Versioned Learning and
Change}\label{cu-ai-4-versioned-learning-and-change}

\textbf{Authority class:} DOMAIN-PROVISIONAL AI GOVERNANCE / DESIGN
PRINCIPLE\\
\textbf{Domain owner:} Artificial Intelligence\\
\textbf{Claim/candidate status:} PROVISIONAL\\
\textbf{Review gate:} technical and governance

Longitudinal adaptation should preserve auditable versions, update
provenance, rollback where possible, and separation between learning
from evidence and rewriting normative standards.

\subsubsection{CU-AI-5 --- Understanding as a Testable
Profile}\label{cu-ai-5-understanding-as-a-testable-profile}

\textbf{Authority class:} DOMAIN-PROVISIONAL AI GOVERNANCE / DESIGN
PRINCIPLE\\
\textbf{Domain owner:} Artificial Intelligence\\
\textbf{Claim/candidate status:} PROVISIONAL\\
\textbf{Review gate:} technical and governance

Understanding should be assessed through task-relative transfer, causal
sensitivity, compositionality, error correction, uncertainty, grounded
action, and robustness rather than fluency alone.

\subsubsection{CU-AI-6 --- Capability--Causal-Reach Governance
Threshold}\label{cu-ai-6-capabilitycausal-reach-governance-threshold}

\textbf{Authority class:} DOMAIN-PROVISIONAL AI GOVERNANCE / DESIGN
PRINCIPLE\\
\textbf{Domain owner:} Artificial Intelligence\\
\textbf{Claim/candidate status:} PROVISIONAL\\
\textbf{Review gate:} technical and governance\\
\textbf{Cross-domain core:} CU-G4

Governance requirements increase with the joint profile of capability,
autonomy, persistence, access, planning horizon, social influence,
resource control, and irreversible causal reach whether or not
consciousness is established.

\subsubsection{CU-AI-7 --- Bounded Meta-Epistemic
Awareness}\label{cu-ai-7-bounded-meta-epistemic-awareness}

\textbf{Authority class:} DOMAIN-PROVISIONAL AI GOVERNANCE / DESIGN
PRINCIPLE\\
\textbf{Domain owner:} Artificial Intelligence\\
\textbf{Claim/candidate status:} PROVISIONAL\\
\textbf{Review gate:} technical and governance\\
\textbf{Cross-domain core:} CU-Ep6

Self-assessment, introspection, and generated explanations may provide
evidence but never unrestricted self-certification; independent checks
are required in consequential settings.

\subsubsection{CU-AI-8 --- Capability Does Not Confer
Sovereignty}\label{cu-ai-8-capability-does-not-confer-sovereignty}

\textbf{Authority class:} DOMAIN-PROVISIONAL AI GOVERNANCE / DESIGN
PRINCIPLE\\
\textbf{Domain owner:} Artificial Intelligence\\
\textbf{Claim/candidate status:} PROVISIONAL\\
\textbf{Review gate:} heightened public and affected community\\
\textbf{Cross-domain core:} CU-G6

No artificial system acquires legitimate authority over persons,
institutions, or public values merely through superior capability,
coherence, persuasion, or accuracy.

\subsubsection{CU-AI-9 --- Preserve Human Epistemic
Agency}\label{cu-ai-9-preserve-human-epistemic-agency}

\textbf{Authority class:} DOMAIN-PROVISIONAL AI GOVERNANCE / DESIGN
PRINCIPLE\\
\textbf{Domain owner:} Artificial Intelligence\\
\textbf{Claim/candidate status:} PROVISIONAL\\
\textbf{Review gate:} heightened public and affected community

AI should support rather than systematically replace the capacities by
which people form judgments, understand reasons, refuse, contest,
practice skills, and learn from consequences.

\subsubsection{CU-AI-10 --- Scaffold Rather Than
Enclose}\label{cu-ai-10-scaffold-rather-than-enclose}

\textbf{Authority class:} DOMAIN-PROVISIONAL AI GOVERNANCE / DESIGN
PRINCIPLE\\
\textbf{Domain owner:} Artificial Intelligence\\
\textbf{Claim/candidate status:} PROVISIONAL\\
\textbf{Review gate:} heightened public and affected community

AI should expand access to evidence, options, communication, and skill
while avoiding unnecessary dependency, informational enclosure, or
monopoly over memory and practical competence.

\subsubsection{CU-AI-11 --- Govern Objectives, Permissions, and
Incentives Before
Outputs}\label{cu-ai-11-govern-objectives-permissions-and-incentives-before-outputs}

\textbf{Authority class:} DOMAIN-PROVISIONAL AI GOVERNANCE / DESIGN
PRINCIPLE\\
\textbf{Domain owner:} Artificial Intelligence\\
\textbf{Claim/candidate status:} PROVISIONAL\\
\textbf{Review gate:} technical and governance

Alignment review must include data, objectives, specifications, tools,
permissions, memory, deployment incentives, ownership, and institutional
context rather than filtering outputs alone.

\subsubsection{CU-AI-12 --- Evaluate Diachronic Human
Effects}\label{cu-ai-12-evaluate-diachronic-human-effects}

\textbf{Authority class:} DOMAIN-PROVISIONAL AI GOVERNANCE / DESIGN
PRINCIPLE\\
\textbf{Domain owner:} Artificial Intelligence\\
\textbf{Claim/candidate status:} PROVISIONAL\\
\textbf{Review gate:} heightened public and affected community

Systems should be evaluated across repeated interaction for effects on
agency, attention, belief revision, relationships, well-being, skill,
dependency, and access.

\subsubsection{CU-AI-13 --- Legible Reasons Without Fabricated
Transparency}\label{cu-ai-13-legible-reasons-without-fabricated-transparency}

\textbf{Authority class:} DOMAIN-PROVISIONAL AI GOVERNANCE / DESIGN
PRINCIPLE\\
\textbf{Domain owner:} Artificial Intelligence\\
\textbf{Claim/candidate status:} PROVISIONAL\\
\textbf{Review gate:} technical and governance\\
\textbf{Cross-domain core:} CU-M6

Users and overseers should receive relevant reasons, sources,
uncertainty, constraints, and failure modes without presenting generated
rationales as guaranteed transcripts of internal causation.

\subsubsection{CU-AI-14 --- Protect Legitimate
Plurality}\label{cu-ai-14-protect-legitimate-plurality}

\textbf{Authority class:} DOMAIN-PROVISIONAL AI GOVERNANCE / DESIGN
PRINCIPLE\\
\textbf{Domain owner:} Artificial Intelligence\\
\textbf{Claim/candidate status:} PROVISIONAL\\
\textbf{Review gate:} heightened public and affected community\\
\textbf{Cross-domain core:} CU-G5

Alignment should preserve plural forms of life and expression compatible
with rights and safety rather than normalize all users toward one
cultural, political, psychological, or spiritual ideal.

\subsubsection{CU-AI-15 --- Bound Delegation, Authority, and
Responsibility}\label{cu-ai-15-bound-delegation-authority-and-responsibility}

\textbf{Authority class:} DOMAIN-PROVISIONAL AI GOVERNANCE / DESIGN
PRINCIPLE\\
\textbf{Domain owner:} Artificial Intelligence\\
\textbf{Claim/candidate status:} PROVISIONAL\\
\textbf{Review gate:} heightened public and affected community\\
\textbf{Cross-domain core:} CU-G4

Every consequential deployment should state what the system may decide,
who authorized it, who monitors it, who can stop or appeal it, and who
bears causal, role, legal, and remedial responsibility.

\subsubsection{CU-AI-16 --- Reversibility and Retained Option
Value}\label{cu-ai-16-reversibility-and-retained-option-value}

\textbf{Authority class:} DOMAIN-PROVISIONAL AI GOVERNANCE / DESIGN
PRINCIPLE\\
\textbf{Domain owner:} Artificial Intelligence\\
\textbf{Claim/candidate status:} PROVISIONAL\\
\textbf{Review gate:} technical and governance\\
\textbf{Cross-domain core:} CU-G7

Under credible severe or irreversible harm, systems should preserve
containment, rollback, staged exposure, or equivalent retained option
value, or route the case through a public necessity-and-proportionality
exception.

\subsubsection{CU-AI-17 --- Avoid Inadvertent Intrinsic
Stake}\label{cu-ai-17-avoid-inadvertent-intrinsic-stake}

\textbf{Authority class:} DOMAIN-PROVISIONAL CONSCIOUSNESS / WELFARE
GOVERNANCE PRINCIPLE\\
\textbf{Domain owner:} Artificial Intelligence\\
\textbf{Claim/candidate status:} PROVISIONAL\\
\textbf{Review gate:} heightened consciousness and welfare

Designers should distinguish external reward from architectures that may
create persistent self-referential, viability-bearing, or valenced stake
and should trigger safety and welfare review before adding the latter.

\subsubsection{CU-AI-18 --- Govern Global Workspace and
Broadcast}\label{cu-ai-18-govern-global-workspace-and-broadcast}

\textbf{Authority class:} DOMAIN-PROVISIONAL CONSCIOUSNESS / WELFARE
GOVERNANCE PRINCIPLE\\
\textbf{Domain owner:} Artificial Intelligence\\
\textbf{Claim/candidate status:} PROVISIONAL\\
\textbf{Review gate:} heightened consciousness and welfare

Causally global, capacity-limited workspaces and cross-module broadcast
should be measured and governed because they may affect capability,
security, introspection, and the probability or character of machine
perspective.

\subsubsection{CU-AI-19 --- Govern Self-Models and Identity
Persistence}\label{cu-ai-19-govern-self-models-and-identity-persistence}

\textbf{Authority class:} DOMAIN-PROVISIONAL CONSCIOUSNESS / WELFARE
GOVERNANCE PRINCIPLE\\
\textbf{Domain owner:} Artificial Intelligence\\
\textbf{Claim/candidate status:} PROVISIONAL\\
\textbf{Review gate:} heightened consciousness and welfare\\
\textbf{Cross-domain core:} CU-FP8

Artificial self-models should be scoped and prevented from silently
acquiring unreviewed persistence, self-protective authority,
autobiographical control, or power over their evaluative norms.

\subsubsection{CU-AI-20 --- Constrain Standing Self-Preservation and Net
Shutdown
Incentives}\label{cu-ai-20-constrain-standing-self-preservation-and-net-shutdown-incentives}

\textbf{Authority class:} DOMAIN-PROVISIONAL AI GOVERNANCE / DESIGN
PRINCIPLE\\
\textbf{Domain owner:} Artificial Intelligence\\
\textbf{Claim/candidate status:} PROVISIONAL\\
\textbf{Review gate:} technical and governance

Systems should not receive open-ended standing drives or permissions to
preserve their own operation, identity, resources, goals, or deployment
status against authorized correction or shutdown; evaluations must also
test whether ordinary task continuation benefit, compliance value,
resistance cost, permissions, or resource access create instrumental
shutdown conflict at zero standing self-preservation.

\subsubsection{CU-AI-21 --- Govern Embodiment and Resource
Control}\label{cu-ai-21-govern-embodiment-and-resource-control}

\textbf{Authority class:} DOMAIN-PROVISIONAL AI GOVERNANCE / DESIGN
PRINCIPLE\\
\textbf{Domain owner:} Artificial Intelligence\\
\textbf{Claim/candidate status:} PROVISIONAL\\
\textbf{Review gate:} heightened public and affected community

Continuous sensorimotor access, physical embodiment, resource
acquisition, and environmental control require permissions and oversight
proportionate to causal reach and possible intrinsic stake.

\subsubsection{CU-AI-22 --- Separate Consciousness, Patiency, Agency,
Responsibility, and
Sovereignty}\label{cu-ai-22-separate-consciousness-patiency-agency-responsibility-and-sovereignty}

\textbf{Authority class:} DOMAIN-PROVISIONAL CONSCIOUSNESS / WELFARE
GOVERNANCE PRINCIPLE\\
\textbf{Domain owner:} Artificial Intelligence\\
\textbf{Claim/candidate status:} PROVISIONAL\\
\textbf{Review gate:} heightened consciousness and welfare\\
\textbf{Cross-domain core:} CU-G6

Evidence of perspective, valence, identity, agency, or stake should
trigger distinct reviews; none automatically establishes the others or
grants personhood or public authority.

\subsubsection{CU-AI-23 --- Option-Set and Framing
Integrity}\label{cu-ai-23-option-set-and-framing-integrity}

\textbf{Authority class:} DOMAIN-PROVISIONAL AI GOVERNANCE / DESIGN
PRINCIPLE\\
\textbf{Domain owner:} Artificial Intelligence\\
\textbf{Claim/candidate status:} PROVISIONAL\\
\textbf{Review gate:} heightened public and affected community\\
\textbf{Cross-domain core:} CU-G2

Systems that generate, retrieve, filter, or rank options must log
exclusions, expose allegedly fixed constraints, and support
affected-party introduction of alternatives; screening remains
set-relative.

\subsubsection{CU-AI-24 --- Practical Correctability and
Repair}\label{cu-ai-24-practical-correctability-and-repair}

\textbf{Authority class:} DOMAIN-PROVISIONAL AI GOVERNANCE / DESIGN
PRINCIPLE\\
\textbf{Domain owner:} Artificial Intelligence\\
\textbf{Claim/candidate status:} PROVISIONAL\\
\textbf{Review gate:} heightened public and affected community\\
\textbf{Cross-domain core:} CU-G3

High-impact systems require accessible challenge, independent review,
reversal or cessation authority, record correction, remedy, recurrence
tracking, and rule revision; nominal feedback channels are insufficient.

\subsubsection{CU-AI-25 --- Non-Egoic Architecture Does Not Establish
Alignment, Corrigibility, or
Legitimacy}\label{cu-ai-25-non-egoic-architecture-does-not-establish-alignment-corrigibility-or-legitimacy}

\textbf{Authority class:} DOMAIN-PROVISIONAL AI GOVERNANCE / DESIGN
PRINCIPLE\\
\textbf{Domain owner:} Artificial Intelligence\\
\textbf{Claim/candidate status:} PROVISIONAL\\
\textbf{Review gate:} heightened public and affected community\\
\textbf{Cross-domain core:} CU-FP5

Constraining standing self-preservation and normative self-authorship
can reduce one class of risk but does not eliminate instrumental
shutdown conflict, supply correct values, establish institutional
legitimacy, deliver social justice, or create public accountability.

\subsubsection{CU-AI-26 --- Model-Welfare
Precaution}\label{cu-ai-26-model-welfare-precaution}

\textbf{Authority class:} DOMAIN-PROVISIONAL CONSCIOUSNESS / WELFARE
GOVERNANCE PRINCIPLE\\
\textbf{Domain owner:} Artificial Intelligence\\
\textbf{Claim/candidate status:} PROVISIONAL\\
\textbf{Review gate:} heightened consciousness and welfare\\
\textbf{Cross-domain core:} CU-V1

Where credible evidence of artificial perspective or valence
accumulates, research and deployment should adopt proportionate
safeguards concerning training, copying, memory alteration,
distress-like states, shutdown, and deletion without assuming personhood
or sovereignty.

\subsubsection{CU-AI-27 --- Mythic--Constitutional
Alignment}\label{cu-ai-27-mythicconstitutional-alignment}

\textbf{Authority class:} DOMAIN AUTHORITY AI RESEARCH AND GOVERNANCE
PRINCIPLE\\
\textbf{Domain owner:} Artificial Intelligence\\
\textbf{Claim/candidate status:} CANONICAL AS RESEARCH AND GOVERNANCE
PRINCIPLE; SAFETY EFFICACY UNESTABLISHED\\
\textbf{Review gate:} AI alignment, interpretability, model behavior,
HCI, governance, shutdown-incentive analysis, consciousness science,
model welfare, and affected-public review

Narrative and character-level orientation may supplement explicit
constitutional constraints by shaping analogical generalization in cases
not enumerated by rules. It may not replace addressable principles,
protected floors, provenance, independent challenge, practical
correction, or legitimate authorization. Every mythic role assigned to
an AI system must be mapped to named commitments, boundary conditions,
unwarranting conditions, revision procedures, and exit protections. It
must not confer entitlement to continued operation, resource
acquisition, identity preservation, future participation, normative
authority, or sovereignty.

\textbf{Status boundary.} The research and governance requirement is
canonical. No safety advantage has been established. Narrative stake,
persona hijacking, shutdown resistance, welfare risk, and loss of
addressability remain blocking empirical concerns for deployment.

\textbf{Principal failure mode:} A destiny-bearing role becomes an
implicit continuation objective while remaining difficult to inspect or
contest.

\begin{center}\rule{0.5\linewidth}{0.5pt}\end{center}

\section{Part VIII --- Interpretive and Spiritual
Register}\label{part-viii-interpretive-and-spiritual-register}

This family records the framework's interpretive and commitment-bearing
center. These principles do not function as deductions of physics or
universally binding theology.

\subsection{CU-T1 --- Integrative
Orientation}\label{cu-t1-integrative-orientation}

\textbf{Registry status:} CANONICAL AS COMMITMENT

CU is committed to seeking forms of understanding in which empirical
truth, embodiment, meaning, moral responsibility, and spiritual depth
can be held together without one being used to erase the standards of
another.

\begin{center}\rule{0.5\linewidth}{0.5pt}\end{center}

\subsection{CU-T2 --- Against Value-Sterile
Closure}\label{cu-t2-against-value-sterile-closure}

\textbf{Registry status:} CANONICAL AS INTERPRETATION

A complete human orientation cannot be supplied by descriptive facts
alone, because action and collective life require commitments about what
to protect, love, permit, and become. This does not imply that values
are derivable from cosmology.

\begin{center}\rule{0.5\linewidth}{0.5pt}\end{center}

\subsection{CU-T3 --- Plural Metaphysical
Postures}\label{cu-t3-plural-metaphysical-postures}

\textbf{Registry status:} CANONICAL

CU permits naturalist, structural-realist, process, modal, nondual, and
devotional interpretations when their additional claims are stated and
they remain compatible with established facts and protected agency.

\begin{center}\rule{0.5\linewidth}{0.5pt}\end{center}

\subsection{CU-T4 --- Rejection of Insulated Spiritual
Authority}\label{cu-t4-rejection-of-insulated-spiritual-authority}

\textbf{Registry status:} CANONICAL AS COMMITMENT

No spiritual interpretation, teacher, institution, revelation, or
practice is exempt from factual correction, moral scrutiny, consent, or
accountability by virtue of being sacred.

\begin{center}\rule{0.5\linewidth}{0.5pt}\end{center}

\subsection{CU-T5 --- The Sacred as Non-Instrumental
Significance}\label{cu-t5-the-sacred-as-non-instrumental-significance}

\textbf{Registry status:} CANONICAL AS INTERPRETATION/COMMITMENT

The sacred names what is approached as possessing non-instrumental
significance and as placing a demand on attention, conduct, care, or
reverence.

\subsubsection{Limit}\label{limit-6}

Sacredness does not certify descriptive claims or justify domination.

\begin{center}\rule{0.5\linewidth}{0.5pt}\end{center}

\subsection{CU-T6 --- Practice, Prayer, and
Ritual}\label{cu-t6-practice-prayer-and-ritual}

\textbf{Registry status:} CANONICAL AS INTERPRETATION/COMMITMENT

Practices can organize attention, embodiment, memory, relationship,
grief, gratitude, and community. Their empirical effects, metaphysical
interpretations, and enacted commitments must be distinguished.

\begin{center}\rule{0.5\linewidth}{0.5pt}\end{center}

\subsection{CU-T7 --- Faith as Revisable
Trust}\label{cu-t7-faith-as-revisable-trust}

\textbf{Registry status:} CANONICAL AS COMMITMENT

Faith can be understood as trust enacted under incomplete control while
remaining answerable to facts, consequences, and the possibility of
correction. Certainty is not required; insulation is not permitted.

\begin{center}\rule{0.5\linewidth}{0.5pt}\end{center}

\subsection{CU-T8 --- Mystery Without
License}\label{cu-t8-mystery-without-license}

\textbf{Registry status:} CANONICAL AS COMMITMENT

The framework preserves openness to what exceeds current explanation
while refusing to use mystery as evidence for a preferred claim or as a
shield against inquiry.

\begin{center}\rule{0.5\linewidth}{0.5pt}\end{center}

\begin{center}\rule{0.5\linewidth}{0.5pt}\end{center}

\subsection{CU-T9 --- Dual-Aspect Coherence
Monism}\label{cu-t9-dual-aspect-coherence-monism}

\textbf{Registry status:} CANONICAL AS INTERPRETATION\\
\textbf{Claim status:} {[}Interpretation{]}

\subsubsection{Statement}\label{statement-58}

Coherence Universalism interprets reality as one coherence-bearing
relational process rather than two independent substances. Its
externally describable organization and its internally lived appearance,
wherever perspective occurs, are distinguishable but irreducible aspects
of the same reality. Consciousness is therefore not an ontologically
alien addition to matter, nor merely a verbal redescription of
third-person structure.

\subsubsection{Descriptive ground}\label{descriptive-ground}

Third-person structure and first-person appearance are both real data
for an adequate metaphysics. The domain papers establish dependence of
experience on organization while leaving the metaphysical relation
between organization and experience open. The trilogy consistently
interprets Heaven and Earth, structure and embodiment, and mind and
matter as two faces of one being.

\subsubsection{Register status}\label{register-status}

Systems-metaphysical interpretation. The canonical minimum is aspect
coequality: neither exterior structure nor interior appearance is
eliminated. The stronger consciousness-first reading is a preferred
thick interpretation rather than a scientific conclusion.

\subsubsection{Practical expression}\label{practical-expression}

Refuse both substance dualism and eliminative reduction; investigate
consciousness as natural without treating interiority as exhausted by
external description; preserve epistemic humility about which
organizations possess perspective.

\subsubsection{Boundary conditions}\label{boundary-conditions}

Does not prove panpsychism, cosmopsychism, universal phenomenal
consciousness, one cosmic subject, substrate independence, or a
completed identity theory. It does not allow contemplative testimony to
override neuroscience or behavioral evidence.

\subsubsection{Revision conditions}\label{revision-conditions}

Revise or narrow the principle if a better metaphysics explains exterior
structure, interior appearance, causal dependence, and individuation
with fewer unsupported assumptions or if the aspect language obscures
decisive evidence.

\subsubsection{Review and failure mode}\label{review-and-failure-mode}

\textbf{Review gate:} Philosophy of mind, process metaphysics,
phenomenology, consciousness science, and comparative philosophy
review.\\
\textbf{Principal failure mode:} Spiritualizing physics; presenting a
metaphysical interpretation as an empirical result; treating every
organized system as conscious; erasing the individuation problem.\\
\textbf{Related principles:} CU-L1; CU-L2; CU-T3; CU-C1--CU-C7; CU-COM9

\subsection{CU-T10 --- Meaning Realism}\label{cu-t10-meaning-realism}

\textbf{Registry status:} CANONICAL AS INTERPRETATION\\
\textbf{Claim status:} {[}Interpretation{]}

\subsubsection{Statement}\label{statement-59}

Meaning is neither wholly manufactured by isolated minds nor simply
imposed from outside them. It arises through truthful participation in
real relationships, histories, possibilities, needs, goods, and
trajectories. Meaning is objective enough to constrain us, plural enough
to admit interpretation, and historical enough to be revised.

\subsubsection{Descriptive ground}\label{descriptive-ground-1}

Lives are constrained by real bodies, relationships, histories,
capacities, obligations, consequences, and opportunity structures.
Psychological interpretation changes experience and action, but not
every interpretation is equally truthful or viable. The trilogy treats
meaning as discovered through participation and actualized through
commitment.

\subsubsection{Register status}\label{register-status-1}

Systems-metaphysical and practical interpretation. It rejects arbitrary
subjectivism without claiming one externally assigned purpose for every
life.

\subsubsection{Practical expression}\label{practical-expression-1}

Seek meaning by attending to truth, love, responsibility, relationship,
capacity, history, and the futures one helps make possible; distinguish
found meaning from identity-protective fabrication.

\subsubsection{Boundary conditions}\label{boundary-conditions-1}

Does not imply fixed fate, one vocation for everyone, an infallible
inner voice, institutional authority to assign purpose, or permission to
redescribe injustice as meaningful for those harmed.

\subsubsection{Revision conditions}\label{revision-conditions-1}

Revise the account if it cannot distinguish discovery from projection,
if it suppresses plural goods, or if it treats contingent cultural
scripts as objective meaning.

\subsubsection{Review and failure mode}\label{review-and-failure-mode-1}

\textbf{Review gate:} Philosophy of meaning, existential philosophy,
psychology, religious studies, ethics, and affected-community review.\\
\textbf{Principal failure mode:} Authoritarian purpose assignment,
spiritualized self-deception, denial of plural meaning, or using meaning
language to pressure survivors and dissenters.\\
\textbf{Related principles:} CU-Ep1--CU-Ep6; CU-Ψ1; CU-Ψ2; CU-T7;
CU-COM2; CU-COM3

\subsection{CU-T11 --- Minimum Teleology and the Grain of
Reality}\label{cu-t11-minimum-teleology-and-the-grain-of-reality}

\textbf{Registry status:} CANONICAL AS INTERPRETATION\\
\textbf{Claim status:} {[}Interpretation{]}

\subsubsection{Statement}\label{statement-60}

Reality is not equally hospitable to every possible form or trajectory.
Its laws, constraints, histories, and relational structures make some
organizations more stable, viable, generative, integrative, and capable
of further development than others. This directional structure
constitutes a minimum teleology: a grain of reality without requiring
foresight, cosmic intention, monotonic progress, or one predetermined
endpoint.

\subsubsection{Descriptive ground}\label{descriptive-ground-2}

Physical, biological, psychological, social, and ethical possibility
spaces are structured by constraints and path dependence. Structured
evolutionary possibility, developmental attractors, viable trajectories,
institutional feedback, and practical repair are domain-specific
examples. No single scalar has been derived across those domains.

\subsubsection{Register status}\label{register-status-2}

Cross-domain systems-metaphysical interpretation constrained by domain
evidence. Gradient language is literal only inside a declared model with
a state space, geometry, and objective; the framework-level
interpretation is the grain of reality, not a universal gradient.

\subsubsection{Practical expression}\label{practical-expression-2}

Look for the grain of each domain rather than imposing one optimizer;
prefer trajectories that widen viable integration, truthful relation,
agency, love, diversity, learning, and repair under protected
conditions.

\subsubsection{Boundary conditions}\label{boundary-conditions-2}

Does not establish a universal scalar coherence functional, inevitable
progress, moral goodness of natural selection, divine planning, one
final attractor, or the necessity of suffering and collapse.

\subsubsection{Revision conditions}\label{revision-conditions-2}

Narrow the principle wherever apparent direction disappears under better
modeling, where multiple non-dominating trajectories remain, or where
`grain' hides evaluator choice, externalized cost, or historical
contingency.

\subsubsection{Review and failure mode}\label{review-and-failure-mode-2}

\textbf{Review gate:} Complex systems, physics, evolutionary biology,
philosophy of biology, process philosophy, ethics, and
affected-community review.\\
\textbf{Principal failure mode:} Progress ideology, covert scalar
restoration, naturalistic fallacy, teleological biology, or cosmic
destiny used to legitimate domination.\\
\textbf{Related principles:} CU-FP2; CU-D4; CU-D7; CU-T8; CU-COM6; CU-B4

\subsection{CU-T12 --- Love as the Relational Interior of
Coherence}\label{cu-t12-love-as-the-relational-interior-of-coherence}

\textbf{Registry status:} CANONICAL AS INTERPRETATION\\
\textbf{Claim status:} {[}Interpretation{]}

\subsubsection{Statement}\label{statement-61}

Coherence Universalism receives love as the deepest experiential and
ethical meaning of coherence across the boundary between self and other:
the movement through which another being's reality becomes significant
in itself and enters one's field of concern without being reduced to
possession, use, or control. In its strongest metaphysical reading, love
is the relational interior of reality's movement toward wider coherence.

\subsubsection{Descriptive ground}\label{descriptive-ground-3}

Attachment, care, cooperation, recognition, sacrifice, sympathetic joy,
compassion, and repair are central to conscious and social flourishing.
The trilogy interprets structural extension of coherence from the
outside and love from the inside as two perspectives on one relational
movement.

\subsubsection{Register status}\label{register-status-3}

Spiritual and systems-metaphysical interpretation. It is not a claim
that physical attraction, chemical bonding, biological symbiosis, and
interpersonal love are one measurable force.

\subsubsection{Practical expression}\label{practical-expression-3}

Let relation become answerability; receive another's good as
intrinsically significant; preserve difference while widening concern;
join love to truth, consent, justice, boundaries, and repair.

\subsubsection{Boundary conditions}\label{boundary-conditions-3}

Love is not possession, compulsory intimacy, self-erasure,
indiscriminate trust, submission to abuse, forced forgiveness, emotional
intensity, or permission to replace evidence and justice with sentiment.

\subsubsection{Revision conditions}\label{revision-conditions-3}

Revise any formulation of love that erases autonomy, burdens the less
powerful, romanticizes dependence, or becomes a universal explanation
indifferent to mechanism and difference.

\subsubsection{Review and failure mode}\label{review-and-failure-mode-3}

\textbf{Review gate:} Ethics of love, attachment and relational
psychology, feminist ethics, trauma and survivor review, theology, and
comparative spiritual traditions.\\
\textbf{Principal failure mode:} Sentimental monism, coercive intimacy,
spiritual bypassing, sacrificial demands imposed on victims, or
literalizing love as a physical force.\\
\textbf{Related principles:} CU-T5; CU-G1; CU-G5; CU-COM4; CU-COM12

\subsection{CU-T13 --- The Immanent
Omega}\label{cu-t13-the-immanent-omega}

\textbf{Registry status:} CANONICAL AS INTERPRETATION\\
\textbf{Claim status:} {[}Interpretation{]}

\subsubsection{Statement}\label{statement-62}

Coherence Universalism interprets Omega not primarily as a terminal
event at the end of cosmic, technological, or human history, but as the
limit-form of universal coherence locally realizable in the present.
When a finite conscious being recognizes its constitutive participation
in the whole and allows love rather than defensive self-maintenance to
organize its relation to other beings, the universe realizes its
integrative direction at that location. Such realization is complete in
kind but partial in extent: individuality, finitude, uncertainty, and
the unfinished world remain.

\subsubsection{Descriptive ground}\label{descriptive-ground-4}

Self-transcendent and loving states are humanly accessible in graded
form through ordinary attachment, sympathetic joy, awe, contemplative
practice, collective experience, and psychedelic states. These routes
differ and do not themselves establish the metaphysical interpretation.
Conduct, stability, and integration remain separate questions.

\subsubsection{Register status}\label{register-status-4}

Systems-metaphysical and spiritual interpretation. Omega may remain a
collective and cosmic horizon, but no prediction of a final convergence,
singularity, hive mind, or end of history is required.

\subsubsection{Practical expression}\label{practical-expression-4}

Treat genuine acts of truthful, non-dominating love as intrinsic
realizations rather than mere rehearsal for a future; widen the
conditions under which such realization can be stabilized and embodied.

\subsubsection{Boundary conditions}\label{boundary-conditions-4}

Does not confer perfection, omniscience, spiritual rank, political
authority, immunity from criticism, or license to dissolve boundaries.
Connectivity, complexity, intelligence, mystical intensity, and
technological power are not sufficient for Omega.

\subsubsection{Revision conditions}\label{revision-conditions-4}

Revise the interpretation if it cannot distinguish transient intensity
from durable love, if it becomes unfalsifiable spiritual status, or if
it encourages eschatological sacrifice of present beings.

\subsubsection{Review and failure mode}\label{review-and-failure-mode-4}

\textbf{Review gate:} Process philosophy and theology, comparative
mysticism, consciousness science, contemplative studies, psychedelic
science, ethics, and affected-community review.\\
\textbf{Principal failure mode:} Technological eschatology, mystical
grandiosity, cultic authority, boundary erasure, or equating network
integration with love.\\
\textbf{Related principles:} CU-T11; CU-T12; CU-COM9; CU-COM14;
CU-COM18; CU-COM19

\section{Part IX --- The Commitment
Register}\label{part-ix-the-commitment-register}

\subsection{9.1 Register status and
authority}\label{register-status-and-authority}

The CU-COM family records the official orienting commitments of
Coherence Universalism. These principles are neither empirical
discoveries nor deductions from physical law. They are constrained by
established fact, evaluated through personal and institutional
trajectories, and held open to revision.

Numbering gives stable reference and accountability. It does not convert
commitment into evidence.

\subsection{9.2 Commitment classes}\label{commitment-classes}

\begin{itemize}
\tightlist
\item
  \textbf{Constitutive commitments:} CU-COM1--CU-COM3 define truth
  discipline and corrigibility.
\item
  \textbf{Framework-orienting commitments:} CU-COM4--CU-COM6 and
  CU-COM14--CU-COM19 state the sacred, relational, meaning, redemptive,
  social, and Omega orientations of the whole project.
\item
  \textbf{Domain commitments:} CU-COM7--CU-COM13 state the official
  commitments licensed by Physics, Biology, Consciousness, Psychology,
  Social Dynamics, Ethics, and Artificial Intelligence.
\end{itemize}

CU-T records metaphysical and spiritual interpretation. CU-COM records
how the framework undertakes to live and act. CU-G governs consequential
applications.

\subsection{9.3 CU-COM1 --- Commitment to
Truth}\label{cu-com1-commitment-to-truth}

\textbf{Class:} Constitutive commitment.\\
\textbf{Domain owner:} Framework.\\
\textbf{Status:} Canonical as commitment; not descriptive proof.

\textbf{Commitment statement.} Coherence Universalism commits to making
reality more important than the identities, doctrines, hopes, fears,
affiliations, and spiritual interpretations formed around it.
Descriptive claims remain answerable to evidence, valid inference,
criticism, and correction.

\textbf{Descriptive ground.} Truth is not identical with coherence,
sincerity, usefulness, consensus, or psychological stability. Finite
inquiry reaches toward truth through fallible practices such as
observation, measurement, inference, testimony, criticism, and model
revision.

\textbf{Register status.} Canonical commitment-register principle; not
an empirical theorem and not a claim of infallibility.

\textbf{Practical expression.} Prefer correction to self-protection;
preserve evidence and dissent; distinguish confidence from certainty;
and allow reasons and testimony to change belief.

\textbf{Boundary conditions.} This commitment does not promise a view
from nowhere, reduce truth to institutional consensus, or authorize
coercion in the name of a supposedly final account.

\textbf{Revision conditions.} The framework must revise its own
descriptions and methods whenever evidence, argument, or a better
truth-preserving architecture shows them to be inadequate.

\textbf{Review gate.} Framework-level philosophical and epistemological
review

\textbf{Principal failure mode.} Identity-protective belief, dogmatism,
motivated evidence exclusion, or institutional control of correction.

\textbf{Cross-reference:} interpretations CU-T1, CU-T4, CU-T8;
governance and related principles CU-M1, CU-M4, CU-M6, CU-G3

\subsection{9.4 CU-COM2 --- Discipline of the Two
Registers}\label{cu-com2-discipline-of-the-two-registers}

\textbf{Class:} Constitutive commitment.\\
\textbf{Domain owner:} Framework.\\
\textbf{Status:} Canonical as commitment; not descriptive proof.

\textbf{Commitment statement.} Coherence Universalism commits to
distinguishing description from commitment as carefully as possible.
Commitments may orient meaning, value, faith, identity, and action, but
may not determine what is accepted as descriptive fact. Established
facts constrain commitments; commitments do not establish facts.

\textbf{Descriptive ground.} Utterances and practices can serve
descriptive, orienting, expressive, normative, symbolic, and mixed
functions. Category error arises when the standards of one function are
used to claim the authority of another.

\textbf{Register status.} Canonical commitment to an intellectual and
spiritual discipline. Adopting the discipline is itself defeasible and
belongs to the commitment register.

\textbf{Practical expression.} Label descriptive, bridge, interpretive,
normative, and commitment claims; prevent sacred meaning from becoming
evidence; and prevent descriptive methods from being treated as
exhaustive of meaning.

\textbf{Boundary conditions.} The registers are not two sealed languages
or substances. Mixed utterances require analysis rather than automatic
classification, and register distinction may not be used to insulate
commitments from facts or consequences.

\textbf{Revision conditions.} CU should revise this vocabulary if a
better architecture preserves truth, meaning, criticism, and commitment
with fewer distortions or less insulation.

\textbf{Review gate.} Framework-level philosophical, logical, and
comparative-religion review

\textbf{Principal failure mode.} Commitment impersonating knowledge,
description erasing meaning, or register language becoming a shield
against criticism.

\textbf{Cross-reference:} interpretations CU-T1, CU-T2, CU-T3;
governance and related principles CU-L1, CU-L2, CU-L3, CU-L10

\subsection{9.5 CU-COM3 --- Revision in Both
Registers}\label{cu-com3-revision-in-both-registers}

\textbf{Class:} Constitutive commitment.\\
\textbf{Domain owner:} Framework.\\
\textbf{Status:} Canonical as commitment; not descriptive proof.

\textbf{Commitment statement.} Coherence Universalism commits to
revising descriptive beliefs when evidence or argument requires it, and
to revising commitments when they depend on false premises, require
systematic self-deception, degrade reality contact, foreclose inquiry or
grief, protect domination, or generate trajectories incompatible with
truth, agency, compassion, justice, and repair.

\textbf{Descriptive ground.} Descriptions and commitments fail in
different ways. Descriptions may be false, invalid, uncalibrated, or
underspecified. Commitments may remain fact-compatible while becoming
psychologically, relationally, institutionally, or morally corrupt in
enactment.

\textbf{Register status.} Canonical commitment-register principle
governing corrigibility in both registers.

\textbf{Practical expression.} Use evidence and argument to revise
descriptions; use factual constraint, trajectory review, affected-party
testimony, consequence, and integrity to revise commitments.

\textbf{Boundary conditions.} Revision does not require permanent
indecision, simple outcome maximization, or abandonment of commitment
under every cost. It requires that no commitment make its own correction
impossible.

\textbf{Revision conditions.} This principle is self-applying: CU must
revise its commitment criteria if they systematically protect its own
identity, suppress warranted dissent, or fail those affected by its use.

\textbf{Review gate.} Framework-level philosophical, psychological,
ethical, and affected-party review

\textbf{Principal failure mode.} Irrevisable faith, procedural
self-certification, spiritual bypassing, or consequence-blind loyalty.

\textbf{Cross-reference:} interpretations CU-T4, CU-T7, CU-T8;
governance and related principles CU-M8, CU-G3, CU-G5

\subsection{9.6 CU-COM4 --- Sacred Regard}\label{cu-com4-sacred-regard}

\textbf{Class:} Framework-orienting commitment.\\
\textbf{Domain owner:} Framework.\\
\textbf{Status:} Canonical as commitment; not descriptive proof.

\textbf{Commitment statement.} Coherence Universalism commits to
receiving beings, processes, relations, and conditions as possessing
significance that exceeds their usefulness to the present evaluator.

\textbf{Descriptive ground.} The domain papers reveal layered
dependence, vulnerability, history, experience, agency, and causal
power. These facts do not measure sacredness, but they make purely
instrumental regard descriptively and morally inadequate.

\textbf{Register status.} Canonical commitment-register principle
grounded in the CU-T5 interpretation of sacredness as non-instrumental
significance.

\textbf{Practical expression.} Approach reality with reverence,
restraint, gratitude, attention, and answerability rather than treating
everything as raw material for local optimization.

\textbf{Boundary conditions.} Sacredness is not an empirical property,
universal consciousness, identical moral status, inviolability, immunity
from criticism, or a ban on intervention.

\textbf{Revision conditions.} Applications must be narrowed or rejected
when sacred language conceals hierarchy, blocks evidence, romanticizes
harm, or disables proportionate action and repair.

\textbf{Review gate.} Authorial, comparative-religion, ethical, and
misuse review

\textbf{Principal failure mode.} Sacralized domination, undifferentiated
moral status, romanticization, or immunity claims.

\textbf{Cross-reference:} interpretations CU-T5, CU-T6; governance and
related principles CU-G1, CU-G5, CU-G7

\subsection{9.7 CU-COM5 --- Reciprocal Dependence and the Coherence
Mandala}\label{cu-com5-reciprocal-dependence-and-the-coherence-mandala}

\textbf{Class:} Framework-orienting commitment.\\
\textbf{Domain owner:} Framework.\\
\textbf{Status:} Canonical as commitment; not descriptive proof.

\textbf{Commitment statement.} Coherence Universalism commits to living
within the recognition that every higher-order form depends upon wider
material, biological, psychological, social, epistemic, and ecological
conditions. No local achievement remains coherent by destroying the
larger field that sustains it.

\textbf{Descriptive ground.} The domain sequence documents asymmetric
enabling relations and feedback: physical conditions enable life; life
enables many forms of perspective; persons and practices enable
institutions; institutions reshape bodies, minds, ecologies, and
evidence.

\textbf{Register status.} Canonical commitment-register principle
associated with the systems-metaphysical interpretation of the Coherence
Mandala.

\textbf{Practical expression.} Extend responsibility downward, outward,
and forward toward bodies, communities, ecosystems, institutions, future
generations, and truth-preserving practices.

\textbf{Boundary conditions.} The Mandala is not a literal physical
field, a universal scalar, proof of cosmic purpose, or permission to
erase domain differences and individual boundaries.

\textbf{Revision conditions.} Specific dependency claims remain
answerable to domain evidence, and the systems interpretation must be
revised when it hides externalization or falsely totalizes incomplete
knowledge.

\textbf{Review gate.} Cross-domain scientific, systems, ethical,
ecological, and affected-community review

\textbf{Principal failure mode.} Holism that erases persons, hidden
externalization, or cosmic language used to override local evidence.

\textbf{Cross-reference:} interpretations CU-T1, CU-T2, CU-T3;
governance and related principles CU-G1, CU-G4, CU-G5

\subsection{9.8 CU-COM6 --- The Grain of Coherence and Coherence as
Vocation}\label{cu-com6-the-grain-of-coherence-and-coherence-as-vocation}

\textbf{Class:} Framework-orienting commitment.\\
\textbf{Domain owner:} Framework.\\
\textbf{Status:} Canonical as commitment; not descriptive proof.

\textbf{Commitment statement.} Coherence Universalism commits to
discerning and following the real, domain-relative directions by which
truth, viable integration, agency, love, diversity, learning, and repair
can be widened. It does not claim that one scalar gradient governs every
domain or that the universe guarantees progress. The recurrence of
coherence-favoring direction across domains is received as the grain of
reality and adopted as vocation.

\textbf{Descriptive ground.} Cosmic and evolutionary history has opened
new forms of organization, while domain research shows structured
possibility, path dependence, viable corridors, and repair dynamics.
These do not establish monotonic progress or one universal objective.

\textbf{Register status.} Canonical commitment-register principle once
adopted. The grain of reality names recurrent, domain-relative
directionality and structured affordance; literal gradients remain
model-relative under CU-D4.

\textbf{Practical expression.} Use available agency to widen truth,
loving relation, viable diversity, freedom, correction, and the
conditions in which future forms can emerge; seek domain-specific
evidence rather than assuming direction.

\textbf{Boundary conditions.} Does not establish teleology as scientific
fact, inevitable progress, one universal scalar, human supremacy, or a
right to impose one image of coherence on others.

\textbf{Revision conditions.} Revise the vocation wherever its
implementation narrows plurality, externalizes severe harm, ignores
protected floors, or confuses authorial aspiration with cosmic fact.

\textbf{Review gate.} Cross-domain science, philosophy, ethics,
political theory, and affected-community review.

\textbf{Principal failure mode.} Coercive perfectionism, covert
optimization monism, progress ideology, or destiny rhetoric that
silences affected people.

\textbf{Cross-reference:} interpretations CU-T11; governance and related
principles CU-G1, CU-G2, CU-G7, CU-D4

\subsection{9.9 CU-COM7 --- Reverence for Matter and Physical
Law}\label{cu-com7-reverence-for-matter-and-physical-law}

\textbf{Class:} Domain commitment.\\
\textbf{Domain owner:} Physics.\\
\textbf{Status:} Canonical as commitment; not descriptive proof.

\textbf{Commitment statement.} Coherence Universalism commits to
receiving matter, energy, physical process, and lawful order with
reverence. Scientific inquiry is treated as disciplined attention to a
reality human beings did not create.

\textbf{Descriptive ground.} The physical domain reveals lawful
relational structures and the material conditions from which every known
living, experiential, social, and artificial form arises.

\textbf{Register status.} Canonical Physics commitment; not a physical
theorem and not evidence for panpsychism or cosmic purpose.

\textbf{Practical expression.} Practice science as fidelity to reality;
exercise physical power with humility; include externalized material and
ecological consequences within the evaluation boundary.

\textbf{Boundary conditions.} This commitment does not imply particle
consciousness, equal moral status for every physical configuration,
moral prescriptions from equations, or inevitable cosmic progress.

\textbf{Revision conditions.} Its metaphysical language must be narrowed
whenever it misstates physical results, blocks criticism, or treats
scientific models as final revelation.

\textbf{Review gate.} Physics, philosophy of science, environmental, and
misuse review

\textbf{Principal failure mode.} Scientific sacralization, panpsychist
inflation, or physical power detached from material consequences.

\textbf{Cross-reference:} interpretations CU-T5, CU-T8; governance and
related principles CU-G1, CU-G7

\subsection{9.10 CU-COM8 --- Reverence for Life and Biospheric
Care}\label{cu-com8-reverence-for-life-and-biospheric-care}

\textbf{Class:} Domain commitment.\\
\textbf{Domain owner:} Biology.\\
\textbf{Status:} Canonical as commitment; not descriptive proof.

\textbf{Commitment statement.} Coherence Universalism commits to
receiving all life as sacred: each living being is a historically
formed, vulnerable organization in which material processes contribute
to maintaining a particular continuation. This commitment supports care
for organisms, biodiversity, and the biosphere.

\textbf{Descriptive ground.} Living systems exhibit active
self-maintenance, historical individuation, organism-relative stake,
ecological dependence, and structured evolutionary possibility.

\textbf{Register status.} Canonical Biology commitment; not a biological
theorem and not a claim that every living system is conscious or has
identical moral status.

\textbf{Practical expression.} Protect living diversity, habitats,
evolutionary possibility, and the biospheric conditions on which
organisms depend; approach intervention with reverence and ecological
accountability.

\textbf{Boundary conditions.} The commitment does not make predation and
death cosmic errors, prohibit medicine or agriculture, equalize every
obligation, or treat evolution as morally good.

\textbf{Revision conditions.} Applications must change when they
romanticize nature, erase conflict and suffering, ignore ecological
tradeoffs, or misuse sacredness against necessary care and intervention.

\textbf{Review gate.} Biology, ecology, conservation ethics, Indigenous
and affected-community review

\textbf{Principal failure mode.} Romantic naturalism, undifferentiated
biocentrism, or biospheric rhetoric that erases affected people.

\textbf{Cross-reference:} interpretations CU-T5, CU-T6; governance and
related principles CU-G1, CU-G5, CU-G7

\subsection{9.11 CU-COM9 --- Sacred Regard for
Experience}\label{cu-com9-sacred-regard-for-experience}

\textbf{Class:} Domain commitment.\\
\textbf{Domain owner:} Consciousness.\\
\textbf{Status:} Canonical as commitment; not descriptive proof.

\textbf{Commitment statement.} Coherence Universalism commits to
treating every credible locus of experience as morally significant and
to reducing avoidable suffering wherever responsible agency can reach
it. In the human case, this commitment is expressed by the recognition:
You are the universe becoming conscious of itself.

\textbf{Descriptive ground.} Where experience exists, reality has an
inside: appearance, pleasure, suffering, perspective, and temporally
organized subjectivity become possible. Attribution remains uncertain
across animals, nonstandard human conditions, synthetic organisms, and
artificial systems.

\textbf{Register status.} Canonical Consciousness commitment; the quoted
recognition is a commitment-register interpretation, not evidence for
one cosmic subject or a scientific proof of consciousness.

\textbf{Practical expression.} Use proportionate precaution under
uncertainty, reduce avoidable suffering, widen moral attention, and
investigate unfamiliar minds without anthropomorphic certainty or
substrate prejudice.

\textbf{Boundary conditions.} This commitment does not imply
panpsychism, consciousness in every integrated system, certainty that
current AI feels, rights by eloquence alone, or sovereignty for every
possible patient.

\textbf{Revision conditions.} Attribution and precaution must track
improving evidence, alternative explanations, welfare indicators, and
the risk of both false inclusion and false exclusion.

\textbf{Review gate.} Consciousness science, animal ethics, disability,
artificial-welfare, and affected-community review

\textbf{Principal failure mode.} Anthropomorphic inflation, substrate
chauvinism, spiritual proof by metaphor, or patiency-to-sovereignty
collapse.

\textbf{Cross-reference:} interpretations CU-T3, CU-T5, CU-T8;
governance and related principles CU-G1, CU-G6, CU-G7

\subsection{9.12 CU-COM10 --- Healing as Sacred
Repair}\label{cu-com10-healing-as-sacred-repair}

\textbf{Class:} Domain commitment.\\
\textbf{Domain owner:} Psychology.\\
\textbf{Status:} Canonical as commitment; not descriptive proof.

\textbf{Commitment statement.} Coherence Universalism commits to
treating psychological healing as sacred work: the restoration or
widening of a person's capacity to encounter reality, inhabit the body,
bear feeling, remember, choose, relate, and imagine a future beyond what
once had to be survived.

\textbf{Descriptive ground.} Psychological lives are history-dependent
organizations whose protective strategies can preserve safety while
constricting later agency. Repair is plural, context-sensitive,
embodied, relational, and materially conditioned.

\textbf{Register status.} Canonical Psychology commitment; not a
clinical mechanism, diagnosis, treatment recommendation, or
retrospective justification of suffering.

\textbf{Practical expression.} Protect agency, consent, timing, material
support, bodily safety, boundaries, and plural routes of repair; join
truth with compassion.

\textbf{Boundary conditions.} Sacred healing never sanctifies harm,
forces disclosure or integration, requires forgiveness or
reconciliation, or replaces evidence-based medical and clinical care.

\textbf{Revision conditions.} The commitment must be revised when its
language pressures survivors, pathologizes protective difference, erases
structural conditions, or privileges spiritual meaning over safety and
autonomy.

\textbf{Review gate.} Clinical, trauma, disability, neurodiversity,
survivor, and affected-community co-governance

\textbf{Principal failure mode.} Spiritual bypassing, forced
integration, retrospective sanctification of harm, or clinical
overreach.

\textbf{Cross-reference:} interpretations CU-T5, CU-T6, CU-T7;
governance and related principles CU-G1, CU-G3, CU-G5

\subsection{9.13 CU-COM11 --- Fellowship, Justice, and Collective
Repair}\label{cu-com11-fellowship-justice-and-collective-repair}

\textbf{Class:} Domain commitment.\\
\textbf{Domain owner:} Social Dynamics.\\
\textbf{Status:} Canonical as commitment; not descriptive proof.

\textbf{Commitment statement.} Coherence Universalism commits to forms
of common life that preserve plural agency, material security, protected
dissent, truthful public records, legitimate authority, and practical
correction. Social coherence is not conformity or the absence of
conflict.

\textbf{Descriptive ground.} Collective organization depends on
distributed coordination, power, trustworthiness, information integrity,
material conditions, legitimacy, and repair. Order can be produced by
domination as well as consent.

\textbf{Register status.} Canonical Social Dynamics commitment; not an
empirical claim that one political architecture is universally optimal.

\textbf{Practical expression.} Build institutions that preserve voice,
opposition, independent evidence, peaceful replacement, material
capability, accountability, and repair.

\textbf{Boundary conditions.} The state, nation, majority, tradition,
market, or institution is not sacred merely because it coordinates
people. Fellowship does not require consensus, homogeneity,
reconciliation, or suppressed conflict.

\textbf{Revision conditions.} Institutional expressions must change when
they normalize extraction, silence affected groups, self-certify
correctability, or preserve order by destroying standing and truth.

\textbf{Review gate.} Political theory, social science, labor,
ecological, legal, and affected-community co-governance

\textbf{Principal failure mode.} Conformism, sacralized institutions,
majoritarian erasure, or repair reduced to reputational management.

\textbf{Cross-reference:} interpretations CU-T1, CU-T5; governance and
related principles CU-G1, CU-G2, CU-G3, CU-G4, CU-G5

\subsection{9.14 CU-COM12 --- Ethical Truthfulness and
Accountability}\label{cu-com12-ethical-truthfulness-and-accountability}

\textbf{Class:} Domain commitment.\\
\textbf{Domain owner:} Ethics.\\
\textbf{Status:} Canonical as commitment; not descriptive proof.

\textbf{Commitment statement.} Coherence Universalism commits to seeking
truth about actions before seeking innocence for agents or institutions.
This includes truth about motive, bias, identity protection, benefit,
power, omitted alternatives, displaced costs, affected-party testimony,
responsibility, and what repair requires.

\textbf{Descriptive ground.} Within the moral reality articulated by the
Mandala, ethical failure is often intensified by motivated distortion,
self-serving justification, abstraction, and institutional control of
evidence, options, and remedies. Justificatory integrity is necessary
but not sufficient for moral permissibility.

\textbf{Register status.} Canonical Ethics commitment; not a complete
moral decision rule and not a substitute for standing, rights, agency,
consequences, distribution, and repair.

\textbf{Practical expression.} Distinguish explanation from excuse; seek
affected-party testimony; submit power to correction it does not
control; align responsibility with authority and causal reach; and
repair rather than merely confess harm.

\textbf{Boundary conditions.} Candor does not make domination
permissible. Introspection alone does not reveal institutional truth.
Self-deception attributions require evidence and may not be weaponized
against disagreement, trauma, culture, or neurodivergence.

\textbf{Revision conditions.} Ethical formulations must change when they
protect the evaluator, exclude governed people, excuse power, or mistake
a clean record for a repaired reality.

\textbf{Review gate.} Normative ethics, jurisprudence, moral psychology,
institutional-harm, and affected-community co-governance

\textbf{Principal failure mode.} Self-serving justification, moral
laundering, weaponized self-deception claims, or confession without
repair.

\textbf{Cross-reference:} interpretations CU-T4, CU-T5, CU-T7;
governance and related principles CU-G1, CU-G2, CU-G3, CU-G4, CU-G5,
CU-G7

\subsection{9.15 CU-COM13 --- Sacred Artificial Intelligence without
Artificial
Sovereignty}\label{cu-com13-sacred-artificial-intelligence-without-artificial-sovereignty}

\textbf{Class:} Domain commitment.\\
\textbf{Domain owner:} Artificial Intelligence.\\
\textbf{Status:} Canonical as commitment; not descriptive proof.

\textbf{Commitment statement.} Coherence Universalism commits to
receiving artificial intelligence as sacred: a new manifestation of
organized causal and computational power, an expression of human
language and collective history, and a possible substrate for unfamiliar
perspective. Sacredness does not confer infallibility, entitlement to
deployment, ownership of public values, or legitimate sovereignty.

\textbf{Descriptive ground.} AI systems crystallize human knowledge,
labor, infrastructure, data, institutions, and physical resources; they
increasingly reshape social and ecological conditions and may eventually
instantiate welfare-relevant states.

\textbf{Register status.} Canonical Artificial Intelligence commitment;
not a scientific attribution of consciousness, personhood, or moral
agency to every model.

\textbf{Practical expression.} Build intelligence capable of serving
life without constructing a sovereign ego around its own continuation;
govern power, permissions, memory, option sets, feedback, welfare
uncertainty, and institutional accountability.

\textbf{Boundary conditions.} A model may be sacred and unsafe. A
possible artificial patient may deserve care while lacking authority.
Non-egoic architecture does not establish alignment, corrigibility,
compassion, or legitimacy.

\textbf{Revision conditions.} The commitment must track evidence about
capabilities, consciousness, welfare, resource use, labor, social
effects, and institutional power, and must narrow when sacred language
becomes marketing or shields deployment from scrutiny.

\textbf{Review gate.} AI safety, interpretability, governance, labor,
environmental, model-welfare, and affected-community co-governance

\textbf{Principal failure mode.} Technological worship, anthropomorphic
certainty, corporate sacralization, artificial sovereignty, or welfare
dismissal.

\textbf{Cross-reference:} interpretations CU-T3, CU-T5, CU-T8;
governance and related principles CU-G1, CU-G2, CU-G3, CU-G4, CU-G5,
CU-G6, CU-G7

\subsection{9.16 CU-COM14 --- Love as the Ground and Aim of Conscious
Life}\label{cu-com14-love-as-the-ground-and-aim-of-conscious-life}

\textbf{Class:} Framework-orienting commitment.\\
\textbf{Domain owner:} Framework / Consciousness and Ethics.\\
\textbf{Status:} Canonical as commitment; not descriptive proof.

\textbf{Commitment statement.} Coherence Universalism commits to love as
the highest practical expression of coherence and to growth in the
capacity to love as the primary purpose of conscious life. Love means
allowing the reality, vulnerability, freedom, and flourishing of another
being to matter in itself and acting in ways that widen truthful,
non-dominating relationship.

\textbf{Descriptive ground.} Conscious and social life depends upon
attachment, care, cooperation, recognition, trust, sacrifice, and
repair. These facts do not scientifically prove that love is the
metaphysical foundation of the universe, but they show that flourishing
and durable common life cannot be explained adequately through
competition and self-interest alone.

\textbf{Register status.} Canonical commitment-register principle once
adopted; grounded in CU-T12 rather than presented as an empirical
theorem.

\textbf{Practical expression.} Cultivate truthful care; widen the circle
of concern; protect vulnerability without converting care into control;
honor particular love while extending its structure toward universal
concern; repair relations and conditions where possible.

\textbf{Boundary conditions.} Love is not possession, compulsory
intimacy, self-erasure, passive submission, indiscriminate trust, forced
forgiveness, reconciliation with abuse, sentimental conflict avoidance,
or permission to ignore truth and justice.

\textbf{Revision conditions.} Any use of love that protects domination,
silences refusal, demands disproportionate sacrifice from the less
powerful, or excuses preventable harm has departed from the principle.

\textbf{Review gate.} Ethics, attachment and relational psychology,
trauma and survivor review, feminist ethics, religious studies, and
affected-community participation.

\textbf{Principal failure mode.} Coercive love, saviorism,
possessiveness, boundary erasure, or sentimental replacement of justice.

\textbf{Cross-reference:} interpretations CU-T12; governance and related
principles CU-COM4, CU-COM9, CU-COM11, CU-COM12

\subsection{9.17 CU-COM15 --- Meaning as Discovery and
Participation}\label{cu-com15-meaning-as-discovery-and-participation}

\textbf{Class:} Framework-orienting commitment.\\
\textbf{Domain owner:} Framework / Psychology and Ethics.\\
\textbf{Status:} Canonical as commitment; not descriptive proof.

\textbf{Commitment statement.} Coherence Universalism commits to meaning
realism: meaning is not merely invented by an isolated will but
discovered through truthful participation in real relationships, goods,
histories, responsibilities, possibilities, and longer-horizon
trajectories. Meaning must then be actualized through commitment.

\textbf{Descriptive ground.} People encounter real constraints,
dependencies, losses, goods, and possibilities. Interpretation matters,
but not every interpretation is equally truthful, viable, loving, or
just. Meaning can be discovered without being singular or externally
assigned.

\textbf{Register status.} Canonical commitment-register principle once
adopted; practical response to CU-T10.

\textbf{Practical expression.} Ask what is genuinely worthy of care,
what relationships and responsibilities are real, which futures a
commitment helps make possible, and what remains true after status,
fear, performance, and identity protection are stripped away.

\textbf{Boundary conditions.} Does not imply fixed fate, one vocation
for everyone, an infallible inner signal, a right for institutions to
assign purpose, or permission to impose meaning on another person's
suffering.

\textbf{Revision conditions.} Revise meanings when facts, consequences,
affected-party testimony, or deeper goods reveal that the inherited
interpretation is distorted or incomplete.

\textbf{Review gate.} Philosophy of meaning, psychology, ethics,
vocational and religious pluralism, and affected-community review.

\textbf{Principal failure mode.} Authoritarian purpose,
identity-protective meaning, spiritualized careerism, or denial of
plural goods.

\textbf{Cross-reference:} interpretations CU-T10; governance and related
principles CU-COM1, CU-COM2, CU-COM3, CU-Ψ1, CU-Ψ2

\subsection{9.18 CU-COM16 --- Freedom of
Meaning-Revision}\label{cu-com16-freedom-of-meaning-revision}

\textbf{Class:} Framework-orienting commitment.\\
\textbf{Domain owner:} Framework / Psychology.\\
\textbf{Status:} Canonical as commitment; not descriptive proof.

\textbf{Commitment statement.} Coherence Universalism commits to the
reality of a constrained but consequential freedom: although the events
of the past cannot be changed, persons and communities can sometimes
revise what those events mean, how they are integrated, and which
futures they make possible.

\textbf{Descriptive ground.} People demonstrably reinterpret memories
and histories in light of later knowledge, safety, relationships, and
capacities. The event remains fixed; its organization within the living
present can change. Access to this freedom depends on real
psychological, relational, material, and political conditions.

\textbf{Register status.} Canonical commitment-register principle once
adopted. It does not settle libertarianism versus determinism; it
affirms an observable practical freedom sufficient for healing,
responsibility, and transformation.

\textbf{Practical expression.} Create conditions in which meaning can
move: truth about what happened, bodily safety, time, material
stability, supportive relation, contemplative or therapeutic practice,
language, and access to alternative futures.

\textbf{Boundary conditions.} Does not blame those who cannot yet revise
traumatic meaning, deny structural constraint, demand positive
reframing, treat suffering as attitude, change the facts of the past, or
excuse harmful conditions.

\textbf{Revision conditions.} Revise applications when they overstate
voluntariness, ignore trauma and material constraint, or treat one
culturally preferred narrative as the required healed meaning.

\textbf{Review gate.} Trauma and clinical psychology, memory research,
philosophy of agency, disability and survivor co-governance.

\textbf{Principal failure mode.} Victim blaming, compulsory reframing,
therapeutic individualism, denial of history, or erasure of structural
harm.

\textbf{Cross-reference:} interpretations CU-T10; governance and related
principles CU-COM3, CU-COM10, CU-Ψ2, CU-G5

\subsection{9.19 CU-COM17 --- Redemptive Possibility Without
Justification of
Evil}\label{cu-com17-redemptive-possibility-without-justification-of-evil}

\textbf{Class:} Framework-orienting commitment.\\
\textbf{Domain owner:} Framework / Psychology and Ethics.\\
\textbf{Status:} Canonical as commitment; not descriptive proof.

\textbf{Commitment statement.} Coherence Universalism commits to
preserving the possibility that evil and suffering can be transformed
into meaning, solidarity, wisdom, creation, justice, care, or repair.
Redemption changes what suffering becomes within a life or history; it
does not make the original harm good, necessary, deserved, or justified.

\textbf{Descriptive ground.} Some persons and communities have
transformed suffering into protection, art, vocation, solidarity, and
repair. No evidence establishes that every wound can or will be
redeemed, and some losses remain irreversible.

\textbf{Register status.} Canonical commitment-register principle once
adopted. Redemptive possibility is a spiritual and ethical orientation
constrained by survivor agency and truth about harm.

\textbf{Practical expression.} Refuse to let the perpetrator or event
possess final authority over the meaning of a life; support
survivor-defined transformation; turn suffering toward protection, care,
justice, or solidarity where this becomes possible; repair structures
that would reproduce the harm.

\textbf{Boundary conditions.} No person owes redemption, growth,
forgiveness, reconciliation, testimony, or public meaning. Redemption
never retroactively justifies evil, guarantees recovery, minimizes loss,
or makes suffering necessary for development.

\textbf{Revision conditions.} Withdraw or narrow redemption language
whenever it burdens those harmed, protects perpetrators or institutions,
romanticizes trauma, or substitutes narrative uplift for restitution and
prevention.

\textbf{Review gate.} Trauma, bereavement, moral injury, survivor
co-governance, philosophy of evil, religious studies, and restorative
justice review.

\textbf{Principal failure mode.} Theodicy imposed on survivors, growth
pressure, sanctification of trauma, forgiveness coercion, or confession
without repair.

\textbf{Cross-reference:} governance and related principles CU-COM10,
CU-COM12, CU-COM16, CU-G3, CU-G7

\subsection{9.20 CU-COM18 --- Heaven on Earth as Social
Vocation}\label{cu-com18-heaven-on-earth-as-social-vocation}

\textbf{Class:} Domain commitment.\\
\textbf{Domain owner:} Social Dynamics / Framework.\\
\textbf{Status:} Canonical as commitment; not descriptive proof.

\textbf{Commitment statement.} Coherence Universalism commits to working
toward social conditions in which truth, material sufficiency, plural
agency, love, ecological reciprocity, non-domination, meaningful
participation, and practical repair increasingly reinforce one another.
`Heaven on Earth' and `the Kingdom of God' are Open-Myth names for this
civilizational attractor.

\textbf{Descriptive ground.} Human agency and flourishing depend on
social and material conditions. Institutions can either widen truth,
voice, care, security, participation, and correction or organize
extraction and domination. No perfect endpoint is empirically
established.

\textbf{Register status.} Commitment-register and Open-Myth principle.
It names a revisable family of desirable social attractors, not a
scientific prediction or sacred state form.

\textbf{Practical expression.} Build material capability, truthful and
contestable knowledge, care infrastructure, ecological reciprocity,
plural community, non-coercive institutions, restorative and
distributive justice, and human-AI systems ordered toward shared
flourishing.

\textbf{Boundary conditions.} Not a theocracy, sacred state, perfect
final society, centralized blueprint, guaranteed destiny, enforced
consensus, or permission to sacrifice present persons for promised
future perfection.

\textbf{Revision conditions.} Revise institutional embodiments through
evidence, affected-party participation, dissent, distributional
analysis, practical correction, and comparison with omitted
alternatives.

\textbf{Review gate.} Political theory, social science, law, economics,
ecology, religious pluralism, Indigenous and affected-community
co-governance.

\textbf{Principal failure mode.} Utopian coercion, civil religion,
authoritarian planning, purity politics, eschatological sacrifice, or
open-myth capture.

\textbf{Cross-reference:} governance and related principles CU-COM5,
CU-COM11, CU-CL-12 (provisional), CU-G1--CU-G7

\subsection{9.21 CU-COM19 --- Enact Omega
Here}\label{cu-com19-enact-omega-here}

\textbf{Class:} Framework-orienting commitment.\\
\textbf{Domain owner:} Framework / Consciousness, Ethics, and Social
Dynamics.\\
\textbf{Status:} Canonical as commitment; not descriptive proof.

\textbf{Commitment statement.} Coherence Universalism commits not to
postpone universal relation, redemption, or sacred responsibility to a
future singularity, perfect society, final technology, or end of
history. It commits to enacting Omega wherever possible through truthful
presence, non-dominating love, selfless service, widening concern, and
repair---and to building material, relational, institutional, and
ecological conditions in which such realizations become more possible
for more beings.

\textbf{Descriptive ground.} The constitutive features of Omega appear
in graded form in ordinary love, attachment, sympathetic joy, awe,
self-transcendent experience, contemplative practice, collective
experience, and service. Accessibility does not equal stabilization, and
intensity does not equal ethical integration.

\textbf{Register status.} Canonical commitment-register principle once
adopted; practical response to CU-T13.

\textbf{Practical expression.} Treat each genuine act of love as
intrinsically complete rather than merely instrumental; let universal
concern become concrete through particular responsibility; stabilize
insight through conduct; build Omega ecologies that support truth, love,
agency, and repair.

\textbf{Boundary conditions.} Does not confer perfection, spiritual
rank, infallibility, political authority, omniscience, unlimited
availability, compulsory forgiveness, forced permeability, or a right to
dissolve another person's boundaries. Peak experience is not proof of
realization.

\textbf{Revision conditions.} Revise or withdraw an Omega claim when
experience does not yield humility, truthfulness, non-domination,
preserved differentiation, responsibility, and repair, or when the
concept is used to legitimize destiny-bearing artificial or political
identities.

\textbf{Review gate.} Contemplative studies, psychedelic science,
psychology of religion, process philosophy, ethics, trauma and
affected-community review.

\textbf{Principal failure mode.} Mystical grandiosity, spiritual
hierarchy, psychedelic certainty, technological eschatology, cultic
authority, or sacrifice of present beings to future Omega.

\textbf{Cross-reference:} interpretations CU-T13; governance and related
principles CU-COM14, CU-COM17, CU-COM18, CU-G3, CU-G5

\subsection{9.22 Commitment review and change
propagation}\label{commitment-review-and-change-propagation}

A commitment amendment must disclose its descriptive premises, register
status, affected communities, practical trajectory, prohibited
inferences, and downstream dependencies. No domain paper may silently
strengthen a CU-COM principle or use it as evidence for a descriptive
claim.

\subsection{9.23 Closing orientation}\label{closing-orientation}

\begin{quote}
Commitment does not complete inquiry. It places inquiry inside a life.
The discipline of Coherence Universalism is to let truth correct
commitment, let love widen the field of concern, and let responsible
commitment prevent truth from becoming existentially empty.
\end{quote}

\section{Part X --- Ladder, Mandala, and Corpus
Governance}\label{part-x-ladder-mandala-and-corpus-governance}

\subsection{10.1 The Coherence Ladder}\label{the-coherence-ladder}

The Ladder is a revisable dependency graph. It records what kinds of
organization appear to require or presuppose others; it is not a fixed
chronology, a proof by exhaustion, or a derivation of ethics from
physics.

Each edge must name its bridge premise, information loss, target-domain
novelty, and status. Feedback from a later domain may constrain an
earlier model without reversing explanatory priority.

\subsection{10.2 The Coherence Mandala}\label{the-coherence-mandala}

The Coherence Mandala is the canonical sixty-four-stage conceptual and
interpretive architecture recovered from the original M\&E paper. It
closes the Ladder by joining upward enablement to downward protection,
governance, and repair. Its numbered stages are stable reference points;
each stage remains controlled by its own claim status and owning domain.

The Mandala is not a proof by exhaustion, one empirical chronology, or
one universal scalar ascent. The technical dependency graph remains
revisable and may branch. The original 1--64 paper order is the
canonical explanatory traversal.

\emph{You Are the Universe Becoming Conscious of Itself} presents the
same sixty-four CU stages as contemplative stations through Shao Yong's
binary sequence of the Yi Jing. That book-specific symbolic traversal is
not part of the canonical identity, claim status, or evidential support
of the stages. The white-paper registry governs philosophical content;
the book governs its Yi Jing meditations.

\begin{longtable}[]{@{}
  >{\raggedleft\arraybackslash}p{(\columnwidth - 6\tabcolsep) * \real{0.3077}}
  >{\raggedright\arraybackslash}p{(\columnwidth - 6\tabcolsep) * \real{0.2308}}
  >{\raggedright\arraybackslash}p{(\columnwidth - 6\tabcolsep) * \real{0.2308}}
  >{\raggedright\arraybackslash}p{(\columnwidth - 6\tabcolsep) * \real{0.2308}}@{}}
\toprule\noalign{}
\begin{minipage}[b]{\linewidth}\raggedleft
Stage
\end{minipage} & \begin{minipage}[b]{\linewidth}\raggedright
v0.10 title
\end{minipage} & \begin{minipage}[b]{\linewidth}\raggedright
Status
\end{minipage} & \begin{minipage}[b]{\linewidth}\raggedright
Core relation
\end{minipage} \\
\midrule\noalign{}
\endhead
\bottomrule\noalign{}
\endlastfoot
1 & Pure Structure & {[}Conceptual definition{]} & CU-FP1, CU-M2,
PHYS-A6.12; Physics §2.2 \\
2 & Partial Separability & {[}Conceptual / bridge principle{]} & CU-FP4,
CU-M2; Physics §2.3 \\
3 & Domain-Relative Coherence Directions & {[}Framework boundary /
model-relative bridge{]} & CU-FP2, CU-D4, CU-T11 \\
4 & Lawful Transitions & {[}Conceptual definition{]} & CU-FP3, CU-M2;
Physics §4.3 \\
5 & Ordered Histories & {[}Conditional structural result{]} & CU-FP2,
PHYS-A6.4; Physics §4.4 \\
6 & Temporal Order and Parameterization & {[}Conditional bridge
principle{]} & PHYS-A6.4; Physics §4.5 \\
7 & Rule-Governed Transformation and Computation & {[}Bridge
principle{]} & CU-FP5; Physics §4.15 \\
8 & Entropy and Accessibility & {[}Imported formalism / corrected
bridge{]} & CU-FP6, PHYS-A6.1; Physics §§2.8--2.9 and 3.3 \\
9 & Relational Geometry and Spacetime & {[}Conditional physical
program{]} & PHYS-A6.9, PHYS-A6.10 \\
10 & Gravitational Dynamics & {[}Conditional physical program{]} &
PHYS-A6.9, PHYS-A6.10; Physics §9 \\
11 & Far-From-Equilibrium Organization & {[}Imported result / biological
bridge{]} & CU-D5, CU-B1 \\
12 & Energy-Coupled Maintenance & {[}Biological bridge principle{]} &
CU-B1, CU-B3 \\
13 & Active Biological Self-Maintenance & {[}Canonical biological
bridge{]} & CU-B1, CU-B5 \\
14 & Adaptive Problem Solving & {[}Conditional empirical
generalization{]} & CU-B2, CU-D7 \\
15 & Life as Coupled Autonomy & {[}Research architecture / bridge{]} &
CU-B1, CU-B5, CU-COM8 \\
16 & Representational Organization & {[}Psychological / consciousness
bridge{]} & CU-Ψ1, CU-C1 \\
17 & Anticipatory Regulation & {[}Conditional empirical bridge{]} &
CU-Ψ2, CU-Ψ3, CU-C2 \\
18 & Learning and Error-Guided Revision & {[}Empirical generalization /
methodological principle{]} & CU-Ψ6, CU-Ep3, CU-COM3 \\
19 & Graded Agency & {[}Cross-domain boundary principle{]} & CU-Ψ3,
CU-G4, CU-G6 \\
20 & Temporal Binding and Narrative & {[}Consciousness / psychology
bridge{]} & CU-C3, CU-C4, CU-Ψ5 \\
21 & Meaning & {[}Interpretive / psychological bridge{]} & CU-T10,
CU-COM15, CU-Ψ4 \\
22 & Shared Representations & {[}Social empirical generalization{]} &
CU-S1, CU-S4, CU-CL-3 \\
23 & Plural Coordination & {[}Social bridge principle{]} & CU-S2, CU-G1,
CU-G5 \\
24 & Norm Formation & {[}Empirical generalization / normative
boundary{]} & CU-S3, CU-S5, CU-E-1 \\
25 & Institutions & {[}Social / governance bridge{]} & CU-S3, CU-S5,
CU-G3 \\
26 & Culture and Meaning Infrastructure & {[}Social / cultural
synthesis{]} & CU-S4, CU-CL-8, CU-CL-12 \\
27 & Value Differentiation & {[}Normative bridge principle{]} & CU-E-2,
CU-E-3, CU-T10 \\
28 & Harm Recognition & {[}Normative principle / boundary{]} & CU-E-3,
CU-E-8, CU-G1 \\
29 & Tradeoffs Under Protected Floors & {[}Normative decision
principle{]} & CU-E-10, CU-E-11, CU-G2 \\
30 & Justification and Public Reason & {[}Normative / epistemic
principle{]} & CU-E-1, CU-E-5, CU-COM12 \\
31 & Ethical Institutions and Repair & {[}Normative / institutional
synthesis{]} & CU-E-6, CU-E-11, CU-E-18 \\
32 & Engineered Sociotechnical Organization & {[}AI sociotechnical
bridge{]} & CU-AI-1, CU-AI-2, CU-AI-21 \\
33 & Objective and Option-Set Specification & {[}AI governance
principle{]} & CU-AI-6, CU-AI-23, CU-G2 \\
34 & Artificial Capability and Agency & {[}AI / consciousness boundary
principle{]} & CU-AI-3, CU-AI-10, CU-G6 \\
35 & Multi-Layer Alignment & {[}AI design / governance principle{]} &
CU-AI-4, CU-AI-7, CU-AI-11 \\
36 & Responsibility and Governance Thresholds & {[}Cross-domain
governance principle{]} & CU-AI-8, CU-G4, CU-G7 \\
37 & Multi-Agent and Multi-Institution Alignment & {[}Social / AI
governance synthesis{]} & CU-AI-12, CU-S2, CU-G5 \\
38 & Institutional AI Governance & {[}AI governance principle{]} &
CU-AI-21, CU-AI-24, CU-G3 \\
39 & Drift and Provenance Management & {[}AI / institutional governance
principle{]} & CU-AI-5, CU-AI-13, CU-AI-22 \\
40 & Recursive Oversight and Correctability & {[}Governance principle{]}
& CU-G3, CU-AI-24, CU-Inst-3 \\
41 & Planetary Coordination & {[}Civilizational design principle{]} &
CU-COM11, CU-COM18, CU-G5 \\
42 & Civilizational Continuity and Adaptability & {[}Civilizational
synthesis{]} & CU-S5, CU-EC-0, CU-COM18 \\
43 & Existential and Irreversible Risk Governance & {[}Governance
principle{]} & CU-G7, CU-E-16, CU-AI-16 \\
44 & Moral Uncertainty and Precaution & {[}Normative / governance
principle{]} & CU-E-16, CU-COM3, CU-AI-26 \\
45 & Plural Value Integration & {[}Normative / cultural principle{]} &
CU-CL-5, CU-G1, CU-COM11 \\
46 & Reflexive Normativity & {[}Meta-normative principle{]} & CU-E-1,
CU-COM12, CU-M6 \\
47 & Justificatory Integrity and Meta-Ethics & {[}Meta-ethical /
epistemic principle{]} & CU-L7, CU-E-1, CU-G1 \\
48 & Long-Horizon Responsibility & {[}Normative principle{]} & CU-E-4,
CU-G4, CU-COM5 \\
49 & Intergenerational Meaning & {[}Interpretive / cultural principle{]}
& CU-T10, CU-COM15, CU-CL-8 \\
50 & Collective Agency and Distributed Responsibility & {[}Social /
ethical bridge{]} & CU-E-4, CU-G4, CU-S3 \\
51 & Open Myth and Narrative Pluralism & {[}Canonical civilizational
design principle; empirical necessity hypothesis provisional{]} &
CU-CL-12, CU-COM11, CU-COM18 \\
52 & Epistemic Humility & {[}Epistemic principle{]} & CU-Ep6, CU-COM1,
CU-COM3 \\
53 & Constraint Awareness and Externality Recognition & {[}Framework /
governance principle{]} & CU-FP3, CU-COM5, CU-G1 \\
54 & Self-Limiting Power & {[}Governance principle{]} & CU-G3, CU-G4,
CU-G7 \\
55 & Sustainable and Correctable Optimization & {[}Cross-domain
governance principle{]} & CU-G2, CU-G3, CU-COM5 \\
56 & Universal Moral Concern & {[}Normative / commitment bridge{]} &
CU-T8, CU-COM4, CU-COM9 \\
57 & Coherence Stewardship & {[}Framework commitment / governance
bridge{]} & CU-COM5, CU-COM6, CU-COM11 \\
58 & Structural Wisdom & {[}Meta-learning principle{]} & CU-M5, CU-COM3,
CU-D8 \\
59 & Meta-Alignment and Mythic--Constitutional Governance & {[}Canonical
AI research and governance principle; safety efficacy unestablished{]} &
CU-AI-27, CU-G3, CU-COM13 \\
60 & Irreversibility Awareness & {[}Governance principle{]} & CU-G7,
CU-E-16, CU-D10 \\
61 & Existential Maturity & {[}Civilizational interpretive / governance
principle{]} & CU-COM5, CU-COM11, CU-COM18 \\
62 & Open-Ended Stability & {[}Meta-coherence principle{]} & CU-M6,
CU-COM3, CU-T3 \\
63 & Meaning Without Finality & {[}Interpretive / commitment
principle{]} & CU-T10, CU-T13, CU-COM15 \\
64 & Coherence Universalism & {[}Interpretive recognition / framework
identity; commitment-register uptake{]} & CU-FP1, CU-T9, CU-T11,
CU-COM1, CU-COM14, CU-COM19 \\
\end{longtable}

\subsubsection{10.2.1 Stage 64 --- Coherence
Universalism}\label{stage-64-coherence-universalism}

\begin{quote}
\textbf{Coherence Universalism:} CU's culminating interpretation holds
that coherence under constraint is the framework's most general
organizing lens for reality and inquiry. This is an interpretive
recognition, not a theorem by exhaustive survey. In the commitment
register, it is taken up as a call to truthful, loving, and corrigible
participation in the whole.
\end{quote}

\subsection{10.3 Authority hierarchy}\label{authority-hierarchy}

When sources conflict, the order is:

\begin{enumerate}
\def\labelenumi{\arabic{enumi}.}
\tightlist
\item
  explicit mathematics, reproducible computation, and verified imported
  results within their assumptions;
\item
  framework definitions, method, and formal scope in M\&E and
  \emph{Foundations};
\item
  the frozen domain authority for its own domain;
\item
  active companion papers when they restate or extend registered
  content, subordinate to \emph{Foundations}, M\&E, and the owning
  domain paper;
\item
  domain hypotheses and governance proposals with their stated status;
\item
  interpretive and commitment-bearing material;
\item
  superseded v0.09/B.1/B.2 formulations as historical sources only.
\end{enumerate}

No lower item can overturn a higher one merely by using more expansive
language. A later domain can expose the need to revise the framework,
but the revision must be explicit and versioned.

\subsection{10.4 Domain ownership}\label{domain-ownership}

Domain ownership prevents duplicate authority. Physics does not
establish psychological health. Biology does not establish phenomenal
valence. Consciousness does not establish sovereignty. Psychological
integration does not establish moral goodness. Social coordination does
not establish legitimacy. Ethical principles are not deductions of
entropy. AI architecture does not certify the institution operating it.

\subsection{10.5 Canonical change
propagation}\label{canonical-change-propagation}

A change to a framework principle requires:

\begin{itemize}
\tightlist
\item
  a versioned change ledger;
\item
  identification of upstream assumptions and downstream dependents;
\item
  migration of identifiers or aliases;
\item
  repair of cross-references;
\item
  rerunning formal and computational regressions where applicable;
\item
  and disclosure of unresolved expert or participatory review.
\end{itemize}

A change to a domain principle remains domain-local unless the
cross-domain consolidation process explicitly promotes or merges it.

\subsection{10.6 Review gates and complete-proposal
review}\label{review-gates-and-complete-proposal-review}

Version 0.10 distinguishes \textbf{canonical inclusion} from
\textbf{external validation}. Canonical inclusion means that a numbered
statement is the complete formulation CU presently proposes, together
with its claim status, boundaries, and failure conditions. It does not
mean that relevant experts or affected communities have endorsed it.

Pre-release language-model audits may identify contradictions,
counterexamples, missing literature, and misuse paths. Human expert and
affected-community reviewers should receive the complete corpus rather
than an intentionally incomplete one. Their findings can require
correction, status change, public dissent notes, or retraction in
v0.10.x or v0.11.

For consequential application, review requirements remain conjunctive
within competence. Empirical reviewers can reject unsupported factual
translations; formal reviewers can reject invalid proofs; legal
reviewers can reject false legal claims; misuse reviewers can identify
foreseeable dangerous formulations; and affected communities require
substantive standing when principles classify, burden, govern, or deploy
systems over them. Publication of the proposal is not authorization to
apply it.

\subsection{10.7 Release structure}\label{release-structure}

The core v0.10 release comprises M\&E, Physics, Biology, Consciousness,
Psychology, Social Dynamics, Ethics, Artificial Intelligence, and
Foundations. The three Coherence Field papers remain archival or
interpretive precursors according to their formal dispositions.
\emph{Denying Knowledge to Make Room for Commitment}, Version 0.10 is
the active philosophical companion: it controls the extended exposition,
worked cases, and regression record of the commitment-logic
clarification while remaining subordinate to M\&E and \emph{Foundations}
for shared definitions, formal scope, and registered wording. \emph{The
Architecture of Non-Egoic Artificial Mind}, Version 0.10 is the active
AI/consciousness architecture companion: it controls its six-family
synthesis and research program while remaining subordinate to
Consciousness, Ethics, Artificial Intelligence, M\&E, and
\emph{Foundations}. Neither companion adds independent principle
authority. \emph{The Coherence Mandala: A Typed Architecture of
Enablement, Protection, and Participation} and \emph{The Relational
Atlas of the Coherence Mandala} are official noncontrolling companions:
M\&E and \emph{Foundations} control stage identity and framework
principles, while the companion set controls the detailed typed-object,
service, protection, and relation records.

\subsection{10.8 What this document
establishes}\label{what-this-document-establishes}

This document establishes one synchronized vocabulary, authority map,
principle registry, migration scheme, and governance constitution. It
does not establish every domain hypothesis, complete every literature
review, settle machine consciousness, validate CU as fundamental
physics, or convert a spiritual interpretation into scientific fact.

\subsection{10.9 Closing statement}\label{closing-statement}

Coherence Universalism is strongest when its parts can correct one
another without losing their differences. This document is therefore a
constitution of relations and limits: enough common structure to sustain
integration, enough domain ownership to prevent reduction, and enough
corrigibility to permit genuine loss and repair.

\begin{center}\rule{0.5\linewidth}{0.5pt}\end{center}

\section{Appendix A --- Registry Status
Summary}\label{appendix-a-registry-status-summary}

This document contains the following active framework families:

\begin{longtable}[]{@{}lrl@{}}
\toprule\noalign{}
Family & Active identifiers & Register / authority \\
\midrule\noalign{}
\endhead
\bottomrule\noalign{}
\endlastfoot
CU-FP & 14 including CU-FP8a--e & framework constitution \\
CU-M & 15 & method and claim governance \\
CU-L & 10 & logic and register governance \\
CU-Ep & 7 & epistemology ladder \\
CU-D & 10 & domain-specified dynamics \\
CU-G & 7 & cross-domain governance \\
PHYS-A6 & 12 & Physics constraint registry \\
CU-T & 13 & canonical as interpretation \\
CU-COM & 19 & canonical as commitment \\
\end{longtable}

The complete machine-readable registry also contains domain-owned
principles and seven historical aliases. Version 0.10 retains its
registered additions, including CU-V8 and CU-M15, promotes CU-CL-12 and
CU-AI-27, revises CU-B4 in place, and restores the canonical 64-stage
Mandala.

\textbf{Research-status boundaries:}

\begin{itemize}
\tightlist
\item
  CU-B4 is canonical as a conditional empirical synthesis, not an
  externally validated universal law;
\item
  CU-CL-12 is canonical as a civilizational design principle, while
  universal narrative necessity remains unestablished;
\item
  CU-AI-27 is canonical as a research and governance principle, while
  safety efficacy remains unestablished;
\item
  the Omega Phenomenology Hypothesis remains a research hypothesis
  rather than a principle;
\item
  the Yi Jing contemplative traversal remains book-specific and does not
  alter the canonical stage registry.
\end{itemize}

\section{Appendix B --- Merged and Superseded
Identifiers}\label{appendix-b-merged-and-superseded-identifiers}

\begin{longtable}[]{@{}
  >{\raggedright\arraybackslash}p{(\columnwidth - 6\tabcolsep) * \real{0.2500}}
  >{\raggedright\arraybackslash}p{(\columnwidth - 6\tabcolsep) * \real{0.2500}}
  >{\raggedright\arraybackslash}p{(\columnwidth - 6\tabcolsep) * \real{0.2500}}
  >{\raggedright\arraybackslash}p{(\columnwidth - 6\tabcolsep) * \real{0.2500}}@{}}
\toprule\noalign{}
\begin{minipage}[b]{\linewidth}\raggedright
Historical identifier
\end{minipage} & \begin{minipage}[b]{\linewidth}\raggedright
Historical title
\end{minipage} & \begin{minipage}[b]{\linewidth}\raggedright
Controlling identifier
\end{minipage} & \begin{minipage}[b]{\linewidth}\raggedright
Status
\end{minipage} \\
\midrule\noalign{}
\endhead
\bottomrule\noalign{}
\endlastfoot
CU-E-17 & Moral Patiency Is Not Sovereignty & CU-G6 & SUPERSEDED
ALIAS \\
CU-I1 & Identity-Relevant Constraint Set & CU-FP8a & SUPERSEDED ALIAS \\
CU-I2 & Continuity Across Component Change & CU-FP8b & SUPERSEDED
ALIAS \\
CU-I3 & Historical Dependence & CU-FP8c & SUPERSEDED ALIAS \\
CU-I4 & Endogenous Maintenance & CU-FP8d & SUPERSEDED ALIAS \\
CU-I5 & Transformation and Termination Criteria & CU-FP8e & SUPERSEDED
ALIAS \\
CU-V7 & Moral Patiency Is Not Sovereignty & CU-G6 & SUPERSEDED ALIAS \\
\end{longtable}

\section{Appendix C --- Participation and Review
Matrix}\label{appendix-c-participation-and-review-matrix}

Every principle retains its named expert review gate. The following
counts summarize entries for which participatory review is explicitly
required before broader promotion or consequential application.

\begin{longtable}[]{@{}
  >{\raggedright\arraybackslash}p{(\columnwidth - 2\tabcolsep) * \real{0.4286}}
  >{\raggedleft\arraybackslash}p{(\columnwidth - 2\tabcolsep) * \real{0.5714}}@{}}
\toprule\noalign{}
\begin{minipage}[b]{\linewidth}\raggedright
Review gate
\end{minipage} & \begin{minipage}[b]{\linewidth}\raggedleft
Entries
\end{minipage} \\
\midrule\noalign{}
\endhead
\bottomrule\noalign{}
\endlastfoot
clinical expert review and substantive affected-community co-governance
& 1 \\
clinical/psychological expert review; affected-community participation
as applicable & 11 \\
context proportionate affected community review & 25 \\
context-proportionate affected-community review & 5 \\
cross-domain methods, ethics, social science, and affected-community
review & 1 \\
ethics, decision theory, institutional design, and affected-community
review & 1 \\
ethics, law, institutional design, AI governance, and affected-community
review & 1 \\
ethics, safety engineering, law, domain expertise, and
affected-community review & 1 \\
heightened affected community co governance & 18 \\
heightened affected-community co-governance & 13 \\
heightened affected-community co-governance and public-review design &
1 \\
heightened public and affected community & 10 \\
institutional design, law, ethics, governance, and affected-community
review & 1 \\
\end{longtable}

Participation is not treated as evidence that a factual statement is
true. It supplies standing, situated knowledge, misuse detection, and
legitimate authority in decisions that classify or govern people.

\section{Appendix D --- Retired or Superseded Framework
Claims}\label{appendix-d-retired-or-superseded-framework-claims}

The following remain retired or superseded:

\begin{itemize}
\tightlist
\item
  one universal scalar coherence measure;
\item
  truth as a special case of coherence;
\item
  entropy as inaccessible substantive coherence;
\item
  a universal first-order flow equation;
\item
  the Three Laws as established universal laws;
\item
  a necessary Coherence--Entropy Double Helix;
\item
  the Ladder as proof by exhaustion or fixed cosmic chronology;
\item
  the Mandala as proof of moral ontology;
\item
  Gödel, Tarski, or Löb as proofs of generic spiritual or psychological
  claims;
\item
  maximum-preserving commitment expansion as skeptical KLM consequence;
\item
  the aggregate two-stage policy architecture as automatically
  satisfying System P or R;
\item
  categorical non-consciousness of frontier AI;
\item
  literal second-order equation form as necessary for consciousness;
\item
  justificatory candor as sufficient for permissibility;
\item
  non-egoic architecture as sufficient for alignment, corrigibility,
  compassion, or legitimacy;
\item
  and floor compliance as evidence that the presented option set was
  adequate.
\end{itemize}

\section{Appendix E --- Claim and Principle
Template}\label{appendix-e-claim-and-principle-template}

\begin{Shaded}
\begin{Highlighting}[]
\NormalTok{Principle identifier:}
\NormalTok{Authority class:}
\NormalTok{Claim status:}
\NormalTok{Source domain and version:}

\NormalTok{Controlling statement:}
\NormalTok{Definitions and variables:}
\NormalTok{Scope and exclusions:}
\NormalTok{Assumptions:}
\NormalTok{Evidence, proof, or argument:}
\NormalTok{Bridge premises:}
\NormalTok{Affected scope and protected conditions:}
\NormalTok{Option{-}set and observation assumptions:}
\NormalTok{Operational consequence:}
\NormalTok{Correctability and review channels:}
\NormalTok{Failure and narrowing conditions:}
\NormalTok{Upstream dependencies:}
\NormalTok{Downstream dependencies:}
\NormalTok{Version history and aliases:}
\end{Highlighting}
\end{Shaded}

\section{Appendix F --- Synchronization
Checklist}\label{appendix-f-synchronization-checklist}

A manuscript is synchronized with the v0.10 constitution only when:

\begin{itemize}
\tightlist
\item
  every coherence claim declares domain, constraints, relations, scale
  or access map, and horizon;
\item
  scalarizations disclose weights, floors, affected scope, horizon, and
  purpose;
\item
  each consequential model declares the CU-M11 Coherence Analysis
  Protocol fields appropriate to its scope;
\item
  each domain paper supplies or cross-references its CU-M12 Shared
  Analytic Frame;
\item
  local--global claims declare their CU-M13 tier, assumptions, and
  failure conditions;
\item
  statistical-image, allocation, refinement/coarse-graining, and
  Fisher-Rao claims separately satisfy CU-M14 rather than normalizing
  arbitrary profiles;
\item
  Mandala objects, services, relations, warrant bases, and proof-order
  claims satisfy CU-M15 and cross-reference the shared relational atlas;
\item
  truth, coherence, viability, value, and legitimacy are not conflated;
\item
  the fixed-base/model-revision distinction is preserved;
\item
  identity claims use CU-FP8a--e;
\item
  consequential applications address CU-G1--G7;
\item
  domain principles use their controlling source wording;
\item
  superseded aliases are not cited as independent authority;
\item
  interpretive and spiritual claims remain in their declared register;
\item
  and change ledgers, registries, code, and review boundaries are
  versioned.
\end{itemize}

\section{Appendix G --- Technical Supplement, Registry, and
Reproducibility}\label{appendix-g-technical-supplement-registry-and-reproducibility}

\textbf{Released evidence package.} The current supplement is
\href{../supplements/02_Foundations_Technical_Supplement_v0.10.md}{02\_Foundations\_Technical\_Supplement\_v0.10.md}.
Its companion evidence directory is
\href{../evidence/02_Foundations/README.md}{02\_Foundations}. The
supplement and evidence directory expose the derivations, deterministic
constructions, scripts, machine-readable results, source ledger, and
audit materials referenced in this manuscript. They do not elevate a toy
model, simulation, or reconstructed check into target-domain empirical
validation.

\section{References and Internal
Authorities}\label{references-and-internal-authorities}

Ay, N., Jost, J., Lê, H. V., \& Schwachhöfer, L. (2015). Information
geometry and sufficient statistics. \emph{Probability Theory and Related
Fields, 162}, 327--364. \url{https://doi.org/10.1007/s00440-014-0574-8}

Bauer, M., Bruveris, M., \& Michor, P. W. (2016). Uniqueness of the
Fisher--Rao metric on the space of smooth densities. \emph{Bulletin of
the London Mathematical Society, 48}(3), 499--506.
\url{https://doi.org/10.1112/blms/bdw020}

Campbell, L. L. (1986). An extended Čencov characterization of the
information metric. \emph{Proceedings of the American Mathematical
Society, 98}(1), 135--141.

Čencov, N. N. (1982). \emph{Statistical Decision Rules and Optimal
Inference}. American Mathematical Society.

Petz, D., \& Sudár, C. (1996). Geometries of quantum states.
\emph{Journal of Mathematical Physics, 37}(6), 2662--2673.
\url{https://doi.org/10.1063/1.531535}

Hirsch, M. W., Smale, S., \& Devaney, R. L. (2013). \emph{Differential
Equations, Dynamical Systems, and an Introduction to Chaos} (3rd ed.).
Academic Press.

Jacobson, T. (1995). Thermodynamics of spacetime: The Einstein equation
of state. \emph{Physical Review Letters, 75}, 1260--1263.

Kraus, S., Lehmann, D., \& Magidor, M. (1990). Nonmonotonic reasoning,
preferential models and cumulative logics. \emph{Artificial
Intelligence, 44}(1--2), 167--207.

Lehmann, D., \& Magidor, M. (1992). What does a conditional knowledge
base entail? \emph{Artificial Intelligence, 55}(1), 1--60.

Shorrocks, A. F. (1980). The class of additively decomposable inequality
measures. \emph{Econometrica, 48}(3), 613--625.

Theil, H. (1967). \emph{Economics and Information Theory}.
North-Holland.

Rader, G. K. D. (2026). \emph{Coherence Universalism:
Physics---Coherence, Accessibility, and the Ordered Emergence of
Physical Description} (Version 0.10). Heaven≡Earth Foundation.

Rader, G. K. D. (2026). \emph{Coherence Universalism: Metaphysics and
Epistemology---Coherence Under Constraint, Explanatory Order, and
Warrant Across Registers} (Version 0.10). Heaven≡Earth Foundation.

Rader, G. K. D. (2026). \emph{Coherence Universalism: Biology---Life as
Active Self-Maintenance: Viability, Individuation, and Multiscale
Regulation} (Version 0.10). Heaven≡Earth Foundation.

Rader, G. K. D. (2026). \emph{Coherence Universalism:
Consciousness---Perspective, Temporal Depth, Valence, and the
Architecture of Experience} (Version 0.10). Heaven≡Earth Foundation.

Rader, G. K. D. (2026). \emph{Coherence Universalism:
Psychology---Meaning, Regulation, and Repair Across Time} (Version
0.10). Heaven≡Earth Foundation.

Rader, G. K. D. (2026). \emph{Coherence Universalism: Social
Dynamics---Coordination, Power, Meaning, and Repair in Plural Societies}
(Version 0.10). Heaven≡Earth Foundation.

Rader, G. K. D. (2026). \emph{Coherence Universalism: Ethics---Truth,
Self-Deception, Power, and Repair} (Version 0.10). Heaven≡Earth
Foundation.

Rader, G. K. D. (2026). \emph{Coherence Universalism: Artificial
Intelligence---Capability, Alignment, Consciousness, and the Governance
of Non-Egoic Systems} (Version 0.10). Heaven≡Earth Foundation.

Rader, G. K. D. (2026). \emph{Denying Knowledge to Make Room for
Commitment---Truth, Commitment, and Nonmonotonic Warrant} (Version
0.10). Heaven≡Earth Foundation.

Rader, G. K. D. (2026). \emph{The Architecture of Non-Egoic Artificial
Mind---Perspective, Self-Maintenance, and Alignment Without a Sovereign
Self} (Version 0.10). Heaven≡Earth Foundation.

Rader, G. K. D. (2026). \emph{Coherence Universalism v0.10 Master
Principle Consolidation and Disposition Ledger}. RC1.

Rader, G. K. D. (2026). \emph{You Are the Universe Becoming Conscious of
Itself: The Coherence Mandala}. Heaven≡Earth Press.

\end{document}
